%==============================================================================
%  Rapport de Projet de Fin d'Études — Data Quality Wevioo
%  Compilation : pdfLaTeX (Overleaf : Menu > Compiler = pdfLaTeX)
%  Compiler deux fois pour mettre à jour la table des matières et les listes.
%==============================================================================
\documentclass[12pt,a4paper,oneside]{report}

%------------------------------------------------------------------------------
%  Encodage, langue, police
%------------------------------------------------------------------------------
\usepackage[T1]{fontenc}
\usepackage[utf8]{inputenc}
\IfFileExists{french.ldf}
  {\usepackage[english,french]{babel}}
  {\usepackage[english]{babel}\babelprovide[import,main]{french}}
\usepackage{mathptmx}          % Times (comme le modèle)
\usepackage[scaled=0.92]{helvet}
\usepackage{courier}
\usepackage{csquotes}
\usepackage{underscore}   % autorise la coupure après « _ » dans \texttt
\setlength{\emergencystretch}{3em}

%------------------------------------------------------------------------------
%  Mise en page
%------------------------------------------------------------------------------
\usepackage[a4paper,top=2.5cm,bottom=2.5cm,left=2.5cm,right=2.5cm,headheight=15pt]{geometry}
\usepackage{setspace}
\onehalfspacing
\setlength{\parindent}{1cm}
\setlength{\parskip}{4pt}

%------------------------------------------------------------------------------
%  Paquets généraux
%------------------------------------------------------------------------------
\usepackage{graphicx}
\usepackage{float}
\usepackage{xcolor}
\usepackage{amsmath,amssymb}
\usepackage{array,tabularx,longtable,booktabs,multirow,makecell}
\usepackage[table]{xcolor}
\usepackage{enumitem}
\setlist{itemsep=2pt,topsep=4pt}
\setlist[itemize]{label=\textbullet}
\usepackage{caption}
\usepackage{listings}
\usepackage{fancyhdr}
\usepackage{titlesec}
\usepackage{titletoc}
\usepackage{etoolbox}
\usepackage[hidelinks]{hyperref}

%------------------------------------------------------------------------------
%  Couleurs
%------------------------------------------------------------------------------
\definecolor{espritred}{RGB}{200,16,46}
\definecolor{captionblue}{RGB}{46,84,150}
\definecolor{headgray}{RGB}{90,90,90}
\definecolor{tablehead}{RGB}{217,226,243}
\definecolor{codebg}{RGB}{247,247,250}
\definecolor{codekw}{RGB}{0,51,153}
\definecolor{codestr}{RGB}{163,21,21}
\definecolor{codecom}{RGB}{0,128,0}
\definecolor{wbox}{RGB}{235,241,250}
\definecolor{wboxborder}{RGB}{46,84,150}

%------------------------------------------------------------------------------
%  Légendes (style du modèle : bleu, italique, petit)
%------------------------------------------------------------------------------
\captionsetup{font={small,it,color=captionblue},labelfont={small,it,color=captionblue},labelsep=colon,justification=centering}
\captionsetup[table]{position=bottom}
\addto\captionsfrench{%
  \renewcommand{\figurename}{Figure}%
  \renewcommand{\tablename}{Tableau}%
  \renewcommand{\listfigurename}{Liste des Figures}%
  \renewcommand{\listtablename}{Liste des Tableaux}%
  \renewcommand{\contentsname}{Sommaire}%
  \renewcommand{\bibname}{Bibliographie}%
}
\AtBeginDocument{%
  \renewcommand{\figurename}{Figure}%
  \renewcommand{\tablename}{Tableau}%
  \renewcommand{\listfigurename}{Liste des Figures}%
  \renewcommand{\listtablename}{Liste des Tableaux}%
  \renewcommand{\contentsname}{Sommaire}%
  \renewcommand{\bibname}{Bibliographie}%
  \renewcommand{\chaptername}{Chapitre}%
}
% Numérotation continue des figures et tableaux (Figure 1, 2, 3 ...)
\counterwithout{figure}{chapter}
\counterwithout{table}{chapter}
\renewcommand{\lstlistingname}{Extrait de code}

%------------------------------------------------------------------------------
%  Titres de chapitres et sections
%------------------------------------------------------------------------------
\titleformat{\chapter}[display]
  {\normalfont\bfseries\centering}
  {\LARGE\chaptertitlename\ \thechapter}
  {8pt}
  {\titlerule[1.2pt]\vspace{6pt}\Huge}
  [{\vspace{4pt}\titlerule[1.2pt]}]
\titlespacing*{\chapter}{0pt}{20pt}{30pt}


\renewcommand{\thesection}{\Roman{section}}
\renewcommand{\thesubsection}{\Roman{section}.\arabic{subsection}}
\renewcommand{\thesubsubsection}{\Roman{section}.\arabic{subsection}.\arabic{subsubsection}}
\titleformat{\section}{\normalfont\Large\bfseries\color{black}}{\thesection.}{0.6em}{}
\titleformat{\subsection}{\normalfont\large\bfseries}{\thesubsection}{0.6em}{}
\titleformat{\subsubsection}{\normalfont\normalsize\bfseries\itshape}{\thesubsubsection}{0.6em}{}
\titlespacing*{\section}{0pt}{18pt}{8pt}
\titlespacing*{\subsection}{0pt}{14pt}{6pt}
\setcounter{secnumdepth}{3}
\setcounter{tocdepth}{2}

%------------------------------------------------------------------------------
%  En-têtes et pieds de page (nom du chapitre en haut, n° de page en bas à droite)
%------------------------------------------------------------------------------
\pagestyle{fancy}
\fancyhf{}
\renewcommand{\chaptermark}[1]{\markboth{\ifnum\value{chapter}>0 \chaptername\ \thechapter\ : \fi #1}{}}
\fancyhead[L]{\small\color{headgray}\leftmark}
\fancyfoot[R]{\small\thepage}
\renewcommand{\headrulewidth}{0.4pt}
\renewcommand{\footrulewidth}{0.4pt}
\fancypagestyle{plain}{%
  \fancyhf{}\fancyfoot[R]{\small\thepage}%
  \renewcommand{\headrulewidth}{0pt}\renewcommand{\footrulewidth}{0.4pt}}

%------------------------------------------------------------------------------
%  Listings (code)
%------------------------------------------------------------------------------
\lstdefinestyle{pfe}{
  backgroundcolor=\color{codebg},
  basicstyle=\ttfamily\footnotesize,
  keywordstyle=\color{codekw}\bfseries,
  stringstyle=\color{codestr},
  commentstyle=\color{codecom}\itshape,
  numbers=left, numberstyle=\tiny\color{gray}, numbersep=6pt,
  frame=single, rulecolor=\color{gray!40},
  breaklines=true, breakatwhitespace=false, showstringspaces=false,
  tabsize=2, columns=fullflexible, keepspaces=true,
  xleftmargin=14pt, framexleftmargin=14pt,
  captionpos=b, aboveskip=8pt, belowskip=8pt,
  literate={é}{{\'e}}1 {è}{{\`e}}1 {à}{{\`a}}1 {ç}{{\c c}}1 {ê}{{\^e}}1 {ô}{{\^o}}1 {É}{{\'E}}1 {→}{{$\rightarrow$}}1
}
\lstset{style=pfe}
\lstdefinelanguage{json}{
  morestring=[b]", stringstyle=\color{codestr},
  literate={:}{{\textcolor{codekw}{:}}}1 {,}{{\textcolor{codekw}{,}}}1
           {\{}{{\textcolor{codekw}{\{}}}1 {\}}{{\textcolor{codekw}{\}}}}1
           {[}{{\textcolor{codekw}{[}}}1 {]}{{\textcolor{codekw}{]}}}1
}
\lstdefinelanguage{yaml}{
  keywords={true,false,null}, morecomment=[l]{\#}, morestring=[b]", morestring=[b]'
}

%------------------------------------------------------------------------------
%  Emplacement de figure à remplir
%  \figplaceholder[largeur]{fichier.png}{Légende}{label}{Indication de contenu}
%  -> Déposez simplement l'image dans figures/ avec ce nom : elle s'affichera
%     automatiquement à la place du cadre.
%------------------------------------------------------------------------------
\NewDocumentCommand{\figplaceholder}{O{0.85} m m m m}{%
  \begin{figure}[H]
    \centering
    \IfFileExists{figures/#2}%
      {\includegraphics[width=#1\linewidth,height=0.7\textheight,keepaspectratio]{figures/#2}}%
      {\fcolorbox{espritred}{gray!6}{%
        \parbox[c]{#1\linewidth}{\centering\vspace*{1.1cm}
          {\large\bfseries\color{espritred} FIGURE À INSÉRER}\\[6pt]
          {\small\ttfamily figures/\detokenize{#2}}\\[8pt]
          \begin{minipage}{0.88\linewidth}\centering\small\itshape #5\end{minipage}\par\vspace*{1.1cm}}}}
    \caption{#3}\label{#4}
  \end{figure}}

% Logo d'outil : déposez figures/logos/<nom>.png et il remplace le cadre
\newcommand{\toolicon}[1]{\IfFileExists{figures/logos/#1}%
  {\raisebox{-0.35\height}{\includegraphics[height=0.9cm,width=1.6cm,keepaspectratio]{figures/logos/#1}}}%
  {\raisebox{-0.35\height}{\fcolorbox{espritred}{gray!6}{\parbox[c][0.75cm][c]{1.5cm}{\centering\tiny\color{espritred}logo\\ \detokenize{#1}}}}}}

% Espace réservé à une section dont le contenu sera fourni plus tard
\newcommand{\espacereserve}[1]{%
  \begin{center}
  \fcolorbox{espritred}{gray!4}{\parbox[c][#1][c]{0.92\linewidth}{\centering
    {\color{espritred}\bfseries Section réservée}\\[4pt]
    {\small\itshape Le contenu de cette section sera communiqué ultérieurement.}}}
  \end{center}}

% Encadré « à compléter » pour les informations à vérifier
\newcommand{\acompleter}[1]{\textcolor{espritred}{\textbf{[#1]}}}

% Encadré de synthèse
\newcommand{\encadre}[2]{%
  \begin{center}
  \fcolorbox{wboxborder}{wbox}{\parbox{0.92\linewidth}{\textbf{#1}\par\vspace{3pt}#2}}
  \end{center}}

% En-tête de tableau coloré
\newcommand{\thd}[1]{\cellcolor{tablehead}\textbf{#1}}
\newcolumntype{L}[1]{>{\raggedright\arraybackslash}p{#1}}
\newcolumntype{C}[1]{>{\centering\arraybackslash}p{#1}}
\renewcommand{\arraystretch}{1.25}

% Nom de la solution (modifiable ici : il sera mis à jour dans tout le rapport)
\newcommand{\projet}{DQ-Insight}
% Informations générales (à compléter)
\newcommand{\etudiant}{Ahmed FATHALLAH}
\newcommand{\encadrantacad}{Rim ALOUANE}
\newcommand{\encadrantpro}{Ines BEN AMOR}
\newcommand{\ecole}{ESPRIT}
\newcommand{\anneeuniv}{2025 -- 2026}

%==============================================================================
\begin{document}
%==============================================================================

%------------------------------ Page de garde ---------------------------------
%------------------------------------------------------------------------------
%  PAGE DE GARDE — provisoire
%  À remplacer plus tard par la page de garde officielle de l'école
%  (ex. : \includepdf{figures/page_de_garde.pdf} avec le paquet pdfpages).
%------------------------------------------------------------------------------
\begin{titlepage}
  \thispagestyle{empty}
  \centering
  {\large \ecole\par}
  \vspace{0.5cm}
  {\large Année universitaire \anneeuniv\par}
  \vspace{1.5cm}
  {\Huge\bfseries\color{espritred} PROJET DE FIN D'ÉTUDES\par}
  \vspace{0.6cm}
  {\Large Diplôme National d'Ingénieur\par}
  {\large Spécialité : Génie logiciel\par}
  \vspace{1.2cm}
  \rule{\linewidth}{1pt}\par\vspace{0.4cm}
  {\LARGE\bfseries Conception et mise en place d'une plateforme intelligente d'évaluation, de correction et d'interrogation de la qualité des données basée sur le profilage automatique, les grands modèles de langage locaux et un graphe de connaissances\par}
  \vspace{0.4cm}\rule{\linewidth}{1pt}\par
  \vspace{1.5cm}
  {\large Réalisé par : \textbf{\etudiant}\par}
  \vspace{1cm}
  {\large Encadré par :\par}
  {\large Encadrant académique : \encadrantacad\par}
  {\large Encadrant entreprise : \encadrantpro\par}
  \vfill
  {\large Organisme d'accueil : \textbf{Wevioo}\par}
\end{titlepage}


%------------------------------ Pages liminaires ------------------------------
\pagenumbering{roman}
\chapter*{Dédicaces}
\markboth{Dédicaces}{}
\vspace*{1cm}
\begin{center}
\begin{minipage}{0.82\linewidth}
\centering\itshape\large
\textbf{À mes chers parents},\\
pour leur amour inconditionnel, leurs sacrifices et leur soutien constant tout au long de mon parcours. Que ce travail soit le témoignage de ma profonde reconnaissance.

\vspace{1cm}
\textbf{À mes frères et sœurs},\\
pour leurs encouragements et leur présence à chaque étape.

\vspace{1cm}
\textbf{À mes amis et camarades de promotion},\\
pour les moments partagés, l'entraide et la motivation mutuelle.

\vspace{1cm}
\textbf{À tous ceux qui m'ont aidé},\\
de près ou de loin, à mener à bien ce projet.

\vspace{1.5cm}
\upshape\hfill \etudiant
\end{minipage}
\end{center}

\chapter*{Remerciements}
\markboth{Remerciements}{}

Au terme de ce projet de fin d'études, je tiens à exprimer ma sincère gratitude à toutes les personnes qui ont contribué, de près ou de loin, à sa réalisation.

J'adresse mes vifs remerciements à l'organisme d'accueil, \textbf{Wevioo}, pour m'avoir offert l'opportunité de réaliser ce stage dans un environnement professionnel stimulant, et pour la confiance qui m'a été accordée sur un sujet à la croisée de l'ingénierie des données et de l'intelligence artificielle.

Je remercie tout particulièrement mon encadrant professionnel, \encadrantpro, pour son suivi attentif, ses orientations techniques pertinentes et sa disponibilité tout au long du stage.

Je tiens également à exprimer ma profonde reconnaissance à mon encadrant(e) académique, \encadrantacad, pour ses conseils éclairés, sa rigueur et son accompagnement dans la rédaction de ce rapport.

Mes remerciements s'adressent aussi à l'ensemble de l'équipe de Wevioo pour son accueil chaleureux et son aide précieuse au quotidien.

J'exprime enfin ma gratitude aux membres du jury, qui ont accepté d'évaluer ce travail, ainsi qu'à l'ensemble du corps enseignant de \ecole{} pour la qualité de la formation dispensée.

\vspace{1cm}
\hfill \etudiant

\chapter*{Résumé}
\markboth{Résumé}{}
\addcontentsline{toc}{chapter}{Résumé}

Dans un contexte où les données constituent un actif stratégique, leur qualité conditionne directement la fiabilité des analyses, des tableaux de bord et des modèles décisionnels. Pourtant, l'évaluation de la qualité des données reste, dans de nombreuses organisations, une activité manuelle, longue et dépendante de l'expertise individuelle : les outils de profilage mesurent sans interpréter, et les outils de validation exigent des règles écrites à l'avance.

Ce projet de fin d'études, réalisé au sein de \textbf{Wevioo}, répond à cette problématique par la conception et la mise en place de \textbf{\projet}, une plateforme web d'évaluation intelligente de la qualité des données, fondée sur un principe : \textbf{le code mesure et contrôle, le modèle interprète et rédige}. Un traitement orchestré par \textbf{Apache Airflow} lit des fichiers hétérogènes, calcule les métriques de chaque table, détermine le rôle de chaque colonne à partir d'indices explicites et mesure les formats, l'ordre des dates et les liens entre tables. Un \textbf{grand modèle de langage exécuté localement} (Ollama) interprète ces mesures ; ses affirmations sont confrontées aux mesures, avec une analyse déterministe de repli.

Un serveur \textbf{Spring Boot} versionne les dépôts, sécurise l'accès par jetons, évalue des \textbf{règles de qualité} à chaque exécution et alerte en cas d'échec, applique des \textbf{corrections} sur une copie sans jamais modifier l'original, et construit un \textbf{graphe de connaissances} dans \textbf{Neo4j}. Une interface \textbf{Angular} restitue les résultats en temps réel. Un service de questions en langage naturel (\textbf{GraphRAG}) répond sur les données, leur qualité et l'usage de l'outil ; les requêtes Cypher écrites par le modèle sont contrôlées (lecture seule, périmètre d'une version, plan d'exécution, durée) avant d'être exécutées.

La plateforme est déployable par une commande Docker Compose et couverte par 292 tests automatisés. Une comparaison de trois modèles de langage sur un jeu de questions rejouable éclaire le choix du modèle.

\vspace{0.4cm}
\noindent\textbf{Mots-clés :} qualité des données, profilage de données, grands modèles de langage, LLM local, Ollama, GraphRAG, Neo4j, Spring Boot, Angular, Apache Airflow, Docker, gouvernance des données.


\cleardoublepage
\markboth{Sommaire}{}
\tableofcontents

\cleardoublepage
\phantomsection
\addcontentsline{toc}{chapter}{Liste des Figures}
\markboth{Table des figures}{}
\listoffigures

\cleardoublepage
\phantomsection
\addcontentsline{toc}{chapter}{Liste des Tableaux}
\markboth{Liste des tableaux}{}
\listoftables

\cleardoublepage
\chapter*{Liste des Abréviations}
\markboth{Liste des Abréviations}{}
\addcontentsline{toc}{chapter}{Liste des Abréviations}

{\renewcommand{\arraystretch}{1.15}
\begin{longtable}{>{\leavevmode\bfseries\color{espritred}}p{2.6cm} c p{10cm}}
API     & = & Application Programming Interface \\
CI      & = & Continuous Integration (intégration continue) \\
CSV     & = & Comma-Separated Values \\
CU      & = & Cas d'Utilisation \\
DAG     & = & Directed Acyclic Graph (graphe orienté acyclique) \\
DAMA    & = & Data Management Association \\
DMBOK   & = & Data Management Body of Knowledge \\
EDA     & = & Exploratory Data Analysis (analyse exploratoire des données) \\
ETL     & = & Extract, Transform, Load \\
FCM     & = & Factors -- Criteria -- Metrics \\
GenAI   & = & Generative Artificial Intelligence (IA générative) \\
GraphRAG & = & RAG appliqué à un graphe de connaissances \\
GPU     & = & Graphics Processing Unit \\
HTML    & = & HyperText Markup Language \\
HTTP    & = & HyperText Transfer Protocol \\
IA      & = & Intelligence Artificielle \\
ID      & = & Identifiant \\
JPA     & = & Java Persistence API \\
ISO     & = & International Organization for Standardization \\
JSON    & = & JavaScript Object Notation \\
JWT     & = & JSON Web Token \\
LLM     & = & Large Language Model (grand modèle de langage) \\
LLMOps  & = & Large Language Model Operations \\
MoSCoW  & = & Must, Should, Could, Won't \\
MLOps   & = & Machine Learning Operations \\
RAG     & = & Retrieval-Augmented Generation (génération augmentée par la recherche) \\
RAM     & = & Random Access Memory \\
REST    & = & Representational State Transfer \\
RGPD    & = & Règlement Général sur la Protection des Données \\
SSE     & = & Server-Sent Events \\
SQL     & = & Structured Query Language \\
UML     & = & Unified Modeling Language \\
UI      & = & User Interface \\
VRAM    & = & Video Random Access Memory \\
YAML    & = & YAML Ain't Markup Language \\
\end{longtable}
\addtocounter{table}{-1}}


%------------------------------ Corps du rapport ------------------------------
\cleardoublepage
\pagenumbering{arabic}
\chapter*{Introduction Générale}
\markboth{Introduction Générale}{}
\addcontentsline{toc}{chapter}{Introduction Générale}

Une entreprise vit de ses données : tableaux de bord, rapports réglementaires, modèles d'apprentissage. Mais ces données ne valent que ce que vaut leur qualité, et on la regarde rarement tant que tout semble fonctionner. Il suffit pourtant d'une colonne décalée à l'export, d'un identifiant en double ou d'une date mal formée pour fausser un indicateur et retarder une décision.

Dans la pratique, contrôler la qualité reste un travail manuel. On ouvre chaque fichier, on calcule quelques statistiques, on repère ce qui cloche, puis on traduit le tout en actions pour les équipes métier. Les outils de profilage font bien la première étape, mais ils rendent de longs rapports techniques que peu de gens lisent jusqu'au bout, sans dire quoi faire. Les outils de validation par règles demandent de connaître les règles à l'avance, ce qui est rarement le cas devant un jeu de données qu'on découvre. Et aucun des deux ne corrige les données ni ne permet de les interroger sans écrire de requête.

Les grands modèles de langage changent la donne, à condition de les utiliser avec précaution. Ils savent lire des informations structurées et les expliquer simplement, et on peut désormais les faire tourner en local avec Ollama, sans envoyer de données sensibles à l'extérieur : c'est essentiel pour l'assurance ou la banque. Mais ils se trompent avec aplomb. Ils citent une colonne qui n'existe pas, attribuent un pourcentage à la mauvaise colonne, écrivent une requête qui modifie la base.

Ce projet de fin d'études, réalisé chez \textbf{Wevioo}, part de ce constat. J'ai construit \textbf{\projet}, une plateforme web qui suit une règle simple : \textbf{le code mesure et contrôle, le modèle explique}. Pour chaque jeu de données, elle donne un score de qualité expliqué, le rôle de chaque colonne, et une analyse en langage naturel dont chaque affirmation est vérifiée contre les mesures. On peut y définir des règles surveillées à chaque exécution, corriger les données sur une copie, construire un graphe dans Neo4j et poser des questions en français ou en anglais, dont les requêtes sont contrôlées avant d'être exécutées. L'ensemble est orchestré par Apache Airflow, repose sur Spring Boot et Angular, et se lance avec une seule commande Docker Compose.

\section*{Problématique}
\addcontentsline{toc}{section}{Problématique}

La question à laquelle ce projet répond est la suivante :

\encadre{Question de recherche}{\itshape Comment automatiser l'évaluation de la qualité de données tabulaires pour obtenir non seulement des indicateurs fiables, mais aussi une interprétation métier, une surveillance par des règles, des corrections sûres et une interrogation en langage naturel, tout en gardant les données en local, les affirmations du modèle vérifiables et les résultats traçables ?}

Elle se décompose en problèmes plus précis. Comment déterminer le rôle d'une colonne sans s'en remettre au seul modèle ? Comment confronter ce que dit le modèle aux mesures, et garantir qu'un résultat exploitable sort toujours ? Comment corriger des données sans jamais toucher à l'original ni laisser le modèle inventer des valeurs ? Comment représenter plusieurs tables sous forme de graphe quand elles n'ont pas d'identifiant déclaré ? Et comment laisser un modèle écrire des requêtes sans lui faire confiance ?

\section*{Organisation du rapport}
\addcontentsline{toc}{section}{Organisation du rapport}

Le rapport compte huit chapitres.

\begin{itemize}
  \item Le \textbf{chapitre 1} présente Wevioo, le contexte du projet, l'étude de l'existant et la méthode de travail.
  \item Le \textbf{chapitre 2} fait l'état de l'art (qualité et profilage des données, modèles de langage, orchestration, applications web, graphes, GraphRAG) et justifie les choix technologiques.
  \item Le \textbf{chapitre 3} spécifie les besoins : acteurs, cas d'utilisation, backlog et exigences non fonctionnelles.
  \item Le \textbf{chapitre 4} décrit la conception : modèle UML, architectures, sécurité et composants d'intelligence artificielle.
  \item Les \textbf{chapitres 5, 6 et 7} racontent la réalisation en sept sprints, du profilage jusqu'aux questions en langage naturel.
  \item Le \textbf{chapitre 8} présente l'évaluation : tests automatisés, tests fonctionnels et comparaison de trois modèles de langage sur un jeu de questions.
\end{itemize}

Le rapport se termine par un bilan du travail et des pistes d'évolution.

\chapter{Contexte Général du Projet}

\section*{Introduction}
\addcontentsline{toc}{section}{Introduction}
Ce premier chapitre pose le cadre général du projet. Nous présentons d'abord l'organisme d'accueil, Wevioo, puis le contexte dans lequel s'inscrit le sujet et la problématique qu'il adresse. Nous étudions ensuite les solutions existantes selon des critères définis, afin d'en dégager les limites, avant de présenter la solution proposée et ses objectifs. Nous terminons par la méthodologie de travail adoptée.

%==============================================================================
\section{Organisme d'accueil}
%==============================================================================

Ce projet de fin d'études a été réalisé au sein de Wevioo. Avant d'aborder le sujet lui-même, il est utile de présenter l'entreprise, son organisation et ses domaines d'activité, qui éclairent les enjeux du projet.

\subsection{Présentation de Wevioo}

\textbf{Wevioo} est un groupe international de conseil et de services numériques, fondé en 1998, qui accompagne ses clients dans leurs enjeux d'innovation, de transformation et de performance opérationnelle~\cite{wevioo}. Il dispose de bureaux à Paris, Tunis, Riyadh, Dubaï et Abidjan, et mobilise plus de 460 collaborateurs.

Le groupe affirme sa vision en ces termes : dans un monde marqué par le numérique, les données et l'intelligence artificielle, la technologie doit avoir du sens et servir de levier à une performance responsable et à la création de valeur durable. Sa mission est de transformer la complexité en décisions claires et en exécution fiable, en combinant compétences métier et excellence technologique (intelligence artificielle, \textit{low-code}, \textit{cloud}, cybersécurité, données). Quatre valeurs structurent son approche : innover, partager, prendre soin (\textit{care}) et livrer (\textit{deliver}).

Wevioo revendique des projets menés dans plus de trente pays de la zone EMEA, avec de grands acteurs de l'assurance, de la banque, de l'industrie et du secteur public.

\figplaceholder[0.45]{logo_wevioo.png}{Logo de Wevioo}{fig:logo}{Logo de Wevioo.}

Après cette présentation générale, nous décrivons l'organisation de l'offre du groupe et situons l'équipe dans laquelle le stage s'est déroulé.

\subsection{Organisation et structure}

Wevioo ne publie pas d'organigramme hiérarchique. Son organisation se lit à travers ses implantations, ses domaines d'expertise et les secteurs qu'il sert (figure~\ref{fig:organigramme}). Notre stage s'est déroulé au sein de l'équipe \textbf{Data}, qui relève du domaine d'expertise \textit{Data \& Analytics}, voisin du domaine \textit{Artificial Intelligence}.

\figplaceholder[1]{organigramme_wevioo.png}{Organisation du groupe Wevioo : implantations, domaines d'expertise et secteurs}{fig:organigramme}{Schéma de l'organisation de l'offre de Wevioo, d'après son site institutionnel, avec l'équipe d'accueil mise en évidence.}

Cette organisation reflète la diversité des services proposés par Wevioo, que nous détaillons à présent.

\subsection{Services}

Wevioo présente son offre à travers douze domaines d'expertise : \textit{Application Services}, \textit{Artificial Intelligence}, \textit{Business Strategy \& Transformation}, \textit{Cloud \& IT Infrastructure}, \textit{Cybersecurity}, \textit{Data \& Analytics}, \textit{Enterprise Architecture}, \textit{Managed Services}, \textit{Nearshore Development}, \textit{Operational Excellence}, \textit{Process Automation} et \textit{Software Development}. Elle s'appuie sur un écosystème de partenaires stratégiques dans le \textit{cloud}, les données, l'intelligence artificielle et le \textit{low-code}.

C'est précisément à l'intersection des domaines \textit{Data \& Analytics} et \textit{Artificial Intelligence} et des besoins des clients du secteur de l'assurance que se situe notre projet, dont nous présentons maintenant le contexte.

%==============================================================================
\section{Contexte général et problématique}
%==============================================================================

Cette section décrit le contexte métier dans lequel le projet a vu le jour, puis formalise la problématique à laquelle il répond.

\subsection{Contexte général du projet}

Les clients de Wevioo, en particulier dans les secteurs de l'\textbf{assurance} et de la finance, manipulent des données issues de multiples systèmes : référentiels clients, contrats et polices, avenants, sinistres, agences et intermédiaires, encaissements. Ces données sont régulièrement extraites sous forme de fichiers tabulaires pour alimenter des projets de migration, de reporting ou d'analyse. Avant toute exploitation, il est indispensable d'en évaluer la qualité : un champ de date mal formé, un identifiant client dupliqué ou une colonne décalée lors de l'export peuvent compromettre l'ensemble d'un projet.

Dans les missions de données, cette phase d'audit de la qualité est traditionnellement conduite par des consultants qui explorent manuellement chaque table. Elle mobilise un temps important, dépend fortement de l'expertise individuelle et produit des livrables hétérogènes d'une mission à l'autre. L'émergence des grands modèles de langage ouvre une perspective nouvelle : automatiser non seulement le calcul des indicateurs, mais aussi leur \textbf{interprétation} et la formulation de \textbf{recommandations}, à condition de maîtriser la fiabilité des résultats et la confidentialité des données.

Ce contexte fait apparaître un ensemble de difficultés concrètes, que nous formalisons sous forme de problématique.

\subsection{Problématique}

L'analyse du processus actuel d'évaluation de la qualité des données fait ressortir plusieurs difficultés :
\begin{itemize}
  \item \textbf{Un effort manuel important} : chaque jeu de données doit être chargé, exploré et documenté individuellement, ce qui devient rapidement ingérable lorsque le portefeuille compte des dizaines de tables.
  \item \textbf{Des rapports techniques difficiles à exploiter} : les outils de profilage produisent des rapports exhaustifs mais sans interprétation, et leur lecture requiert une expertise statistique.
  \item \textbf{L'absence de vision transverse} : l'analyse est menée table par table, sans identification des entités partagées (client, contrat, agence) ni des causes communes aux anomalies observées dans plusieurs sources.
  \item \textbf{Des contraintes fortes de confidentialité} : les données d'assurance contiennent des informations personnelles et commerciales sensibles, ce qui exclut l'envoi de ces données à des services d'IA hébergés par des tiers.
  \item \textbf{Un manque de reproductibilité et de traçabilité} : les analyses manuelles sont difficilement rejouables, et leurs conclusions ne sont ni historisées ni auditables.
\end{itemize}

Pour évaluer dans quelle mesure ces difficultés sont déjà prises en charge par des solutions existantes, nous avons mené une étude de l'existant, présentée dans la section suivante.

%==============================================================================
\section{Étude de l'existant}
%==============================================================================

L'étude de l'existant suit une démarche en quatre temps : nous définissons d'abord les critères d'évaluation, nous étudions ensuite chaque solution individuellement, puis nous les comparons avant d'en faire la critique.

\subsection{Critères d'évaluation}

Les critères suivants découlent directement de la problématique exposée ci-dessus :

\begin{table}[H]
\centering
\small
\begin{tabular}{|C{1.1cm}|L{4.2cm}|L{8.7cm}|}
\hline
\thd{Code} & \thd{Critère} & \thd{Question posée} \\ \hline
C1 & Automatisation du profilage & La solution calcule-t-elle automatiquement les indicateurs de qualité de chaque colonne ? \\ \hline
C2 & Absence de règles préalables & Peut-on l'utiliser sur un jeu de données inconnu, sans écrire de règles au préalable ? \\ \hline
C3 & Interprétation en langage naturel & Produit-elle une lecture compréhensible par un profil non technique ? \\ \hline
C4 & Recommandations priorisées & Propose-t-elle des actions correctives classées par impact et par effort ? \\ \hline
C5 & Vision transverse & Rapproche-t-elle les entités et les problèmes communs à plusieurs tables ? \\ \hline
C6 & Confidentialité & Les données restent-elles sur l'infrastructure locale ? \\ \hline
C7 & Industrialisation & Peut-elle être orchestrée, rejouée et tracée ? \\ \hline
C8 & Coût & Quel est le coût de licence et de mise en œuvre ? \\ \hline
\end{tabular}
\caption{Critères d'évaluation des solutions existantes}
\label{tab:criteres}
\end{table}

Ces critères posés, nous pouvons examiner une à une les principales solutions disponibles.

\subsection{Étude des solutions existantes}

Nous avons retenu six familles de solutions, représentatives des pratiques actuelles, de la plus manuelle à la plus intégrée.

\subsubsection{L'exploration manuelle}

\textbf{Principe.} L'analyste charge chaque fichier dans un tableur ou un environnement d'analyse et calcule lui-même les indicateurs dont il a besoin : valeurs manquantes, valeurs distinctes, distributions, doublons.

\textbf{Fonctionnalités.} Totale liberté d'analyse ; possibilité d'interpréter les résultats au regard du métier et de rédiger des recommandations.

\textbf{Limites.} Approche lente, non reproductible, dont la qualité dépend entièrement de l'expertise de la personne ; elle ne passe pas à l'échelle d'un portefeuille de plusieurs dizaines de tables.

\subsubsection{ydata-profiling}

\textbf{Principe.} Bibliothèque open source~\cite{ydata} qui génère automatiquement un rapport d'analyse exploratoire d'un jeu de données.

\textbf{Fonctionnalités.} Statistiques par colonne, détection de types, alertes automatiques (forte cardinalité, valeurs constantes, fortes corrélations), matrices de corrélation, export structuré.

\textbf{Limites.} Rapport descriptif sans interprétation métier ni recommandation ; temps de calcul élevé sur les tables volumineuses ; analyse limitée à une table à la fois.

\subsubsection{SweetViz}

\textbf{Principe.} Bibliothèque open source~\cite{sweetviz} d'analyse exploratoire, qui produit un rapport interactif et très visuel. C'est la bibliothèque que nous avons \textbf{retenue et exploitée} dans ce projet comme socle du profilage.

\textbf{Fonctionnalités.} Résumé du jeu de données, statistiques détaillées par colonne, valeurs fréquentes, histogrammes et, surtout, une analyse des associations adaptée au type des variables, ainsi que la comparaison de deux jeux de données.

\textbf{Limites.} Les statistiques ne sont pas exportées de manière structurée ; aucune interprétation ni recommandation ; certaines étapes internes sont coûteuses sur les gros fichiers ; analyse limitée à une table.

\subsubsection{Great Expectations}

\textbf{Principe.} Cadre open source~\cite{gx} de validation des données fondé sur des « attentes » déclaratives (par exemple : « cette colonne ne doit pas contenir de valeur nulle »).

\textbf{Fonctionnalités.} Large bibliothèque d'attentes prédéfinies, exécution automatique des contrôles, documentation générée, intégration aux orchestrateurs.

\textbf{Limites.} Les règles doivent être connues et écrites à l'avance par un expert ; l'outil vérifie mais n'explique pas, et il est peu adapté à la découverte d'un nouveau jeu de données.

\subsubsection{Deequ et Soda Core}

\textbf{Principe.} Deequ~\cite{deequ}, développé par Amazon, et Soda Core~\cite{soda} appliquent la même logique de \textbf{contrôles déclaratifs}, à grande échelle pour le premier et à l'aide d'un langage de contrôle dédié pour le second.

\textbf{Fonctionnalités.} Calcul de métriques et vérification de contraintes sur de gros volumes, suggestion automatique de contraintes, suivi des métriques dans le temps.

\textbf{Limites.} Mêmes limites que Great Expectations : écriture préalable des règles, absence d'interprétation et de recommandations ; infrastructure distribuée nécessaire pour Deequ.

\subsubsection{Les suites commerciales}

\textbf{Principe.} Plateformes intégrées (Informatica Data Quality, Talend Data Quality) couvrant l'ensemble du cycle de la qualité des données.

\textbf{Fonctionnalités.} Profilage, règles métier, standardisation, déduplication, tableaux de bord, gestion des anomalies et des processus de correction.

\textbf{Limites.} Coûts de licence élevés, mise en œuvre lourde, inadaptées à une mission d'audit ponctuelle ; intégration encore limitée de l'IA générative en environnement local.

Après cette étude individuelle, il convient de confronter ces solutions entre elles au regard des critères définis.

\subsection{Étude comparative}

Le tableau~\ref{tab:comparatif-existant} confronte les solutions étudiées aux critères définis.

\begin{table}[H]
\centering
\scriptsize
\begin{tabular}{|L{2.1cm}|L{2.2cm}|L{2.3cm}|L{2.4cm}|L{2.4cm}|L{2.4cm}|}
\hline
\thd{Critère} & \thd{Exploration manuelle} & \thd{ydata-profiling / SweetViz} & \thd{Great Exp. / Deequ / Soda} & \thd{Suites commerciales} & \thd{\projet} \\ \hline
C1 Profilage & Calcul manuel, colonne par colonne & Automatique et complet & Métriques ciblées par les règles & Automatique et complet & Automatique, avec rôle de chaque colonne \\ \hline
C2 Sans règles & Oui & Oui & Non : règles à écrire & Non : règles à configurer & Oui, et règles proposées à partir de l'objectif \\ \hline
C3 Interprétation & Rédigée par l'analyste & Aucune & Aucune & Limitée & Générée par un LLM local, vérifiée par les mesures \\ \hline
C4 Recommandations & Selon l'expertise & Aucune & Aucune (échecs de contrôle) & Processus de correction & Règles, alertes et corrections appliquées sur une copie \\ \hline
C5 Vision transverse & Rarement & Non (une table) & Non & Partielle & Modèle de données, graphe Neo4j et questions en langage naturel \\ \hline
C6 Confidentialité & Locale & Locale & Locale & Selon déploiement & Locale (LLM local) \\ \hline
C7 Industrialisation & Non & Non & Oui & Oui & Oui \\ \hline
C8 Coût & Temps humain élevé & Gratuit & Gratuit & Licences élevées & Gratuit (open source) \\ \hline
\end{tabular}
\caption{Comparaison des solutions existantes selon les critères retenus}
\label{tab:comparatif-existant}
\end{table}

Cette comparaison fait ressortir, pour chaque solution, des manques précis que nous analysons dans la critique qui suit.

\subsection{Critique de l'existant}

Le tableau~\ref{tab:critique} synthétise, pour chaque solution, les critères qu'elle ne satisfait pas et l'impact de ces limites dans notre contexte.

\begin{table}[H]
\centering
\small
\begin{tabular}{|L{3.4cm}|L{2.6cm}|L{8cm}|}
\hline
\thd{Solution} & \thd{Critères non satisfaits} & \thd{Impact dans notre contexte} \\ \hline
Exploration manuelle & C1, C5, C7, C8 & Ne passe pas à l'échelle d'un portefeuille ; résultats hétérogènes et non reproductibles d'une mission à l'autre. \\ \hline
ydata-profiling, SweetViz & C3, C4, C5, C7 & Fournissent les bons indicateurs, mais laissent à l'expert tout le travail d'interprétation et de recommandation. \\ \hline
Great Expectations, Deequ, Soda & C2, C3, C4, C5 & Inutilisables lors de la découverte d'un nouveau jeu de données, puisque les règles ne sont pas encore connues. \\ \hline
Suites commerciales & C3, C6, C8 & Coût et lourdeur disproportionnés pour une mission d'audit ; IA générative peu disponible en local. \\ \hline
\end{tabular}
\caption{Critique des solutions existantes}
\label{tab:critique}
\end{table}

Il en ressort qu'aucune solution ne couvre simultanément l'automatisation du profilage, l'interprétation, les recommandations, la vision transverse et la confidentialité. Ce constat justifie la conception d'une nouvelle solution, que nous présentons à présent.

\subsection{Solution proposée}

Pour répondre à ces limites, nous proposons \textbf{\projet}, une plateforme web d'évaluation automatisée et intelligente de la qualité des données tabulaires. Son principe est de \textbf{séparer ce qui se mesure de ce qui s'explique} : le code mesure de façon déterministe l'état de chaque table (valeurs manquantes, doublons, formats, rôle de chaque colonne, liens entre tables), un \textbf{grand modèle de langage exécuté localement} interprète ces mesures et propose des actions, et chacune de ses affirmations est confrontée aux mesures avant d'être affichée. À partir de ces résultats, la plateforme permet de définir des \textbf{règles de qualité} surveillées à chaque exécution, d'appliquer des \textbf{corrections} sur une copie des données, de construire un \textbf{graphe de connaissances} dans Neo4j et de \textbf{poser des questions en langage naturel} sur les données, leur qualité et l'usage de l'outil. Elle fonctionne sans intervention manuelle, sur l'infrastructure de l'entreprise, et conserve la trace de chacun de ses traitements.

\figplaceholder[0.9]{vue_generale_solution.png}{Vue générale de la solution proposée}{fig:vue-generale}{Schéma synthétique et simple : données sources $\rightarrow$ profilage statistique $\rightarrow$ analyse intelligente (LLM local) $\rightarrow$ résultats et recommandations, avec l'orchestration et la traçabilité en transverse.}

La solution s'articule autour de cinq blocs : une interface web, un serveur applicatif, un traitement de profilage orchestré par Airflow, un service de questions et une base de graphe. Leur conception est détaillée au chapitre~4 et leur réalisation, sprint par sprint, aux chapitres~5 à~7. Il convient auparavant de préciser les objectifs que la solution doit atteindre.

\subsection{Objectifs du projet}

L'\textbf{objectif général} du projet est de réduire significativement le temps et l'expertise nécessaires à l'audit de la qualité d'un portefeuille de données, tout en améliorant la lisibilité et la traçabilité des conclusions. Il se décline en objectifs spécifiques :

\begin{itemize}
  \item \textbf{Automatiser l'évaluation de bout en bout} : tout jeu de données déposé doit être lu, profilé, analysé et restitué sans configuration préalable ni écriture de règles.
  \item \textbf{Fiabiliser les mesures} : toutes les statistiques doivent être calculées de façon déterministe et vérifiable, l'intelligence artificielle intervenant pour les interpréter et ses affirmations étant contrôlées.
  \item \textbf{Rendre les résultats intelligibles} pour des profils non techniques : score et note, dimensions de qualité expliquées, recommandations en langage naturel.
  \item \textbf{Intégrer le contexte métier} : permettre à l'utilisateur de décrire le domaine, l'objectif et ce que représente une ligne, pour orienter l'analyse et proposer des règles adaptées.
  \item \textbf{Surveiller et corriger} : évaluer des règles de qualité à chaque exécution, alerter en cas d'échec et corriger les données sur une copie, sans jamais modifier l'original.
  \item \textbf{Offrir une vision transverse} : relier les tables entre elles, construire un graphe de leurs entités et permettre de l'interroger en langage naturel.
  \item \textbf{Garantir la confidentialité} en exécutant les modèles de langage exclusivement sur l'infrastructure locale.
  \item \textbf{Industrialiser et sécuriser} : exécution orchestrée avec reprise sur erreur, authentification, déploiement reproductible et tests automatisés.
  \item \textbf{Assurer la traçabilité} : versionner les jeux de données et conserver la trace de chaque exécution, de chaque correction et de chaque réponse du modèle.
\end{itemize}

Atteindre ces objectifs dans un contexte aussi exploratoire supposait une méthode de travail adaptée, que nous décrivons dans la dernière section de ce chapitre.

%==============================================================================
\section{Méthodologie de travail}
%==============================================================================

Cette section justifie le choix d'une démarche agile, présente le framework Scrum retenu et décrit la manière dont il a été appliqué au projet.

\subsection{La méthodologie Agile}

L'\textbf{agilité} est une approche de gestion de projet fondée sur le Manifeste Agile (2001), qui privilégie les individus et leurs interactions, un produit opérationnel, la collaboration avec le client et l'adaptation au changement. Le produit est construit de manière \textbf{itérative et incrémentale} : il est livré par petites parties fonctionnelles, présentées régulièrement aux parties prenantes, dont les retours orientent la suite du travail.

Cette approche est particulièrement adaptée à notre projet, dont le caractère est fortement exploratoire : la qualité des réponses d'un modèle de langage, la manière de structurer les instructions qui lui sont données ou le choix entre plusieurs modèles ne pouvaient être déterminés qu'en expérimentant. Les besoins se sont par ailleurs précisés au fil des démonstrations à l'encadrant, au point de faire évoluer le périmètre du produit au fil des sprints, des règles de qualité jusqu'aux questions en langage naturel. Une démarche agile nous a permis d'intégrer ces évolutions sans remettre en cause le travail déjà réalisé.

Parmi les cadres de mise en œuvre de l'agilité, nous avons retenu Scrum, présenté ci-après.

\subsection{Le framework Scrum}

\textbf{Scrum}~\cite{scrum} n'est pas une méthodologie en soi, mais un \textbf{framework} (cadre de travail) qui met en œuvre la méthodologie Agile. Il définit un ensemble minimal de rôles, d'artefacts et d'événements, et organise le travail en itérations de durée fixe, au plus d'un mois, appelées \textbf{sprints}, au terme desquelles un incrément utilisable du produit est livré. Nous l'avons retenu pour sa simplicité, son découpage en itérations courtes et la place qu'il accorde aux retours réguliers.

\figplaceholder[0.8]{cycle_scrum.png}{Cycle de vie Scrum}{fig:scrum}{Schéma classique du cycle Scrum : Product Backlog $\rightarrow$ Sprint Planning $\rightarrow$ Sprint Backlog $\rightarrow$ Sprint (avec Daily Scrum) $\rightarrow$ Incrément $\rightarrow$ Sprint Review et Sprint Retrospective.}

Le cadre étant posé, nous décrivons maintenant son application concrète à notre projet, en commençant par la vision du produit qui a précédé les sprints.

\subsection{Application de Scrum au projet}

L'application de Scrum au projet s'est articulée autour d'une vision du produit partagée, d'une répartition claire des rôles, des artefacts et des événements prévus par le Scrum Guide, et d'une planification en sprints de durée égale.

\subsubsection{Vision du projet}

Avant le premier sprint, une phase de cadrage a permis de construire, avec le Product Owner, une vision commune du produit : étude du sujet et de l'existant, identification des utilisateurs, définition des objectifs et premières orientations techniques. Cette phase n'est pas un sprint : elle ne produit pas d'incrément logiciel, mais un livrable de cadrage, le \textbf{Project Vision Statement}, qui guide l'ensemble du développement.

\encadre{Project Vision Statement}{\itshape Pour les équipes data et les responsables métier des organisations du secteur de l'assurance, qui doivent auditer rapidement la qualité de leurs données sensibles, \projet{} est une plateforme web d'évaluation intelligente de la qualité des données qui profile automatiquement chaque table, explique le résultat en langage naturel, surveille des règles, corrige sur une copie et répond aux questions des utilisateurs. Contrairement aux outils de profilage, qui ne font que mesurer, et aux outils de validation, qui exigent des règles écrites à l'avance, notre produit interprète et recommande sans configuration préalable, en exécutant ses modèles d'IA localement pour garantir la confidentialité des données.}

Cette vision a servi de base à la définition de l'objectif du produit (\textit{Product Goal}) et du backlog produit présenté au chapitre~3.

\subsubsection{Rôles}

Le tableau~\ref{tab:roles} présente la répartition des trois responsabilités définies par Scrum.

\begin{table}[H]
\centering
\begin{tabular}{|L{4cm}|L{10cm}|}
\hline
\thd{Rôle Scrum} & \thd{Affectation} \\ \hline
Product Owner & \encadrantpro{} (Wevioo) \\ \hline
Scrum Master & \encadrantpro{} (Wevioo) \\ \hline
Développeurs & \etudiant \\ \hline
\end{tabular}
\caption{Répartition des rôles Scrum}
\label{tab:roles}
\end{table}

\subsubsection{Artefacts et engagements}

Conformément au Scrum Guide, chaque artefact est associé à un engagement qui en fixe l'objectif :
\begin{itemize}
  \item le \textbf{Product Backlog}, associé au \textbf{Product Goal} « disposer d'une plateforme locale qui évalue, explique, corrige et permet d'interroger la qualité de ses données », regroupe sous forme d'epics et de récits utilisateurs l'ensemble des besoins fonctionnels ; il est présenté au chapitre~3 ;
  \item le \textbf{Sprint Backlog}, associé à un \textbf{Sprint Goal}, détaille pour chaque sprint les récits retenus et les tâches associées ; il est présenté en tête de la section consacrée au sprint, aux chapitres~5 à~7 ;
  \item l'\textbf{Incrément}, associé à une \textbf{Definition of Done}, est la partie de la plateforme rendue opérationnelle à la fin de chaque sprint. Un récit est considéré comme terminé lorsqu'il est fonctionnel de bout en bout, testé sur les jeux de données de démonstration, documenté et présenté lors de la revue de sprint.
\end{itemize}

\subsubsection{Événements}

Chaque sprint a débuté par un \textit{Sprint Planning} avec le Product Owner pour fixer le Sprint Goal et sélectionner les récits. Des points de synchronisation, trois fois par semaine, ont joué le rôle de \textit{Daily Scrum}. Chaque sprint s'est terminé par une \textit{Sprint Review}, sous forme de démonstration de l'incrément, puis par une \textit{Sprint Retrospective}, qui a permis d'ajuster la suite du travail.

\subsubsection{Planification des sprints}

Le stage s'est déroulé du 23 février au 23 août 2026, soit vingt-six semaines. Le développement a été organisé en sept sprints de trois semaines, précédés de deux semaines de vision. La figure~\ref{fig:gantt} présente cette planification sous forme de diagramme de Gantt et le tableau~\ref{tab:calendrier} en donne le calendrier.

\figplaceholder[1]{gantt_planning.png}{Diagramme de Gantt de la planification du projet}{fig:gantt}{Diagramme de Gantt (généré).}

\begin{table}[H]
\centering
\small
\begin{tabular}{|L{7.4cm}|C{2.6cm}|L{4cm}|}
\hline
\thd{Phase} & \thd{Semaines} & \thd{Période (2026)} \\ \hline
Vision du projet & 1 à 2 & 23 févr. -- 6 mars \\ \hline
Sprint 1 : ingestion et profilage & 3 à 5 & 9 mars -- 27 mars \\ \hline
Sprint 2 : application web et exécution suivie & 6 à 8 & 30 mars -- 17 avr. \\ \hline
Sprint 3 : contexte, rôles et dimensions & 9 à 11 & 20 avr. -- 8 mai \\ \hline
Sprint 4 : règles de qualité et alertes & 12 à 14 & 11 mai -- 29 mai \\ \hline
Sprint 5 : corrections & 15 à 17 & 1 juin -- 19 juin \\ \hline
Sprint 6 : graphe Neo4j & 18 à 20 & 22 juin -- 10 juil. \\ \hline
Sprint 7 : questions en langage naturel & 21 à 23 & 13 juil. -- 31 juil. \\ \hline
\end{tabular}
\caption{Calendrier des sprints}
\label{tab:calendrier}
\end{table}


\section*{Conclusion}
\addcontentsline{toc}{section}{Conclusion}
Ce chapitre nous a permis de présenter l'organisme d'accueil et de situer le projet dans son contexte. L'étude des solutions existantes, menée selon des critères explicites, a mis en évidence l'absence d'une solution combinant profilage automatique, interprétation intelligente, vision transverse et confidentialité, ce qui justifie la solution proposée. Nous avons enfin présenté la démarche agile adoptée. Le chapitre suivant présente l'état de l'art et les choix technologiques sur lesquels repose cette solution.

\chapter{État de l'Art et Étude Technologique}

\section*{Introduction}
\addcontentsline{toc}{section}{Introduction}
Ce chapitre présente les fondements théoriques sur lesquels repose notre solution. Nous définissons d'abord la notion de qualité des données et ses dimensions, puis les techniques de profilage et les mesures statistiques mobilisées. Nous abordons ensuite les grands modèles de langage, leur exécution locale et les techniques permettant de fiabiliser leurs sorties. Nous présentons les concepts d'orchestration et de conteneurisation, puis ceux d'application web sécurisée, de graphe de connaissances et de génération augmentée par un graphe (GraphRAG). Nous justifions enfin nos choix technologiques à travers des études comparatives.

%==============================================================================
\section{La qualité des données}
%==============================================================================

La qualité des données est l'objet même de notre projet. Nous en donnons d'abord une définition, puis nous présentons les dimensions qui permettent de la mesurer et les enjeux qu'elle représente.

\subsection{Définition}

La qualité des données se définit généralement comme l'\textbf{aptitude des données à l'usage} (\textit{fitness for use}) : une donnée est de bonne qualité si elle permet de répondre correctement au besoin pour lequel elle est exploitée~\cite{wang}. Cette définition souligne le caractère relatif de la qualité : un même jeu de données peut être suffisant pour un tableau de bord agrégé et insuffisant pour une campagne de relance client nominative. Le référentiel DAMA-DMBOK~\cite{dama} fait de la gestion de la qualité des données l'une des fonctions centrales de la gouvernance des données, et la norme ISO/IEC 25012~\cite{iso25012} en propose un modèle formel.

Cette notion générale doit cependant être déclinée en propriétés mesurables, appelées dimensions de la qualité.

\subsection{Dimensions de la qualité des données}

La qualité des données est évaluée selon plusieurs dimensions complémentaires. Le tableau~\ref{tab:dimensions} présente celles qui sont les plus couramment retenues et la manière dont elles sont mesurées dans notre solution.

\begin{table}[H]
\centering
\small
\begin{tabular}{|L{2.6cm}|L{5.6cm}|L{5.8cm}|}
\hline
\thd{Dimension} & \thd{Définition} & \thd{Indicateurs exploités par \projet} \\ \hline
Complétude & Proportion de valeurs renseignées par rapport aux valeurs attendues. & Pourcentage de valeurs manquantes par colonne ; détection des colonnes quasi vides. \\ \hline
Unicité & Absence de doublons d'enregistrements ou d'identifiants. & Pourcentage de valeurs distinctes ; détection des colonnes de type identifiant ; lignes dupliquées. \\ \hline
Validité & Conformité des valeurs à un format, un type ou un domaine attendu. & Type inféré vs type attendu ; dates non analysables ; valeurs numériques stockées comme texte. \\ \hline
Cohérence & Absence de contradiction entre valeurs, colonnes ou sources. & Associations fortes entre colonnes (redondances) ; incohérences inter-tables relevées par l'analyse globale. \\ \hline
Exactitude & Conformité des valeurs à la réalité qu'elles décrivent. & Valeurs extrêmes, asymétrie, aplatissement, proportion de zéros. \\ \hline
Actualité & Fraîcheur des données au regard de leur usage. & Bornes temporelles des colonnes de dates. \\ \hline
\end{tabular}
\caption{Principales dimensions de la qualité des données}
\label{tab:dimensions}
\end{table}

Mesurer ces dimensions n'a de sens que si l'on en comprend la portée pour l'organisation, ce que nous précisons ci-après.

\subsection{Enjeux de la qualité des données}

Une donnée de mauvaise qualité a un coût direct (erreurs opérationnelles, relances inutiles, rapprochements manuels) et un coût indirect (décisions erronées, perte de confiance dans les tableaux de bord, non-conformité réglementaire). Dans le secteur de l'assurance, elle peut affecter la tarification, le calcul des provisions ou la connaissance client. C'est pourquoi l'évaluation de la qualité doit intervenir le plus tôt possible, idéalement dès l'arrivée de nouvelles données, et de façon systématique. Le profilage des données constitue la première étape de cette évaluation.

%==============================================================================
\section{Le profilage des données}
%==============================================================================

Le profilage est la technique qui permet de mesurer les dimensions présentées ci-dessus. Nous en présentons la définition, puis trois notions centrales pour notre solution : l'inférence du type sémantique, les mesures d'association et les statistiques de distribution.

\subsection{Définition et objectifs}

Le \textbf{profilage des données} (\textit{data profiling}) désigne l'ensemble des traitements qui calculent des métadonnées et des statistiques descriptives à partir d'un jeu de données, afin d'en comprendre la structure, le contenu et la qualité~\cite{abedjan}. On distingue trois niveaux d'analyse :
\begin{itemize}
  \item \textbf{l'analyse d'une colonne} : type, cardinalité, valeurs manquantes, distribution, valeurs fréquentes, bornes, quantiles ;
  \item \textbf{l'analyse multi-colonnes} : corrélations et associations, dépendances, redondances ;
  \item \textbf{l'analyse multi-tables} : clés candidates, chevauchements de valeurs et correspondances d'entités entre sources.
\end{itemize}

Le premier niveau repose sur une question préalable : quel est le rôle de chaque colonne ?

\subsection{Inférence du type sémantique}

Le type physique d'une colonne (entier, chaîne de caractères, etc.) ne suffit pas à l'interpréter. Une colonne d'entiers peut représenter une mesure continue, une catégorie codée, un indicateur booléen ou un identifiant. L'\textbf{inférence du type sémantique} consiste à déterminer ce rôle par des heuristiques : ratio de valeurs distinctes, nombre de modalités, capacité à être lue comme une date, présence de mots-clés dans le nom de la colonne. Ce type conditionne les statistiques pertinentes : moyenne et écart-type pour une mesure, fréquences pour une catégorie, bornes temporelles pour une date, unicité pour un identifiant.

Une fois le rôle de chaque colonne établi, il devient possible d'étudier les relations entre colonnes.

\subsection{Mesures d'association}

Pour détecter les redondances et les dépendances entre colonnes, trois mesures d'association sont utilisées selon la nature des variables~\cite{sweetviz}.

\paragraph{Coefficient de corrélation de Pearson (numérique / numérique).} Il mesure l'intensité de la relation linéaire entre deux variables continues $X$ et $Y$ :
\begin{equation}
r_{XY} = \frac{\operatorname{cov}(X,Y)}{\sigma_X\,\sigma_Y} \in [-1, 1].
\end{equation}

\paragraph{Rapport de corrélation $\eta$ (catégorielle / numérique).} Il mesure la part de la variance d'une variable numérique $Y$ expliquée par les modalités $x$ d'une variable catégorielle $X$ :
\begin{equation}
\eta^2 = \frac{\sum_{x} n_x\,(\bar{y}_x - \bar{y})^2}{\sum_{i} (y_i - \bar{y})^2} \in [0, 1].
\end{equation}

\paragraph{Coefficient d'incertitude de Theil $U$ (catégorielle / catégorielle).} Fondé sur l'entropie de Shannon $H$, il mesure la part d'incertitude sur $X$ levée par la connaissance de $Y$. Il est asymétrique, ce qui permet de détecter des dépendances directionnelles (un code postal détermine une ville, mais pas l'inverse) :
\begin{equation}
U(X \mid Y) = \frac{H(X) - H(X \mid Y)}{H(X)} \in [0, 1].
\end{equation}

Une association proche de 1 signale souvent une redondance ou une dépendance fonctionnelle, que l'analyse de qualité doit mettre en évidence. Ces mesures relationnelles sont complétées par des indicateurs propres à chaque colonne numérique.

\subsection{Statistiques de distribution}

Au-delà de la moyenne, de l'écart-type et des quantiles, deux indicateurs de forme sont utiles pour repérer des anomalies : le \textbf{coefficient d'asymétrie} (\textit{skewness}), qui signale une distribution tirée vers des valeurs extrêmes, et le \textbf{coefficient d'aplatissement} (\textit{kurtosis}), qui signale des queues de distribution épaisses, souvent liées à des valeurs aberrantes. L'\textbf{écart interquartile} $\text{IQR} = Q_3 - Q_1$ sert de base à la détection des valeurs atypiques, et le \textbf{coefficient de variation} $\text{CV} = \sigma / \mu$ permet de comparer la dispersion de variables d'échelles différentes.

Le profilage fournit ainsi des mesures fiables, mais il ne les interprète pas. C'est le rôle que nous confions aux grands modèles de langage, présentés dans la section suivante.

%==============================================================================
\section{Les grands modèles de langage}
%==============================================================================

Les grands modèles de langage constituent la composante d'intelligence artificielle de notre solution. Nous en rappelons le principe, puis nous abordons leur exécution locale, les familles de modèles utilisées, les paramètres de génération, les techniques d'ingénierie des instructions et leurs limites.

\subsection{Principe}

Les \textbf{grands modèles de langage} (LLM) sont des réseaux de neurones de très grande taille, fondés sur l'architecture \textit{Transformer}~\cite{vaswani} et entraînés sur d'immenses corpus de textes pour prédire le jeton suivant d'une séquence. Ce pré-entraînement, complété par un ajustement sur des instructions, leur confère des capacités générales de compréhension, de raisonnement et de génération~\cite{brown}. Des travaux récents ont montré que ces modèles peuvent être appliqués à des tâches de préparation et de nettoyage de données, telles que la détection d'erreurs ou l'appariement d'entités~\cite{narayan}.

Dans notre projet, le LLM ne calcule aucune statistique : ce rôle est confié au profilage, déterministe et vérifiable. Le LLM intervient en aval pour \textbf{interpréter} les statistiques, \textbf{relier} les constats au contexte métier, \textbf{prioriser} les problèmes et \textbf{formuler} des recommandations. Cette répartition limite le risque d'erreurs de calcul du modèle et rend ses conclusions vérifiables. Encore faut-il pouvoir exécuter ces modèles sans exposer les données.

\subsection{Exécution locale des modèles}

L'utilisation de LLM hébergés par des tiers implique de transmettre les données à l'extérieur de l'organisation, ce qui est incompatible avec les exigences de confidentialité des données d'assurance. L'exécution \textbf{locale} des modèles lève cette contrainte. Elle est rendue possible par la disponibilité de modèles ouverts performants et par la \textbf{quantification}, qui réduit la précision des poids du modèle (par exemple sur 4 bits) pour diminuer fortement la mémoire nécessaire, au prix d'une légère perte de qualité. \textbf{Ollama}~\cite{ollama} est un environnement qui simplifie le téléchargement, la gestion et le service de ces modèles sur un poste ou un serveur, et qui permet de contraindre la sortie du modèle à un format structuré.

Parmi les modèles ouverts exécutables localement, notre plateforme n'est liée à aucun modèle en particulier : le modèle est choisi par la variable de configuration \texttt{OLLAMA\_MODEL}, ce qui permet de changer de modèle sans modifier le code.

\subsection{Les modèles utilisés}

La famille \textbf{Qwen}, développée par Alibaba Cloud~\cite{qwen25}, a servi de référence pour le profilage intelligent : \textbf{Qwen 2.5 14B Instruct}, un modèle de 14 milliards de paramètres ajusté sur des instructions et réputé pour sa capacité à respecter un format de sortie, est la valeur par défaut du traitement par table. Le service de questions et la correction en langage naturel utilisent le modèle présent dans Ollama (par exemple \texttt{gemma4:cloud} lors des évaluations). Trois modèles sont comparés sur un jeu de questions rejouable au chapitre~\ref{chap:benchmark}.

Quel que soit le modèle choisi, son comportement dépend de paramètres de génération que nous présentons à présent.

\subsection{Paramètres de génération}

La génération d'un LLM est un processus d'échantillonnage contrôlé par plusieurs paramètres :
\begin{itemize}
  \item la \textbf{température} règle le caractère aléatoire de la génération : une valeur faible rend les réponses plus déterministes, une valeur élevée plus variées ;
  \item le paramètre \textbf{top-p} (\textit{nucleus sampling}) restreint l'échantillonnage aux jetons dont la probabilité cumulée atteint $p$ ;
  \item la \textbf{longueur maximale} borne le nombre de jetons générés, et donc la durée de la réponse.
\end{itemize}

Ces paramètres règlent la manière dont le modèle répond ; le contenu de sa réponse dépend, lui, des instructions qui lui sont fournies.

\subsection{Ingénierie des instructions et sorties structurées}

L'\textbf{ingénierie des instructions} (\textit{prompt engineering}) regroupe les techniques de formulation des consignes données au modèle. Pour notre usage, les principes suivants sont déterminants : attribuer un \textbf{rôle} explicite au modèle (« expert en qualité des données ») ; fournir un \textbf{contexte} structuré et compact ; décomposer la tâche en \textbf{sections} ; imposer un \textbf{schéma de sortie} explicite ; \textbf{prioriser} certains constats pour orienter l'attention du modèle.

Une sortie structurée est indispensable pour que la réponse puisse être exploitée automatiquement (tableaux de bord, stockage, transmission à une application). Elle doit néanmoins être \textbf{validée} : un modèle peut produire une structure invalide, l'encapsuler dans du texte ou omettre des champs. Au-delà de la forme, le contenu même de la réponse peut être erroné, ce qui constitue la principale limite des LLM.

\subsection{Limites : les hallucinations}

Les LLM peuvent produire des contenus plausibles mais faux, appelés \textbf{hallucinations}~\cite{ji} : colonnes inexistantes, relations inventées, chiffres approximatifs. Dans un contexte de qualité des données, ce risque est critique, puisque l'outil doit précisément détecter des erreurs. Plusieurs mesures permettent de le réduire : ancrer le modèle sur des statistiques calculées de manière déterministe, lui fournir l'inventaire exhaustif des colonnes, exiger des preuves pour chaque constat, limiter la taille du contexte et prévoir une validation humaine. Le chapitre~\ref{chap:benchmark} observe, sur trois modèles, ce type d'erreur dans les réponses du service de questions.

Disposer d'un profilage fiable et d'un modèle maîtrisé ne suffit pas : il faut encore enchaîner ces traitements de manière fiable et en garder la trace, ce qui relève de l'orchestration et de l'observabilité.

%==============================================================================
\section{Orchestration, conteneurisation et observabilité}
%==============================================================================

Cette section présente les trois concepts qui permettent d'industrialiser un pipeline d'IA : l'orchestration, la conteneurisation et l'observabilité.

\subsection{Orchestration de flux de données}

Un pipeline de données enchaîne des tâches dépendantes les unes des autres. Un \textbf{orchestrateur} se charge de les exécuter dans le bon ordre, de les planifier, de les relancer en cas d'échec et d'en conserver l'historique. \textbf{Apache Airflow}~\cite{airflow} modélise un pipeline sous forme de \textbf{DAG} (\textit{Directed Acyclic Graph}), un graphe orienté acyclique dont les nœuds sont des tâches et les arcs des dépendances, et propose une interface de suivi de chaque exécution. Pour que ces exécutions soient identiques sur tous les environnements, l'orchestrateur et ses services sont généralement conteneurisés.

\subsection{Conteneurisation}

La \textbf{conteneurisation} consiste à empaqueter une application avec toutes ses dépendances dans une image exécutable de manière identique sur n'importe quel hôte. \textbf{Docker} en est la technologie de référence, et \textbf{Docker Compose}~\cite{docker} permet de décrire un ensemble de services interconnectés et de les démarrer en une seule opération. Une fois le pipeline exécuté de manière reproductible, il reste à en observer le fonctionnement, en particulier celui du modèle de langage.

\subsection{Observabilité et traçabilité}

Un système qui s'appuie sur un modèle de langage doit rendre visible ce que le modèle a produit et pourquoi. Dans notre plateforme, cette traçabilité repose sur quatre mécanismes simples et locaux : l'historique et les journaux de chaque exécution dans Airflow, un \textbf{statut du modèle} enregistré avec chaque analyse (\texttt{ok}, \texttt{unavailable}, \texttt{invalid}, \texttt{disabled}, \texttt{pending}), le \textbf{versionnement} des jeux de données et de leurs résultats, et, pour les questions, la restitution de la requête exécutée et des lignes obtenues avec chaque réponse.

Ces fondements posés, nous pouvons justifier les technologies retenues pour chacune des briques de la solution.

%==============================================================================
\section{Applications web, graphes de connaissances et GraphRAG}
%==============================================================================

Au-delà du traitement des données, la plateforme est une application web utilisée par plusieurs profils. Cette section présente les notions qui la fondent : l'architecture d'une application web sécurisée, les bases de données de graphe, puis l'interrogation d'un graphe en langage naturel.

\subsection{Applications web, API REST et temps réel}

Une application web moderne sépare l'\textbf{interface}, exécutée dans le navigateur, du \textbf{serveur applicatif}, qui expose une \textbf{API REST}~\cite{fielding} : des ressources (jeux de données, versions, règles) manipulées par les verbes HTTP. L'authentification sans état repose sur des \textbf{jetons JWT} signés~\cite{jwt}, présentés à chaque appel. Pour suivre une exécution longue sans interroger le serveur en boucle, celui-ci peut pousser des événements au navigateur par \textbf{Server-Sent Events} (SSE), le navigateur revenant à une interrogation périodique si le flux est interrompu.

\subsection{Bases de données de graphe}

Une base de données de \textbf{graphe} stocke des \textbf{nœuds} (les entités), des \textbf{relations} entre ces nœuds et des \textbf{propriétés} sur les uns et les autres. Elle est adaptée aux données fortement reliées, où l'on cherche des chemins ou des voisinages plutôt que des agrégats par table. \textbf{Neo4j}~\cite{neo4j} est la base de graphe la plus répandue ; elle s'interroge avec le langage déclaratif \textbf{Cypher}, dont la clause \texttt{MATCH} décrit un motif de nœuds et de relations. Un ensemble de fichiers tabulaires peut être converti en graphe : une ligne devient un nœud, une colonne de référence vers une autre table devient une relation.

\subsection{Génération augmentée par la recherche et par le graphe}

La \textbf{génération augmentée par la recherche} (\textit{Retrieval-Augmented Generation}, RAG)~\cite{lewis} consiste à fournir au modèle, avec la question, des informations retrouvées dans une source de référence, afin que sa réponse s'appuie sur elles plutôt que sur sa mémoire. Le \textbf{GraphRAG}~\cite{edge} applique ce principe à un graphe de connaissances : le modèle traduit la question en requête Cypher (approche \textit{text-to-Cypher}), la requête est exécutée, et la réponse est rédigée à partir des seules lignes obtenues. Comme le modèle écrit la requête, celle-ci ne peut pas être de confiance : elle doit être vérifiée (lecture seule, périmètre limité, plan d'exécution contrôlé) avant d'être exécutée. Nous détaillons ces garde-fous au chapitre~4.

Ces fondements posés, nous pouvons justifier les technologies retenues pour chacune des briques de la solution.

%==============================================================================
\section{Étude comparative et choix technologiques}
%==============================================================================

Pour chaque brique de la solution, nous avons comparé les principales alternatives selon des critères adaptés à notre contexte, en particulier la confidentialité, la facilité de déploiement et la maturité de l'écosystème.

\subsection{Choix de l'outil de profilage}

\begin{table}[H]
\centering
\small
\begin{tabular}{|L{3.8cm}|L{3.2cm}|L{3.2cm}|L{3.4cm}|}
\hline
\thd{Critère} & \thd{SweetViz} & \thd{ydata-profiling} & \thd{Module sur mesure} \\ \hline
Rapport visuel prêt à l'emploi & Oui & Oui & Non \\ \hline
Associations entre types mixtes & Oui ($r$, $\eta$, $U$) & Oui & Oui ($r$, $\eta$, $U$) \\ \hline
Export structuré des statistiques & Indirect & Natif & Natif \\ \hline
Types sémantiques & Partiels & Oui & Oui, paramétrables \\ \hline
Performance sur gros fichiers & Moyenne & Faible à moyenne & Maîtrisée \\ \hline
\end{tabular}
\caption{Comparaison des outils de profilage}
\label{tab:cmp-profiling}
\end{table}

\noindent\textbf{Choix retenu :} \textbf{SweetViz}, pour la richesse de ses mesures d'association, complété par un traitement sur mesure (rôles des colonnes, formats, liens entre tables) qui produit les mesures que la plateforme exploite. L'interface de restitution est construite par l'application à partir des résultats structurés. Le profilage établi, il fallait choisir l'environnement d'exécution des modèles de langage.

\subsection{Choix de l'environnement d'exécution des LLM}

\begin{table}[H]
\centering
\small
\begin{tabular}{|L{3.6cm}|L{2.5cm}|L{2.5cm}|L{2.5cm}|L{2.6cm}|}
\hline
\thd{Critère} & \thd{Ollama} & \thd{llama.cpp} & \thd{vLLM} & \thd{API cloud} \\ \hline
Données traitées localement & Oui & Oui & Oui & Non \\ \hline
Facilité d'installation & Très élevée & Moyenne & Faible (serveur GPU) & Très élevée \\ \hline
Interface de service intégrée & Oui & Oui & Oui & Oui \\ \hline
Sortie structurée contrainte & Oui & Oui & Oui & Selon fournisseur \\ \hline
Gestion du catalogue de modèles & Intégrée & Manuelle & Manuelle & -- \\ \hline
Coût & Gratuit & Gratuit & Gratuit & Payant à l'usage \\ \hline
\end{tabular}
\caption{Comparaison des environnements d'exécution de LLM}
\label{tab:cmp-runtime}
\end{table}

\noindent\textbf{Choix retenu :} \textbf{Ollama}, pour sa simplicité de déploiement, son support de la sortie structurée et sa gestion intégrée des modèles quantifiés ; les API cloud ont été écartées pour des raisons de confidentialité. Restait à choisir l'orchestrateur du pipeline.

\subsection{Choix de l'orchestrateur}

\begin{table}[H]
\centering
\small
\begin{tabular}{|L{3.8cm}|L{2.5cm}|L{2.5cm}|L{2.5cm}|L{2.4cm}|}
\hline
\thd{Critère} & \thd{Airflow} & \thd{Prefect} & \thd{Dagster} & \thd{Planificateur système} \\ \hline
Maturité et communauté & Très élevée & Élevée & Moyenne & -- \\ \hline
Interface de suivi & Oui & Oui & Oui & Non \\ \hline
Reprise sur erreur & Oui & Oui & Oui & Manuelle \\ \hline
Déclenchement par API REST & Oui & Oui & Oui & Non \\ \hline
Adoption en entreprise & Très large & Moyenne & Croissante & -- \\ \hline
\end{tabular}
\caption{Comparaison des orchestrateurs}
\label{tab:cmp-orch}
\end{table}

\noindent\textbf{Choix retenu :} \textbf{Apache Airflow}, standard de fait dans l'industrie, déjà répandu dans les environnements des clients et déclenchable par son API REST depuis le serveur applicatif. Le serveur applicatif lui-même est le choix suivant.

\subsection{Choix du serveur applicatif}

\begin{table}[H]
\centering
\small
\begin{tabular}{|L{3.8cm}|L{3.3cm}|L{3.3cm}|L{3.4cm}|}
\hline
\thd{Critère} & \thd{Spring Boot} & \thd{FastAPI} & \thd{Django REST} \\ \hline
Sécurité intégrée (authentification, rôles) & Spring Security & À assembler & Oui \\ \hline
Gestion des transactions et de la persistance & JPA / Hibernate & À assembler & ORM intégré \\ \hline
Typage et robustesse sur un gros code & Fort (Java) & Moyen & Moyen \\ \hline
Adoption en entreprise & Très large & Croissante & Large \\ \hline
Tests d'intégration & Très outillés & Outillés & Outillés \\ \hline
\end{tabular}
\caption{Comparaison des serveurs applicatifs}
\label{tab:cmp-backend}
\end{table}

\noindent\textbf{Choix retenu :} \textbf{Spring Boot}~\cite{spring}, qui apporte la sécurité, les transactions et la persistance dont le serveur a besoin (versions, règles, résultats), dans un écosystème très répandu chez les clients de Wevioo. Les traitements Python restent dans le worker Airflow et dans le service de questions. L'interface est le choix suivant.

\subsection{Choix du framework de l'interface}

\begin{table}[H]
\centering
\small
\begin{tabular}{|L{3.8cm}|L{3.3cm}|L{3.3cm}|L{3.4cm}|}
\hline
\thd{Critère} & \thd{Angular} & \thd{React} & \thd{Vue} \\ \hline
Cadre complet (routage, formulaires, HTTP) & Oui & À assembler & Partiel \\ \hline
Typage statique natif & TypeScript & Optionnel & Optionnel \\ \hline
Gestion d'état réactive & Signaux & Hooks & Réactivité native \\ \hline
Cohérence sur un projet long & Forte (conventions) & Variable & Moyenne \\ \hline
\end{tabular}
\caption{Comparaison des frameworks d'interface}
\label{tab:cmp-front}
\end{table}

\noindent\textbf{Choix retenu :} \textbf{Angular}~\cite{angular} (composants autonomes et signaux), pour son cadre complet et ses conventions strictes, adaptés à une application aux nombreux écrans. Le dernier choix concerne la base de graphe.

\subsection{Choix de la base de graphe}

\begin{table}[H]
\centering
\small
\begin{tabular}{|L{3.8cm}|L{3.3cm}|L{3.3cm}|L{3.4cm}|}
\hline
\thd{Critère} & \thd{Neo4j} & \thd{Memgraph} & \thd{ArangoDB} \\ \hline
Langage de requête & Cypher & Cypher & AQL \\ \hline
Édition communautaire gratuite & Oui & Oui & Oui \\ \hline
Outils d'exploration visuelle & Neo4j Browser & Memgraph Lab & Interface web \\ \hline
Bibliothèque GraphRAG & \texttt{neo4j-graphrag} & Non officielle & Non officielle \\ \hline
Communauté et documentation & Très large & Moyenne & Moyenne \\ \hline
\end{tabular}
\caption{Comparaison des bases de graphe}
\label{tab:cmp-graph}
\end{table}

\noindent\textbf{Choix retenu :} \textbf{Neo4j}, pour Cypher, son navigateur d'exploration et sa bibliothèque \texttt{neo4j-graphrag} qui fournit les composants de génération de requêtes. Le tableau suivant récapitule l'ensemble de ces choix.

\subsection{Récapitulatif de l'environnement technologique}

\begin{table}[H]
\centering
\small
\begin{tabular}{|L{3.8cm}|L{4.2cm}|L{6cm}|}
\hline
\thd{Domaine} & \thd{Technologie} & \thd{Rôle dans le projet} \\ \hline
Interface & Angular, TypeScript & Écrans, suivi en temps réel, commandes clavier \\ \hline
Serveur applicatif & Spring Boot, Java 17, Spring Security & API REST, authentification JWT, règles, corrections \\ \hline
Persistance & H2 (MySQL possible), JPA & Jeux de données, versions, règles, résultats \\ \hline
Traitement des données & Python, pandas, SweetViz & Lecture, profilage, rôles, relations \\ \hline
Orchestration & Apache Airflow, PostgreSQL & Pipeline en neuf phases, reprise sur erreur \\ \hline
Intelligence artificielle & Ollama, modèles ouverts, LangGraph & Analyse contrôlée, règles, corrections, questions \\ \hline
Graphe et questions & Neo4j, FastAPI, \texttt{neo4j-graphrag} & Graphe de chaque version, service de questions \\ \hline
Déploiement et qualité & Docker Compose, GitHub Actions & Déploiement reproductible, intégration continue \\ \hline
Gestion de versions & Git & Versionnement du projet \\ \hline
\end{tabular}
\caption{Récapitulatif des technologies retenues}
\label{tab:recap-techno}
\end{table}

\section*{Conclusion}
\addcontentsline{toc}{section}{Conclusion}
Ce chapitre a présenté les notions de qualité et de profilage des données, les principes des grands modèles de langage et de leur exécution locale, les concepts d'orchestration, d'application web, de graphe de connaissances et de GraphRAG. Les études comparatives nous ont permis de justifier une pile technologique entièrement open source et exécutable localement. Le chapitre suivant est consacré à l'analyse et à la spécification des besoins de la solution.

\chapter{Analyse et Spécification des Besoins}

\section*{Introduction}
\addcontentsline{toc}{section}{Introduction}
L'analyse des besoins fixe ce que la solution doit faire et les qualités qu'elle doit présenter. Dans ce chapitre, nous identifions d'abord les acteurs du système et formalisons leurs interactions par un diagramme de cas d'utilisation global. Nous exprimons ensuite les besoins fonctionnels à travers le backlog produit, organisé en epics, puis nous structurons les besoins non fonctionnels à l'aide d'un arbre FCM. Nous détaillons enfin les principaux cas d'utilisation et la répartition des récits utilisateurs entre les sprints.

%==============================================================================
\section{Identification des acteurs et des cas d'utilisation}
%==============================================================================

Toute spécification commence par l'identification de ceux qui interagissent avec le système. Nous présentons donc les acteurs de \projet{}, puis la vue d'ensemble de leurs interactions sous forme de diagramme de cas d'utilisation.

\subsection{Identification des acteurs}

Un acteur représente un rôle joué par une personne ou un système externe qui interagit avec la solution. Le tableau~\ref{tab:acteurs} présente les acteurs de \projet{}, en distinguant l'acteur principal, qui sollicite le système, des acteurs secondaires, que le système sollicite.

\begin{table}[H]
\centering
\small
\begin{tabular}{|L{3.4cm}|L{1.9cm}|L{8.7cm}|}
\hline
\thd{Acteur} & \thd{Type} & \thd{Responsabilités} \\ \hline
Administrateur (analyste qualité, ingénieur de données) & Principal & S'authentifie, dépose les fichiers sources avec leur contexte métier, suit l'exécution en temps réel, consulte les résultats par table et pour l'ensemble du jeu de données, définit et surveille des règles de qualité, corrige les données sur une copie, construit le graphe et pose des questions en langage naturel. \\ \hline
Serveur applicatif (Spring Boot) & Secondaire & Reçoit les résultats du traitement par des appels sécurisés, les persiste, évalue les règles et les restitue à l'interface. \\ \hline
Ordonnanceur Airflow & Secondaire & Exécute les phases du pipeline de profilage, dans l'ordre et avec reprise sur erreur. \\ \hline
Serveur LLM (Ollama) & Secondaire & Exécute localement les modèles de langage qui interprètent les mesures, rédigent les corrections et répondent aux questions. \\ \hline
Base de graphe (Neo4j) & Secondaire & Stocke le graphe de chaque version et l'interroge en lecture. \\ \hline
Service de questions & Secondaire & Trie les questions, génère et contrôle les requêtes Cypher, rédige les réponses. \\ \hline
\end{tabular}
\caption{Acteurs du système \projet}
\label{tab:acteurs}
\end{table}

Un seul rôle humain est défini à ce stade (\texttt{ADMIN}) : l'architecture de sécurité permet d'en ajouter d'autres sans modifier les écrans, et ce point est repris dans les perspectives. Ces acteurs étant identifiés, leurs interactions avec le système peuvent être formalisées à l'aide du langage UML.

\subsection{Diagramme de cas d'utilisation global}

Le diagramme de cas d'utilisation global (figure~\ref{fig:uc-global}) présente une vue d'ensemble des fonctionnalités offertes par \projet{} et des acteurs qui les sollicitent.

\figplaceholder[0.95]{uc_global.png}{Diagramme de cas d'utilisation global de \projet}{fig:uc-global}{Acteur principal : Administrateur. Cas d'utilisation : S'authentifier ; Déposer des jeux de données et leur contexte ; Suivre l'exécution ; Consulter les résultats ; Définir et surveiller des règles ; Corriger une table ; Construire et supprimer le graphe ; Poser une question ; Comparer deux versions ; Générer le rapport. Acteurs secondaires : Airflow, Ollama, Neo4j, Service de questions.}

Chacune de ces fonctionnalités correspond à un ou plusieurs besoins fonctionnels, que nous formalisons dans le backlog produit.

%==============================================================================
\section{Besoins fonctionnels : le Product Backlog}
%==============================================================================

Conformément au framework Scrum (chapitre~1), les besoins fonctionnels de \projet{} sont exprimés à travers l'artefact prévu à cet effet : le \textbf{Product Backlog}. Nous présentons d'abord son organisation en epics, puis son contenu détaillé.

\subsection{Organisation du backlog en epics}

Un \textbf{epic} est un besoin de haut niveau, trop large pour être réalisé en un seul sprint, qui regroupe plusieurs récits utilisateurs (\textit{user stories}) contribuant à un même objectif. Le backlog de \projet{} est organisé en sept epics, présentés dans le tableau~\ref{tab:epics}.

\begin{table}[H]
\centering
\small
\begin{tabular}{|C{1.1cm}|L{4.6cm}|L{6.6cm}|C{1.5cm}|}
\hline
\thd{Epic} & \thd{Intitulé} & \thd{Objectif} & \thd{Sprint} \\ \hline
E1 & Ingestion et profilage & Lire tout jeu de données et en mesurer la qualité de manière fiable. & 1 \\ \hline
E2 & Application web et exécution suivie & Déposer des fichiers, lancer l'analyse et la suivre en temps réel, de façon sécurisée. & 2 \\ \hline
E3 & Contexte, rôles et dimensions & Comprendre à quoi sert chaque colonne et expliquer le score de qualité. & 3 \\ \hline
E4 & Règles de qualité et alertes & Surveiller des règles à chaque exécution et rendre l'ensemble déployable et testé. & 4 \\ \hline
E5 & Corrections & Corriger les données sans jamais modifier l'original. & 5 \\ \hline
E6 & Graphe de connaissances & Représenter les tables comme des entités et des relations dans Neo4j. & 6 \\ \hline
E7 & Questions en langage naturel & Répondre aux questions sur les données, leur qualité et l'usage de l'outil. & 7 \\ \hline
\end{tabular}
\caption{Epics du Product Backlog}
\label{tab:epics}
\end{table}

Chaque epic est ensuite décomposé en récits utilisateurs, priorisés selon la méthode MoSCoW (\textit{Must}, \textit{Should}, \textit{Could}).

\subsection{Product Backlog détaillé}

\begin{longtable}{|C{1.1cm}|C{1.1cm}|L{8.4cm}|C{1.6cm}|C{1.3cm}|}
\hline
\thd{Epic} & \thd{ID} & \thd{Récit utilisateur} & \thd{Priorité} & \thd{Sprint} \\ \hline
\endfirsthead
\hline
\thd{Epic} & \thd{ID} & \thd{Récit utilisateur} & \thd{Priorité} & \thd{Sprint} \\ \hline
\endhead
\multirow{4}{*}{E1} & US1 & En tant qu'utilisateur, je veux que mes fichiers soient lus quels que soient leur encodage et leur séparateur. & Must & 1 \\ \cline{2-5}
 & US2 & En tant qu'utilisateur, je veux le profil statistique de chaque table : lignes, valeurs manquantes, valeurs distinctes, doublons, quartiles. & Must & 1 \\ \cline{2-5}
 & US3 & En tant qu'utilisateur, je veux connaître les liens entre mes tables, déduits du recouvrement de leurs valeurs. & Should & 1 \\ \cline{2-5}
 & US4 & En tant qu'utilisateur, je veux savoir quelles tables ont changé d'une version à l'autre (nouvelle, modifiée, inchangée, supprimée). & Should & 1 \\ \hline
\multirow{5}{*}{E2} & US5 & En tant qu'administrateur, je veux m'authentifier pour que seules les personnes autorisées accèdent aux données. & Must & 2 \\ \cline{2-5}
 & US6 & En tant qu'utilisateur, je veux déposer un ou plusieurs fichiers CSV, que le système versionne et dont il détecte les dépôts identiques. & Must & 2 \\ \cline{2-5}
 & US7 & En tant qu'utilisateur, je veux suivre l'avancement de l'analyse en temps réel, étape par étape. & Must & 2 \\ \cline{2-5}
 & US8 & En tant qu'utilisateur, je veux qu'une analyse interrompue soit reprise ou signalée clairement, et pouvoir la relancer. & Must & 2 \\ \cline{2-5}
 & US9 & En tant qu'utilisateur, je veux un tableau de bord de tous mes jeux de données, avec score, note et état. & Should & 2 \\ \hline
\multirow{5}{*}{E3} & US10 & En tant qu'utilisateur, je veux indiquer le domaine, l'objectif et ce que représente une ligne, pour orienter l'analyse. & Must & 3 \\ \cline{2-5}
 & US11 & En tant qu'utilisateur, je veux connaître le rôle de chaque colonne (identifiant, courriel, date, montant\ldots) et les indices qui le justifient. & Must & 3 \\ \cline{2-5}
 & US12 & En tant qu'utilisateur, je veux une analyse en langage naturel de chaque table, dont les affirmations sont vérifiées par les mesures. & Must & 3 \\ \cline{2-5}
 & US13 & En tant qu'utilisateur, je veux un score et une note expliqués par quatre dimensions avec leur numérateur et leur dénominateur. & Must & 3 \\ \cline{2-5}
 & US14 & En tant qu'utilisateur, je veux comparer deux versions d'un jeu de données. & Should & 3 \\ \hline
\multirow{5}{*}{E4} & US15 & En tant qu'utilisateur, je veux créer des règles de qualité (valeurs manquantes, unicité, format, ordre de deux dates, score minimal\ldots). & Must & 4 \\ \cline{2-5}
 & US16 & En tant qu'utilisateur, je veux que le système propose des règles à partir de mon objectif, que j'approuve ou écarte. & Should & 4 \\ \cline{2-5}
 & US17 & En tant qu'utilisateur, je veux que les règles soient évaluées à chaque exécution et qu'un échec déclenche une alerte. & Must & 4 \\ \cline{2-5}
 & US18 & En tant qu'utilisateur, je veux un rapport imprimable d'une version, avec la définition de chaque score. & Could & 4 \\ \cline{2-5}
 & US19 & En tant qu'ingénieur, je veux déployer l'ensemble en une commande, avec des tests automatisés exécutés à chaque modification. & Must & 4 \\ \hline
\multirow{3}{*}{E5} & US20 & En tant qu'utilisateur, je veux que le système repère les corrections possibles dans mes données (espaces, marqueurs d'absence, dates, doublons). & Must & 5 \\ \cline{2-5}
 & US21 & En tant qu'utilisateur, je veux décrire une correction en langage naturel et que le modèle la traduise en opérations vérifiées. & Should & 5 \\ \cline{2-5}
 & US22 & En tant qu'utilisateur, je veux un aperçu chiffré, puis appliquer la correction à une copie qui devient une nouvelle version, avec un journal. & Must & 5 \\ \hline
\multirow{3}{*}{E6} & US23 & En tant qu'utilisateur, je veux un modèle de graphe déduit de mes tables (entités, propriétés, relations) que je puisse relire avant la construction. & Must & 6 \\ \cline{2-5}
 & US24 & En tant qu'utilisateur, je veux que les cas courants soient gérés : clé composée, table de liaison, référence vers la même table. & Should & 6 \\ \cline{2-5}
 & US25 & En tant qu'utilisateur, je veux construire et supprimer le graphe d'une version dans Neo4j, avec les conflits et les liens incertains signalés. & Must & 6 \\ \hline
\multirow{4}{*}{E7} & US26 & En tant qu'utilisateur, je veux poser une question en langage naturel et que le système sache si elle porte sur les données, la qualité ou l'usage de l'outil. & Must & 7 \\ \cline{2-5}
 & US27 & En tant qu'utilisateur, je veux que les requêtes générées par le modèle soient contrôlées avant exécution (lecture seule, périmètre de la version). & Must & 7 \\ \cline{2-5}
 & US28 & En tant qu'utilisateur, je veux que chaque réponse indique sa source et, pour le graphe, la requête et les lignes obtenues. & Must & 7 \\ \cline{2-5}
 & US29 & En tant que visiteur, je veux une page d'accueil publique et une interface fluide, avec raccourcis clavier. & Could & 7 \\ \hline
\caption{Product Backlog de \projet}
\label{tab:backlog}
\end{longtable}

Ces récits décrivent ce que le système doit faire. Ils doivent être complétés par les qualités attendues du système, qui font l'objet de la section suivante.

%==============================================================================
\section{Besoins non fonctionnels : arbre FCM}
%==============================================================================

Les besoins non fonctionnels décrivent les qualités attendues du système plutôt que ses fonctionnalités. Pour les structurer, nous nous appuyons sur le modèle \textbf{FCM} (\textit{Factors -- Criteria -- Metrics}) de McCall. Nous justifions d'abord le choix des facteurs retenus, puis nous présentons l'arbre complet.

\subsection{Choix des facteurs de qualité}

Le modèle FCM décompose la qualité d'un logiciel en trois niveaux : des \textbf{facteurs}, perçus par les utilisateurs ; des \textbf{critères}, qui traduisent chaque facteur en propriétés techniques ; et des \textbf{métriques}, qui permettent de vérifier objectivement chaque critère. Parmi les facteurs proposés par McCall et par la norme ISO/IEC 25010, nous n'avons retenu que ceux qui découlent directement des contraintes du projet, présentés dans le tableau~\ref{tab:facteurs}.

\begin{table}[H]
\centering
\small
\begin{tabular}{|L{3cm}|L{11cm}|}
\hline
\thd{Facteur} & \thd{Justification du choix} \\ \hline
Sécurité & Les données d'assurance sont personnelles et sensibles : l'accès doit être authentifié, aucune donnée ne doit quitter l'infrastructure de l'entreprise, et les données sources ne doivent jamais être altérées. C'est la contrainte la plus forte du projet. \\ \hline
Fiabilité & Un outil d'audit de la qualité doit lui-même être fiable : chaque table doit obtenir un résultat exploitable, même lorsque le modèle de langage échoue ou répond de manière incorrecte, et une exécution interrompue ne doit pas rester bloquée. \\ \hline
Efficacité & L'exécution locale de grands modèles est coûteuse en temps et en mémoire ; les résultats doivent apparaître progressivement et l'interface rester fluide. \\ \hline
Maintenabilité & Le périmètre du produit a évolué à chaque sprint : chaque bloc doit pouvoir changer indépendamment et être couvert par des tests automatisés. \\ \hline
Portabilité & La solution doit être déployée chez différents clients et sur différents postes, sans adaptation manuelle. \\ \hline
Exploitabilité & Les résultats d'une IA doivent être vérifiables : l'utilisateur doit pouvoir savoir d'où vient une réponse et ce qui a été modifié dans ses données. \\ \hline
\end{tabular}
\caption{Justification des facteurs de qualité retenus}
\label{tab:facteurs}
\end{table}

Les six facteurs retenus sont déclinés en critères et métriques dans l'arbre présenté ci-après.

\subsection{Arbre FCM}

Le tableau~\ref{tab:fcm} présente l'arbre FCM de \projet{} : chaque facteur est décliné en critères, puis en métriques vérifiables sur la solution réalisée.

\begin{longtable}{|L{2.3cm}|L{3cm}|L{8.7cm}|}
\hline
\thd{Facteur} & \thd{Critère} & \thd{Métrique} \\ \hline
\endfirsthead
\hline
\thd{Facteur} & \thd{Critère} & \thd{Métrique} \\ \hline
\endhead

\multirow{3}{*}{Sécurité} & Authentification & Jeton JWT signé, valable 8~h ; tout appel à l'API (sauf connexion et rappels du pipeline) exige le rôle administrateur. \\ \cline{2-3}
 & Confidentialité & Aucun appel vers un service externe : modèle, graphe et questions sont exécutés localement ; le navigateur n'atteint jamais le service de questions. \\ \cline{2-3}
 & Intégrité des sources & Les corrections s'appliquent à une copie ; les requêtes générées sont en lecture seule et limitées au périmètre d'une version. \\ \hline
\multirow{3}{*}{Fiabilité} & Robustesse & Une table dont l'analyse échoue donne le statut \texttt{PARTIAL}, avec la raison ; une analyse de repli déterministe remplace une réponse invalide du modèle. \\ \cline{2-3}
 & Reprise & Rappels du pipeline réessayés pendant environ 2~min ; au démarrage, les exécutions en cours sont réconciliées avec Airflow ; une surveillance marque \texttt{FAILED} toute exécution bloquée au-delà de 180~min. \\ \cline{2-3}
 & Cohérence & Le téléversement est validé en entier avant toute écriture ; un dépôt identique propose d'ouvrir la version existante. \\ \hline
\multirow{2}{*}{Efficacité} & Réactivité & Résultats publiés table par table dès que leurs métriques sont prêtes ; événements poussés par SSE, avec repli sur une interrogation toutes les 5~s. \\ \cline{2-3}
 & Maîtrise des délais & Requête Cypher limitée à 15~s et à 50 lignes lues ; délais paramétrables pour chaque appel au modèle. \\ \hline
\multirow{2}{*}{Maintenabilité} & Modularité & Cinq blocs indépendants (interface, serveur, worker, service de questions, graphe), chacun avec son README et ses tests. \\ \cline{2-3}
 & Testabilité & 292 tests automatisés (serveur 110, interface 76, worker 67, questions 39), exécutés à chaque modification par l'intégration continue. \\ \hline
\multirow{2}{*}{Portabilité} & Reproductibilité & Déploiement complet par une commande Docker Compose ; image dédiée pour chaque service. \\ \cline{2-3}
 & Configurabilité & Tous les paramètres (adresses, secrets, modèle, délais) sont modifiables par variables d'environnement. \\ \hline
\multirow{2}{*}{Exploitabilité} & Explicabilité & Chaque rôle de colonne garde ses indices ; chaque dimension du score donne son numérateur et son dénominateur ; chaque réponse dit d'où elle vient. \\ \cline{2-3}
 & Auditabilité & Statut du modèle enregistré avec chaque analyse ; journal de chaque correction ; requête et lignes retournées avec chaque réponse du graphe. \\ \hline
\caption{Arbre FCM des besoins non fonctionnels de \projet}
\label{tab:fcm}
\end{longtable}

Les besoins étant désormais spécifiés, tant sur le plan fonctionnel que non fonctionnel, nous pouvons raffiner les principaux cas d'utilisation.

%==============================================================================
\section{Raffinement des cas d'utilisation}
%==============================================================================

Chaque cas d'utilisation principal est raffiné par un diagramme de cas d'utilisation détaillé, suivi de sa description textuelle. Nous commençons par le cas central, le dépôt et l'analyse d'un jeu de données.

\subsection{Cas d'utilisation « Déposer et analyser un jeu de données »}

\figplaceholder[0.9]{uc_deposer.png}{Diagramme de cas d'utilisation détaillé « Déposer et analyser un jeu de données »}{fig:uc-deposer}{Acteur : Administrateur. Cas : Déposer et analyser un jeu de données, avec include Valider les fichiers, Créer la version, Lancer le pipeline, Suivre l'avancement ; extend Renseigner le contexte métier, Relancer une analyse. Acteurs secondaires : Airflow, Ollama.}

\begin{table}[H]
\centering
\small
\begin{tabular}{|L{3.4cm}|L{10.6cm}|}
\hline
\thd{Élément} & \thd{Description} \\ \hline
Cas d'utilisation & Déposer et analyser un jeu de données \\ \hline
Acteurs & Administrateur (principal), Airflow, Serveur LLM \\ \hline
Préconditions & L'administrateur est authentifié ; la plateforme est démarrée. \\ \hline
Postconditions & Une version du jeu de données existe, avec les résultats de chaque table, le score, les dimensions et l'analyse globale. \\ \hline
Scénario nominal &
1. L'administrateur dépose un ou plusieurs fichiers CSV et, s'il le souhaite, renseigne le domaine, l'objectif et ce que représente une ligne.\newline
2. Le système valide tous les fichiers, calcule leur empreinte et crée la version.\newline
3. Le système déclenche le pipeline dans Airflow.\newline
4. Le pipeline profile chaque table, détermine le rôle des colonnes et publie un premier résultat.\newline
5. Le modèle enrichit l'analyse de chaque table ; le résultat est publié de nouveau.\newline
6. Le pipeline propose des règles, produit l'analyse globale et déclare l'exécution terminée.\newline
7. Le système évalue les règles et prévient l'administrateur en cas d'échec. \\ \hline
Scénarios alternatifs &
2a. Un fichier est invalide : il est signalé et rien n'est écrit.\newline
2b. L'envoi est identique à une version existante : le système propose d'ouvrir cette version ou de la relancer.\newline
4a. Une table est illisible : elle est signalée et les autres sont traitées ; le statut final est \texttt{PARTIAL}.\newline
5a. Le modèle est indisponible ou répond de façon invalide : une analyse déterministe est publiée et le statut du modèle l'indique.\newline
*a. L'exécution échoue : l'administrateur peut la relancer. \\ \hline
\end{tabular}
\caption{Description du cas d'utilisation « Déposer et analyser un jeu de données »}
\label{tab:cu-deposer}
\end{table}

Une fois les résultats disponibles, l'utilisateur peut en tirer des règles de qualité qui seront surveillées dans le temps.

\subsection{Cas d'utilisation « Définir et surveiller des règles de qualité »}

\figplaceholder[0.85]{uc_regles.png}{Diagramme de cas d'utilisation détaillé « Définir et surveiller des règles de qualité »}{fig:uc-regles}{Acteur : Administrateur. Cas : Définir et surveiller des règles, avec include Évaluer les règles ; extend Créer une règle, Approuver une règle proposée, Écarter une règle proposée, Suspendre une règle, Vérifier maintenant.}

\begin{table}[H]
\centering
\small
\begin{tabular}{|L{3.4cm}|L{10.6cm}|}
\hline
\thd{Élément} & \thd{Description} \\ \hline
Cas d'utilisation & Définir et surveiller des règles de qualité \\ \hline
Acteur & Administrateur \\ \hline
Préconditions & Au moins une version du jeu de données a été analysée. \\ \hline
Postconditions & Les règles actives sont évaluées à chaque exécution ; un échec produit une alerte. \\ \hline
Scénario nominal &
1. L'administrateur ouvre l'onglet des règles.\newline
2. Il consulte les règles proposées à partir de l'objectif et en approuve certaines.\newline
3. Il crée une règle, en choisissant son type, sa table, sa colonne, son seuil et sa gravité.\newline
4. Il demande une vérification immédiate.\newline
5. Le système affiche, pour chaque règle, le résultat (réussie, en échec, non évaluée) et un message explicatif. \\ \hline
Scénarios alternatifs &
2a. Une règle proposée n'est pas pertinente : l'administrateur l'écarte et elle ne revient pas.\newline
3a. La table visée n'existe pas dans la version : la règle est non évaluée, avec la raison. Si c'est la colonne qui n'existe plus, la règle échoue et le message le dit.\newline
5a. Une règle de gravité critique échoue : elle est comptée à part dans l'alerte. \\ \hline
\end{tabular}
\caption{Description du cas d'utilisation « Définir et surveiller des règles de qualité »}
\label{tab:cu-regles}
\end{table}

Lorsque des règles échouent, l'utilisateur peut corriger les données concernées.

\subsection{Cas d'utilisation « Corriger une table »}

\figplaceholder[0.85]{uc_corriger.png}{Diagramme de cas d'utilisation détaillé « Corriger une table »}{fig:uc-corriger}{Acteur : Administrateur. Cas : Corriger une table, avec include Prévisualiser les corrections, Appliquer sur une copie ; extend Choisir une correction trouvée, Décrire une correction en langage naturel. Acteur secondaire : Ollama.}

\begin{table}[H]
\centering
\small
\begin{tabular}{|L{3.4cm}|L{10.6cm}|}
\hline
\thd{Élément} & \thd{Description} \\ \hline
Cas d'utilisation & Corriger une table \\ \hline
Acteurs & Administrateur, Serveur LLM \\ \hline
Préconditions & La version est analysée ; les fichiers sources sont présents. \\ \hline
Postconditions & Une nouvelle version, dérivée de la première, contient les données corrigées et un journal des changements ; la version d'origine est intacte. \\ \hline
Scénario nominal &
1. L'administrateur choisit une table.\newline
2. Le système liste les corrections qu'il a trouvées (espaces, marqueurs d'absence, formats de date, valeurs écrites de plusieurs façons, doublons).\newline
3. L'administrateur en retient certaines et peut en ajouter une, écrite en langage naturel.\newline
4. Le système affiche un aperçu chiffré : nombre de valeurs touchées, exemples, avertissements.\newline
5. L'administrateur applique ; le système crée la nouvelle version et la fait analyser comme un dépôt. \\ \hline
Scénarios alternatifs &
3a. La demande écrite est trop vague ou vise une colonne inexistante : le système la refuse en expliquant pourquoi.\newline
4a. Une date est ambiguë (\texttt{01/02/2024}) : elle n'est jamais devinée et un avertissement est affiché.\newline
3b. Le modèle est indisponible : seules les corrections trouvées par le système restent proposées. \\ \hline
\end{tabular}
\caption{Description du cas d'utilisation « Corriger une table »}
\label{tab:cu-corriger}
\end{table}

Enfin, l'utilisateur peut interroger ses données sans écrire de requête, ce qui fait l'objet du dernier cas d'utilisation.

\subsection{Cas d'utilisation « Poser une question en langage naturel »}

\figplaceholder[0.85]{uc_questions.png}{Diagramme de cas d'utilisation détaillé « Poser une question en langage naturel »}{fig:uc-questions}{Acteur : Administrateur. Cas : Poser une question, avec include Trier la question ; extend Interroger le graphe, Consulter les résultats de qualité, Consulter le guide d'utilisation. Acteurs secondaires : Service de questions, Neo4j, Ollama.}

\begin{table}[H]
\centering
\small
\begin{tabular}{|L{3.4cm}|L{10.6cm}|}
\hline
\thd{Élément} & \thd{Description} \\ \hline
Cas d'utilisation & Poser une question en langage naturel \\ \hline
Acteurs & Administrateur, Service de questions, Neo4j, Serveur LLM \\ \hline
Préconditions & La version est analysée ; pour une question sur le contenu des données, son graphe est construit. \\ \hline
Postconditions & L'administrateur obtient une réponse qui indique sa source ; pour le graphe, la requête exécutée et les lignes obtenues sont fournies. \\ \hline
Scénario nominal &
1. L'administrateur saisit une question dans l'onglet dédié.\newline
2. Le système la classe : contenu des données, qualité ou usage de l'outil.\newline
3. Pour une question sur les données, le modèle écrit une requête Cypher, qui est contrôlée puis exécutée.\newline
4. Le système rédige la réponse à partir des seules lignes obtenues, des faits de qualité et du guide.\newline
5. L'interface affiche la réponse, sa source, la requête et les lignes. \\ \hline
Scénarios alternatifs &
3a. La requête n'est pas admissible (écriture, autre périmètre) : le modèle la corrige, deux tentatives au plus ; à défaut, la réponse s'appuie sur les faits de qualité et le dit.\newline
3b. Le graphe n'est pas construit : le système demande de le construire.\newline
2a. Le classement est incertain : la question est traitée comme une question de qualité.\newline
*a. Le modèle ou Neo4j est injoignable : un message explicite est renvoyé. \\ \hline
\end{tabular}
\caption{Description du cas d'utilisation « Poser une question en langage naturel »}
\label{tab:cu-questions}
\end{table}

Ces cas d'utilisation couvrent l'ensemble du backlog. Il reste à préciser comment ses récits ont été répartis dans le temps.

%==============================================================================
\section{Répartition du backlog par sprint}
%==============================================================================

La planification présentée sous forme de diagramme de Gantt au chapitre~1 (figure~\ref{fig:gantt}) répartit les epics et les récits utilisateurs entre sept sprints de durée égale, chacun associé à un objectif de sprint (\textit{Sprint Goal}), comme le montre le tableau~\ref{tab:planif}.

\begin{table}[H]
\centering
\small
\begin{tabular}{|L{1.6cm}|L{6.4cm}|C{1.4cm}|L{4.6cm}|}
\hline
\thd{Sprint} & \thd{Sprint Goal} & \thd{Epic} & \thd{Récits utilisateurs} \\ \hline
Sprint 1 & Mesurer automatiquement la qualité de tout jeu de données & E1 & US1 à US4 \\ \hline
Sprint 2 & Déposer des fichiers et suivre l'analyse en temps réel, de façon sécurisée & E2 & US5 à US9 \\ \hline
Sprint 3 & Comprendre les colonnes et expliquer le score & E3 & US10 à US14 \\ \hline
Sprint 4 & Surveiller la qualité par des règles et livrer une plateforme déployable et testée & E4 & US15 à US19 \\ \hline
Sprint 5 & Corriger les données sans toucher à l'original & E5 & US20 à US22 \\ \hline
Sprint 6 & Représenter les tables en graphe dans Neo4j & E6 & US23 à US25 \\ \hline
Sprint 7 & Répondre aux questions en langage naturel & E7 & US26 à US29 \\ \hline
\end{tabular}
\caption{Répartition du backlog par sprint}
\label{tab:planif}
\end{table}

\section*{Conclusion}
\addcontentsline{toc}{section}{Conclusion}
Ce chapitre a permis de formaliser les besoins de la solution : acteurs et diagramme de cas d'utilisation global, besoins fonctionnels organisés en sept epics dans le Product Backlog, besoins non fonctionnels structurés par un arbre FCM justifié, cas d'utilisation détaillés et répartition du backlog par sprint. Le chapitre suivant traduit ces besoins en une conception de la solution.

\chapter{Conception de la Solution}

\section*{Introduction}
\addcontentsline{toc}{section}{Introduction}
Après avoir spécifié les besoins, ce chapitre présente la conception de \projet{}. Nous présentons d'abord une conception globale à l'aide du formalisme UML, qui rapproche les entités du domaine des données de celles de l'intelligence artificielle. Nous décrivons ensuite les architectures logicielle, logique et physique de la plateforme, puis la sécurité qui les traverse. Nous détaillons enfin l'architecture de ses composants d'intelligence artificielle : l'analyse contrôlée par table, les corrections, le graphe de connaissances et le service de questions.

%==============================================================================
\section{Conception globale}
%==============================================================================

La conception globale donne une vue d'ensemble du domaine manipulé par la solution, indépendamment des choix d'implémentation. Elle comprend une vue statique, le diagramme de classes, et une vue dynamique, le diagramme d'activité global d'une exécution.

\subsection{Diagramme de classes global}

Le diagramme de classes de la figure~\ref{fig:classes} modélise les entités manipulées par la plateforme et leurs relations.

\figplaceholder[1]{diag_classes_global.png}{Diagramme de classes global de \projet}{fig:classes}{Diagramme de classes global (généré).}

Ce diagramme s'organise en cinq groupes d'entités, dont le tableau~\ref{tab:classes} précise le rôle et l'analogie avec les concepts classiques d'un système d'IA.

\begin{table}[H]
\centering
\small
\begin{tabular}{|L{2.8cm}|L{4.2cm}|L{7cm}|}
\hline
\thd{Groupe} & \thd{Entités} & \thd{Rôle et analogie avec les concepts d'IA} \\ \hline
Données & Jeu de données, Version, Table de la version, Métrique de colonne & Les \textbf{données d'entrée} du système. Une version est un dépôt daté ; une version corrigée dérive de celle qu'elle corrige. Le contexte métier (domaine, objectif, grain d'une ligne) joue le rôle d'information contextuelle pour le modèle. \\ \hline
Mesures & Métriques, Rôle de colonne, Relation entre tables, Signature de contenu & L'équivalent de l'\textbf{extraction de caractéristiques} : des mesures déterministes qui décrivent chaque colonne, ses indices de rôle, ses liens avec les autres tables et le changement d'une version à l'autre. \\ \hline
Intelligence artificielle & Analyse de table, Analyse globale, Statut du modèle & La chaîne d'\textbf{inférence} : le modèle produit une analyse structurée (problèmes, causes, recommandations), validée contre les mesures ; son statut est conservé. \\ \hline
Qualité & Règle, Résultat de règle, Dimension, Score & La \textbf{surveillance} : des règles manuelles ou proposées, évaluées à chaque exécution, et un score décomposé en dimensions. \\ \hline
Exploitation & Exécution du pipeline, Correction, Graphe, Question & Le \textbf{cycle de vie} : chaque exécution a un statut, chaque correction un journal, chaque graphe un périmètre, chaque question une source. \\ \hline
\end{tabular}
\caption{Groupes d'entités du diagramme de classes et analogies avec les concepts d'IA}
\label{tab:classes}
\end{table}

Cette vue statique décrit ce que manipule la plateforme. La vue dynamique qui suit décrit comment ces entités sont produites au cours d'une exécution.

\subsection{Diagramme d'activité global}

Le diagramme d'activité de la figure~\ref{fig:act-global} décrit le déroulement d'une exécution complète, du dépôt des fichiers à la restitution des résultats, y compris la gestion des échecs.

\figplaceholder[0.8]{act_pipeline_global.png}{Diagramme d'activité global d'une exécution}{fig:act-global}{Diagramme d'activité global (généré).}

Chaque activité de ce diagramme est ensuite détaillée par des diagrammes de séquence et d'activité propres à chaque sprint, aux chapitres~5 à~7. Avant cela, il convient de décrire les architectures qui supportent ces traitements.

%==============================================================================
\section{Architectures de la solution}
%==============================================================================

La plateforme est décrite selon trois points de vue complémentaires : l'architecture logicielle, qui précise le style et les composants ; l'architecture logique, qui organise ces composants en couches ; et l'architecture physique, qui décrit leur déploiement. Nous présentons ensuite la sécurité, qui les traverse.

\subsection{Architecture logicielle}

La plateforme suit un style \textbf{en blocs communiquant par des API} : chaque bloc a une responsabilité unique, s'exécute dans son propre processus et ne connaît des autres que leur interface. La figure~\ref{fig:arch-logicielle} présente les cinq blocs, le serveur de modèles de langage et leurs échanges.

\figplaceholder[1]{arch_logicielle.png}{Architecture logicielle de \projet}{fig:arch-logicielle}{Navigateur, interface Angular, serveur Spring Boot, Airflow (worker), service de questions, Neo4j, Ollama, avec les échanges : REST, rappels du pipeline, SSE, Bolt.}

Le tableau~\ref{tab:blocs} précise la responsabilité de chaque bloc.

\begin{table}[H]
\centering
\small
\begin{tabular}{|L{3.2cm}|L{3.2cm}|L{7.6cm}|}
\hline
\thd{Bloc} & \thd{Technologie} & \thd{Responsabilité} \\ \hline
Interface & Angular, TypeScript & Dépôt, suivi en temps réel, onglets de résultats, tableau de bord, page d'accueil publique. \\ \hline
Serveur applicatif & Spring Boot, Java 17 & Authentification, jeux de données et versions, règles, corrections, graphe, relais des questions, déclenchement et surveillance d'Airflow, événements SSE. \\ \hline
Worker de traitement & Python, Airflow & Lecture et profilage, rôles de colonnes, relations entre tables, analyse par le modèle, règles proposées, analyse globale. \\ \hline
Service de questions & Python, FastAPI & Tri des questions, génération et contrôle des requêtes Cypher, rédaction des réponses. \\ \hline
Base de graphe & Neo4j & Un graphe par version, identifié par une propriété de périmètre. \\ \hline
Serveur de modèles & Ollama & Exécution locale du modèle utilisé par le worker, le serveur et le service de questions. \\ \hline
\end{tabular}
\caption{Blocs de la solution et responsabilités}
\label{tab:blocs}
\end{table}

Un principe structure cet ensemble : le worker ne décide jamais seul de ce qui est publié. Chaque résultat est envoyé au serveur par un appel protégé par un jeton, et les réponses du modèle sont validées contre les mesures avant d'être affichées. Ces blocs peuvent être regroupés selon leur responsabilité, ce qui conduit à l'architecture logique.

\subsection{Architecture logique}

L'architecture logique organise chaque bloc en couches, chacune ne s'appuyant que sur la couche inférieure. La figure~\ref{fig:arch-logique} présente celle du serveur applicatif, le bloc le plus riche.

\figplaceholder[0.9]{arch_logique.png}{Architecture logique en couches du serveur applicatif}{fig:arch-logique}{Couches : contrôleurs REST et sécurité ; services métier (jeux de données, règles, corrections, graphe, questions, dimensions, comparaison) ; domaine et persistance (JPA) ; clients externes (Airflow, Neo4j, Ollama, service de questions). Services transverses : surveillance, reprise, événements.}

Le tableau~\ref{tab:couches} résume le rôle de chaque couche et les autres blocs qui suivent la même logique.

\begin{table}[H]
\centering
\small
\begin{tabular}{|L{3.2cm}|L{10.8cm}|}
\hline
\thd{Bloc} & \thd{Organisation en couches} \\ \hline
Serveur applicatif & Contrôleurs REST $\rightarrow$ services métier $\rightarrow$ domaine et persistance $\rightarrow$ clients externes ; services transverses de surveillance des exécutions bloquées et de reprise au démarrage. \\ \hline
Interface & \texttt{core} (appels HTTP, authentification, événements, état partagé en signaux) $\rightarrow$ \texttt{features} (un dossier par écran) $\rightarrow$ \texttt{shared} (formatage, composants communs, animations). \\ \hline
Worker & Phases du DAG $\rightarrow$ analyseurs (rôles, modèle, règles) et extracteurs (profil, formats) $\rightarrow$ client du serveur et stockage des résultats. \\ \hline
Service de questions & API FastAPI $\rightarrow$ tri et génération $\rightarrow$ garde-fous $\rightarrow$ accès à Neo4j. \\ \hline
\end{tabular}
\caption{Organisation en couches de chaque bloc}
\label{tab:couches}
\end{table}

Cette organisation logique doit enfin être projetée sur une infrastructure concrète, ce que décrit l'architecture physique.

\subsection{Architecture physique}

L'architecture physique (figure~\ref{fig:arch-physique}) décrit le déploiement de la plateforme. Les services sont conteneurisés par Docker Compose et partagent un même réseau et les mêmes secrets ; le serveur de modèles de langage reste sur la machine hôte, afin de bénéficier directement du processeur graphique.

\figplaceholder[1]{arch_physique.png}{Architecture physique (diagramme de déploiement)}{fig:arch-physique}{Conteneurs : frontend (nginx, 4200), backend (8081), airflow-webserver (8080), airflow-scheduler, airflow-init, postgres, neo4j (7474, 7687), graphrag (8090). Sur l'hôte : Ollama (11434), volumes de données.}

Le tableau~\ref{tab:conteneurs} liste les conteneurs et leur rôle.

\begin{table}[H]
\centering
\small
\begin{tabular}{|L{3.8cm}|L{7.8cm}|C{2.2cm}|}
\hline
\thd{Conteneur} & \thd{Rôle} & \thd{Port} \\ \hline
\texttt{frontend} & Interface Angular servie par nginx & 4200 \\ \hline
\texttt{backend} & API Spring Boot et base H2 (volume dédié) & 8081 \\ \hline
\texttt{airflow-webserver} & Interface web d'Airflow & 8080 \\ \hline
\texttt{airflow-scheduler} & Exécution des phases du pipeline & interne \\ \hline
\texttt{airflow-init} & Migration de la base d'Airflow et création du compte administrateur, puis arrêt & aucun \\ \hline
\texttt{postgres} & Base d'Airflow, non publiée sur l'hôte & interne \\ \hline
\texttt{neo4j} & Base de graphe (volume dédié) & 7474, 7687 \\ \hline
\texttt{graphrag} & Service de questions & 8090 \\ \hline
\end{tabular}
\caption{Conteneurs de la plateforme}
\label{tab:conteneurs}
\end{table}

Un script de démarrage permet aussi de lancer chaque service dans sa propre fenêtre pour le développement sous Windows. Reste à décrire la sécurité, qui traverse ces trois vues.

\subsection{Sécurité}

La sécurité repose sur des mesures simples et vérifiables, résumées dans le tableau~\ref{tab:securite}.

\begin{table}[H]
\centering
\small
\begin{tabular}{|L{3.8cm}|L{10.2cm}|}
\hline
\thd{Menace} & \thd{Mesure} \\ \hline
Accès non autorisé & Connexion par identifiant et mot de passe, jeton JWT signé valable 8~h ; toute route de l'API, hors connexion et rappels du pipeline, exige le rôle administrateur. Côté interface, des gardes de routes renvoient un visiteur vers la page de connexion. \\ \hline
Appels falsifiés du pipeline & Les rappels d'Airflow n'ont pas de session : ils exigent un jeton partagé, comparé en temps constant. \\ \hline
Secrets par défaut & Un avertissement est journalisé au démarrage si un secret de développement est utilisé ; l'option de mode strict refuse de démarrer avec eux. \\ \hline
Requêtes générées dangereuses & Le service de questions n'est appelé que par le serveur (jeton dédié) ; toute requête Cypher doit être en lecture seule et limitée au périmètre de la version (voir section~\ref{sec:questions-conception}). \\ \hline
Altération des données & Les corrections s'appliquent à une copie ; les fichiers d'origine ne sont jamais réécrits. \\ \hline
Fuite de données & Le modèle s'exécute localement ; la base du pipeline n'est pas exposée ; les secrets ne sont jamais versionnés. \\ \hline
\end{tabular}
\caption{Menaces et mesures de sécurité}
\label{tab:securite}
\end{table}

Ces architectures décrivent la plateforme dans son ensemble. Ses composants d'intelligence artificielle méritent toutefois une description détaillée.

%==============================================================================
\section{Architecture des composants d'intelligence artificielle}
%==============================================================================

L'intelligence artificielle intervient à quatre endroits : l'analyse des tables, la rédaction des corrections, le choix des noms du graphe et les réponses aux questions. Un même principe les gouverne : \textbf{le code mesure et contrôle, le modèle interprète et rédige}. Nous présentons d'abord ce principe, puis chaque composant.

\subsection{Principes de conception}

Quatre principes ont guidé la conception :
\begin{itemize}
  \item \textbf{Séparer la mesure de l'interprétation} : les statistiques, les rôles de colonnes, les formats et les liens entre tables sont calculés de façon déterministe ; le modèle ne tranche que les cas que les mesures laissent ouverts.
  \item \textbf{Confronter le modèle aux preuves} : chaque affirmation chiffrée est vérifiée contre la colonne qu'elle concerne ; une réponse non fiable est remplacée par une analyse déterministe.
  \item \textbf{Contraindre ce que le modèle peut faire} : il choisit parmi des opérations connues (corrections), propose des règles que le code filtre, et écrit des requêtes que des garde-fous contrôlent.
  \item \textbf{Dire d'où vient chaque résultat} : statut du modèle, indices du rôle des colonnes, journal des corrections, requête et lignes de chaque réponse.
\end{itemize}

\subsection{Pipeline de profilage et d'analyse par table}

Le pipeline s'exécute dans Airflow sous la forme d'un DAG de neuf phases, chacune visible comme une tâche distincte (tableau~\ref{tab:phases}). Les phases se transmettent leur travail par de petits fichiers de manifeste, ce qui permet de relancer l'une d'elles isolément.

\begin{table}[H]
\centering
\small
\begin{tabular}{|C{0.8cm}|L{4.4cm}|L{8.8cm}|}
\hline
\thd{N\textsuperscript{o}} & \thd{Phase} & \thd{Traitement} \\ \hline
1 & Validation de la demande & Contrôle des identifiants et du dossier de la version reçus du serveur. \\ \hline
2 & Notification de démarrage & Contrôle de santé du serveur, puis signal de début d'exécution. \\ \hline
3 & Découverte des tables & Recherche des fichiers CSV de la version ; le fichier de contexte réservé n'est pas une table. \\ \hline
4 & Profilage & Lecture avec détection d'encodage et de séparateur, profil SweetViz ; une table vide est conservée et signalée. \\ \hline
5 & Métriques et analyse par le modèle & Mesures, rôles des colonnes, contrôles de format et de dates, signature ; premier résultat publié ; puis analyse du modèle et second résultat publié. \\ \hline
6 & Publication & Nouvelle tentative pour les résultats refusés par le serveur. \\ \hline
7 & Règles proposées & Règles dérivées de l'objectif de l'utilisateur, au mieux : cette phase n'échoue jamais l'exécution. \\ \hline
8 & Analyse globale & Charge utile compacte, un appel du modèle si activé, sinon analyse déterministe. \\ \hline
9 & Clôture & S'exécute toujours en dernier : statut \texttt{success}, \texttt{partial} ou \texttt{failed}, tables en échec, message lisible. \\ \hline
\end{tabular}
\caption{Phases du DAG de profilage et d'analyse}
\label{tab:phases}
\end{table}

La figure~\ref{fig:pipeline-ia} présente le détail de la phase 5, qui concentre l'intelligence du pipeline.

\figplaceholder[1]{pipeline_ia.png}{Architecture du traitement d'une table (phase 5)}{fig:pipeline-ia}{Mesures déterministes (métriques, rôles par indices, formats, dates, relations, signature) ; publication du premier résultat ; colonnes de rôle incertain soumises au modèle ; analyse du modèle : génération, validation contre les preuves, réparation, repli déterministe ; publication du résultat enrichi.}

\subsubsection{Rôle des colonnes}

Chaque colonne reçoit un rôle parmi seize (identifiant, référence, courriel, téléphone, URL, date, montant, quantité, catégorie, statut, booléen, nom, lieu, code postal, texte libre, inconnu). Le rôle est noté à partir des mots du nom de la colonne, des statistiques de ses valeurs et des formats qu'elles suivent ; les indices sont conservés pour pouvoir l'expliquer. Seules les colonnes qui restent incertaines sont soumises au modèle, avec l'objectif de l'utilisateur et ce que représente une ligne ; la source de la décision est enregistrée.

\subsubsection{Contrôles de format et de dates}

Pour chaque colonne de texte, le worker compte les valeurs qui sont un courriel, un téléphone, une URL, une date ou un nombre, et garde des exemples de valeurs invalides. Pour chaque couple de colonnes de date, il mesure, sur les lignes où les deux sont remplies, à quelle fréquence la première est postérieure à la seconde. Le serveur juge ensuite ces mesures contre les règles approuvées, sans relire les fichiers.

\subsubsection{Relations entre tables}

Une relation est proposée lorsque au moins 85\,\% des valeurs d'une colonne source existent dans une colonne cible presque unique. Une colonne est source si elle est majoritairement remplie et n'est pas l'identifiant de sa propre table. Les correspondances ambiguës, sans indice dans le nom, sont écartées. Chaque relation porte sa confiance, sa couverture et les effectifs.

\subsubsection{Analyse contrôlée par le modèle}

L'analyse d'une table suit un flux contrôlé \textit{générer, valider, réparer, replier}, implémenté avec LangGraph : la réponse doit respecter un schéma, chaque pourcentage cité est comparé à la colonne qu'il concerne, une réponse invalide donne lieu à une tentative de réparation, et, si elle reste non fiable, une analyse déterministe est utilisée. Le statut du modèle (\texttt{ok}, \texttt{unavailable}, \texttt{invalid}, \texttt{disabled}, \texttt{pending}) est publié avec l'analyse.

\subsection{Score et dimensions de qualité}

Le serveur calcule le score d'une table comme 100 moins le pourcentage moyen de valeurs manquantes moins le pourcentage de lignes dupliquées ; le score d'une version est la moyenne pondérée par le nombre de lignes. Les notes vont de A (à partir de 90) à D (moins de 60). Quatre dimensions complètent ce score : complétude, unicité, validité et cohérence. Chacune est affichée avec son numérateur et son dénominateur, ou « non mesurée » avec la raison.

\subsection{Règles de qualité}

Une règle cible une table et, le cas échéant, une colonne. Les types disponibles sont : colonne sans valeur manquante, au plus X\,\% de valeurs manquantes, valeurs uniques, au moins N lignes, au plus X\,\% de doublons, score minimal, baisse maximale du nombre de lignes par rapport à la version précédente, valeurs au bon format (courriel, téléphone, URL, date, nombre) et ordre de deux dates. Chaque règle est d'état \texttt{ACTIVE}, \texttt{SUGGESTED} ou \texttt{DISMISSED}, de gravité \texttt{WARNING} ou \texttt{CRITICAL}, et son résultat est \texttt{PASS}, \texttt{FAIL} ou \texttt{SKIPPED}. Les règles proposées par le pipeline restent des suggestions tant qu'elles ne sont pas approuvées, et une suggestion écartée ne revient pas. Le code ne retient parmi les propositions du modèle que celles qui sont vérifiables sur les données : tables et colonnes réelles, types connus, seuils raisonnables.

\subsection{Corrections}

La correction d'une table obéit à deux règles : l'original n'est jamais modifié, et le modèle ne peut qu'\textbf{assembler des opérations connues}. Les corrections viennent de deux sources : ce que le système trouve (espaces superflus, marqueurs d'absence, dates en plusieurs formats, valeur écrite de plusieurs façons, doublons) ou une demande écrite en langage naturel, que le modèle traduit en opérations choisies parmi douze, chacune vérifiée contre les colonnes. L'aperçu donne le nombre de valeurs touchées, des exemples et des avertissements ; une date ambiguë n'est jamais devinée. L'application crée une nouvelle version, analysée comme un dépôt, avec un journal. Les fichiers corrigés sont écrits en UTF-8, séparés par des virgules, en-tête inchangé.

\subsection{Graphe de connaissances}

Le modèle de graphe est déduit des rôles et des données : une ligne d'une table de comptes devient un nœud \textit{Compte}, et le \textit{Client} qu'elle référence devient un nœud à part entière, dont les propriétés sont les colonnes constantes pour un client. Les relations proviennent des liens trouvés par l'analyse. Quatre cas courants sont traités sans identifiant déclaré :
\begin{itemize}
  \item une table sans identifiant dont les lignes se distinguent par une \textbf{combinaison de références} (une commande et un numéro de ligne) est identifiée par cette combinaison ;
  \item une colonne qui contient les \textbf{numéros de ligne} d'un autre fichier sans identifiant le relie à ces lignes ;
  \item une \textbf{table de deux références} (amis, produit-magasin) devient une relation, sans nœud propre ;
  \item une colonne d'identifiants de la \textbf{même table} (un identifiant de responsable) relie deux lignes de cette table.
\end{itemize}
Lorsque les lignes d'une même clé se contredisent (deux courriels pour un client), le graphe conserve \textbf{toutes} les valeurs, marque le nœud comme conflictuel et n'en choisit aucune. Les références vers un nœud absent, les tables liées à rien et les liens incertains sont signalés avant la construction. Chaque nœud porte une propriété de périmètre (\texttt{d<jeu>v<version>}) : les versions ne se mélangent pas.

\subsection{Service de questions et garde-fous}\label{sec:questions-conception}

Le service de questions commence par \textbf{trier} la question en trois types (tableau~\ref{tab:types-questions}). Seul le premier utilise Neo4j ; les deux autres répondent même avant la construction du graphe. Un tri incertain retombe sur le type « qualité ».

\begin{table}[H]
\centering
\small
\begin{tabular}{|L{1.8cm}|L{4.6cm}|L{5cm}|C{1.6cm}|}
\hline
\thd{Type} & \thd{Exemples} & \thd{Source de la réponse} & \thd{Neo4j} \\ \hline
Données (\textit{graph}) & « Combien de commandes par client ? » & Requête Cypher écrite par le modèle, contrôlée, puis exécutée sur le graphe de la version & oui \\ \hline
Qualité (\textit{quality}) & « Quelles règles échouent et pourquoi ? » & Faits de qualité joints par le serveur : règles et messages, dimensions, constats, rôles de colonnes & non \\ \hline
Aide (\textit{help}) & « Comment appliquer une correction ? » & Guide d'utilisation de l'application & non \\ \hline
\end{tabular}
\caption{Types de questions et sources des réponses}
\label{tab:types-questions}
\end{table}

Comme le modèle écrit la requête, rien de ce qu'il écrit n'est considéré comme sûr. Une requête ne s'exécute que si les conditions suivantes sont remplies :
\begin{itemize}
  \item c'est \textbf{une seule instruction}, en \textbf{lecture seule} (les clauses d'écriture, d'appel de procédure, de chargement ou d'administration sont refusées, hors chaînes et commentaires) ;
  \item elle contient \texttt{MATCH} et \texttt{RETURN} ;
  \item \textbf{chaque nœud} est filtré sur le périmètre de la version interrogée, et aucun autre périmètre n'apparaît ;
  \item Neo4j confirme par \texttt{EXPLAIN} qu'elle ne fait que lire, puis elle s'exécute avec une limite de durée (15\,s) ;
  \item au plus 50 lignes sont lues et 15 renvoyées à l'interface ; le modèle dispose de \textbf{deux tentatives} et reçoit la raison d'un refus pour se corriger.
\end{itemize}
La réponse est rédigée \textbf{uniquement} à partir des lignes obtenues, des faits de qualité et du guide, avec l'instruction de ne rien inventer. La requête exécutée et ses lignes sont renvoyées avec la réponse, pour qu'on puisse la vérifier. Si aucune requête admissible n'est trouvée, la réponse s'appuie sur les faits de qualité seuls et le dit. Le diagramme de séquence correspondant est présenté au chapitre~7.

\subsection{Comparaison avec les approches existantes}

Le tableau~\ref{tab:comparaison-approches} situe la conception retenue par rapport à une approche où le modèle de langage interprète seul un rapport de profilage.

\begin{table}[H]
\centering
\small
\begin{tabular}{|L{3.4cm}|L{5.4cm}|L{5.2cm}|}
\hline
\thd{Axe} & \thd{Modèle interprétant un rapport} & \thd{Conception retenue} \\ \hline
Rôle du modèle & Lire le profil et tout en déduire & Trancher les cas ouverts, rédiger et traduire en opérations connues \\ \hline
Rôle des colonnes & Déduit par le modèle & Calculé par indices explicites, modèle en dernier recours \\ \hline
Contrôle des sorties & Validation du format uniquement & Comparaison de chaque chiffre à sa colonne, repli déterministe \\ \hline
Surveillance & Aucune & Règles évaluées à chaque exécution, alertes \\ \hline
Correction & Hors périmètre & Sur une copie, avec aperçu chiffré et journal \\ \hline
Vision transverse & Un appel sur une vue agrégée & Relations mesurées, graphe, questions contrôlées \\ \hline
\end{tabular}
\caption{Comparaison de la conception retenue avec une interprétation directe par le modèle}
\label{tab:comparaison-approches}
\end{table}

\section*{Conclusion}
\addcontentsline{toc}{section}{Conclusion}
Ce chapitre a présenté la conception de \projet{} : une conception globale fondée sur le formalisme UML ; des architectures logicielle, logique et physique organisées en cinq blocs ; une sécurité fondée sur des mesures simples ; et des composants d'intelligence artificielle gouvernés par un même principe, selon lequel le code mesure et contrôle tandis que le modèle interprète et rédige. Les chapitres~5 à~7 décrivent la réalisation, sprint par sprint, et le chapitre~\ref{chap:benchmark} l'évaluation de la solution.

\chapter{Réalisation : Sprints 1 et 2}

\section*{Introduction}
\addcontentsline{toc}{section}{Introduction}
Ce chapitre et les deux suivants présentent la réalisation de \projet{}, sprint par sprint, à partir de la conception décrite au chapitre~4. Ce chapitre couvre les deux premiers sprints. Après une présentation de l'environnement de travail, nous décrivons le sprint~1, qui a rendu opérationnels la lecture et le profilage statistique des jeux de données, puis le sprint~2, qui a doté la plateforme de son application web : authentification, dépôt versionné, suivi en temps réel et reprise après incident. Pour chaque sprint, nous présentons l'objectif et le backlog, les éléments réalisés, les diagrammes de séquence et d'activité des fonctionnalités clés, puis les résultats obtenus lors de la revue de sprint.

%==============================================================================
\section{Environnement de travail}
%==============================================================================

Avant de décrire les sprints, nous présentons l'environnement matériel et logiciel dans lequel la solution a été développée et évaluée.

\subsection{Environnement matériel}

Le tableau~\ref{tab:materiel} présente la configuration matérielle utilisée. L'exécution locale des modèles de langage en est la composante la plus exigeante : la mémoire graphique et la mémoire vive conditionnent la taille des modèles utilisables et la durée de génération.

\begin{table}[H]
\centering
\begin{tabular}{|L{4cm}|L{10cm}|}
\hline
\thd{Composant} & \thd{Caractéristiques} \\ \hline
Processeur & Intel Core Ultra 9 275HX \\ \hline
Mémoire vive & 16 Go \\ \hline
Carte graphique & NVIDIA GeForce RTX 5070 Ti \\ \hline
Stockage & SSD de 1 To \\ \hline
Système d'exploitation & Windows 11 \\ \hline
\end{tabular}
\caption{Environnement matériel}
\label{tab:materiel}
\end{table}

Sur ce matériel s'appuie un ensemble d'outils logiciels, présentés ci-après.

\subsection{Environnement logiciel}

Le tableau~\ref{tab:logiciel} présente les outils et technologies utilisés tout au long du projet.

\begin{table}[H]
\centering
\begin{tabular}{|C{2cm}|L{3.6cm}|L{8.4cm}|}
\hline
\thd{Logo} & \thd{Outil} & \thd{Usage} \\ \hline
\toolicon{intellij.png} & IntelliJ IDEA & Développement du serveur applicatif Java \\ \hline
\toolicon{pycharm.png} & PyCharm & Développement du worker et du service de questions \\ \hline
\toolicon{vscode.png} & Visual Studio Code & Développement de l'interface \\ \hline
\toolicon{spring.png} & Spring Boot & Serveur applicatif, sécurité, persistance \\ \hline
\toolicon{angular.png} & Angular & Interface web \\ \hline
\toolicon{python.png} & Python & Worker de traitement et service de questions \\ \hline
\toolicon{ollama.png} & Ollama & Exécution locale des modèles de langage \\ \hline
\toolicon{airflow.png} & Apache Airflow & Orchestration du pipeline \\ \hline
\toolicon{neo4j.png} & Neo4j & Base de graphe \\ \hline
\toolicon{docker.png} & Docker & Conteneurisation des services de la plateforme \\ \hline
\toolicon{postgresql.png} & PostgreSQL & Base de métadonnées d'Airflow \\ \hline
\toolicon{pandas.png} & pandas & Manipulation des données tabulaires \\ \hline
-- & SweetViz & Profilage statistique \\ \hline
\toolicon{powershell.png} & PowerShell & Scripts de démarrage sous Windows \\ \hline
\toolicon{git.png} & Git et GitHub Actions & Gestion de versions et intégration continue \\ \hline
\end{tabular}
\caption{Environnement logiciel}
\label{tab:logiciel}
\end{table}

L'environnement étant en place, nous pouvons aborder le premier sprint.

%==============================================================================
\section{Sprint 1 — Ingestion et profilage}
%==============================================================================

Le premier sprint réalise l'epic E1. Nous présentons son objectif et son backlog, puis la lecture tolérante, les mesures par table, les relations entre tables et la signature de contenu, avant d'en modéliser le comportement et d'en présenter les résultats.

\subsection{Sprint Goal et Sprint Backlog}

Le \textbf{Sprint Goal} du sprint~1 est le suivant : \textit{« mesurer automatiquement et de manière fiable la qualité de tout jeu de données »}. Le tableau~\ref{tab:backlog-s1} détaille les récits retenus et leurs tâches.

\begin{table}[H]
\centering
\small
\begin{tabular}{|C{1.1cm}|L{4.6cm}|C{1.2cm}|L{6.4cm}|}
\hline
\thd{US} & \thd{Récit utilisateur} & \thd{ID} & \thd{Tâche} \\ \hline
\multirow{2}{*}{US1} & \multirow{2}{4.6cm}{Lecture de tous les fichiers} & T1.1 & Découverte des fichiers CSV d'une version \\ \cline{3-4}
 & & T1.2 & Lecture tolérante : encodage, séparateur, lignes mal formées \\ \hline
\multirow{2}{*}{US2} & \multirow{2}{4.6cm}{Profil statistique de chaque table} & T1.3 & Profil SweetViz de chaque table \\ \cline{3-4}
 & & T1.4 & Métriques par table et par colonne \\ \hline
US3 & Liens entre tables & T1.5 & Recouvrement de valeurs entre colonnes de tables différentes \\ \hline
US4 & Détection des changements & T1.6 & Signatures de contenu et de schéma \\ \hline
\end{tabular}
\caption{Sprint Backlog du sprint 1}
\label{tab:backlog-s1}
\end{table}

Ces tâches s'enchaînent dans les premières phases du pipeline décrit au chapitre~4. La première difficulté concerne la lecture des fichiers.

\subsection{Lecture tolérante des sources}

Les exports issus des systèmes clients sont hétérogènes : encodages différents, séparateurs différents selon les conventions régionales, lignes mal formées. Une lecture naïve échoue ou, pire, produit une table d'une seule colonne. Le worker essaie donc successivement plusieurs combinaisons d'encodage et de séparateur et retient la lecture qui produit le \textbf{plus grand nombre de colonnes}, critère simple et efficace pour identifier le bon séparateur. Si la lecture stricte échoue, un analyseur indulgent ignore les lignes mal formées et consigne un avertissement dans les métriques : le nombre de lignes d'une table peut alors être inférieur au nombre de lignes du fichier. Les chaînes \texttt{NULL}, \texttt{N/A} ou \texttt{NaN} sont lues comme des valeurs manquantes, comme dans le reste de la plateforme. Un fichier réduit à son en-tête est conservé avec zéro ligne et signalé comme table en échec, sans interrompre le traitement des autres.

\subsection{Mesures par table et par colonne}

Pour chaque table, le worker calcule de façon déterministe le nombre de lignes et de colonnes, les valeurs manquantes, les valeurs distinctes, les lignes dupliquées et, pour les colonnes numériques, les quartiles, la moyenne et l'écart type. Le profil SweetViz fournit en complément les valeurs fréquentes et les associations entre colonnes ; l'application construit elle-même ses écrans à partir de ces résultats structurés et ne génère pas de pages HTML.

\subsection{Relations entre tables et signature de contenu}

Les relations entre tables sont déduites du recouvrement des valeurs (chapitre~4, section~III). Chaque table reçoit par ailleurs deux empreintes : \texttt{content\_hash}, calculée sur ses valeurs, et \texttt{schema\_hash}, calculée sur les noms et les types de ses colonnes. Le serveur les compare à celles de la version précédente pour qualifier chaque table de \textit{nouvelle}, \textit{modifiée}, \textit{inchangée} ou \textit{supprimée}.

\subsection{Modélisation dynamique du profilage}

Le diagramme de séquence de la figure~\ref{fig:seq-profilage} décrit les interactions entre composants lors du profilage d'une version.

\figplaceholder[1]{seq_profilage.png}{Diagramme de séquence — Profilage d'une version}{fig:seq-profilage}{Diagramme de séquence (généré).}

Le diagramme d'activité de la figure~\ref{fig:act-profilage} détaille, pour sa part, le traitement appliqué à chaque table.

\figplaceholder[0.6]{act_lecture_profilage.png}{Diagramme d'activité — Lecture tolérante et profilage d'une table}{fig:act-profilage}{Diagramme d'activité (généré).}

Ces traitements ont été validés lors de la revue de sprint, dont nous présentons les résultats.

\subsection{Résultats du sprint 1}

Pour ne pas exposer de données client dans ce rapport, la validation a été conduite sur des jeux de données de démonstration. Le profilage lit les fichiers quelle que soit leur forme, publie les métriques de chaque table et détecte les liens entre tables. La figure~\ref{fig:tables-profil} présente le profil d'une table tel que l'application le restitue.

\figplaceholder[0.95]{capture_tables.png}{Profil d'une table dans l'application}{fig:tables-profil}{Capture de l'onglet « Tables » : colonnes et leur rôle, valeurs fréquentes, valeurs manquantes, associations.}

\noindent\textbf{Sprint Review.} L'incrément livré permet de profiler automatiquement les fichiers d'une version. La rétrospective a mis en évidence la nécessité de présenter ces résultats dans une application utilisable par plusieurs personnes et de suivre l'exécution en direct, ce qui constitue l'objet du sprint~2.

%==============================================================================
\section{Sprint 2 — Application web et exécution suivie}
%==============================================================================

Le deuxième sprint réalise l'epic E2. Après son objectif et son backlog, nous présentons l'authentification, le dépôt versionné, le déclenchement et le suivi de l'exécution, puis la reprise après incident et le tableau de bord, avant de modéliser les interactions et de présenter les résultats.

\subsection{Sprint Goal et Sprint Backlog}

Le \textbf{Sprint Goal} du sprint~2 est le suivant : \textit{« déposer des fichiers, lancer l'analyse et la suivre en temps réel, de façon sécurisée »}. Le tableau~\ref{tab:backlog-s2} détaille les récits retenus et leurs tâches.

\begin{table}[H]
\centering
\small
\begin{tabular}{|C{1.1cm}|L{4.6cm}|C{1.2cm}|L{6.4cm}|}
\hline
\thd{US} & \thd{Récit utilisateur} & \thd{ID} & \thd{Tâche} \\ \hline
US5 & Authentification & T2.1 & Connexion, jeton JWT, gardes de routes \\ \hline
\multirow{3}{*}{US6} & \multirow{3}{4.6cm}{Dépôt versionné} & T2.2 & Validation complète avant écriture, empreinte SHA-256 \\ \cline{3-4}
 & & T2.3 & Création du jeu de données et de la version \\ \cline{3-4}
 & & T2.4 & Détection d'un dépôt identique \\ \hline
\multirow{3}{*}{US7} & \multirow{3}{4.6cm}{Suivi en temps réel} & T2.5 & Déclenchement d'Airflow après validation de la transaction \\ \cline{3-4}
 & & T2.6 & Rappels protégés du pipeline et statuts \\ \cline{3-4}
 & & T2.7 & Flux d'événements SSE et repli par interrogation \\ \hline
\multirow{2}{*}{US8} & \multirow{2}{4.6cm}{Reprise et relance} & T2.8 & Surveillance des exécutions bloquées et reprise au démarrage \\ \cline{3-4}
 & & T2.9 & Bouton de relance \\ \hline
US9 & Tableau de bord & T2.10 & Liste des jeux de données avec score, note et état \\ \hline
\end{tabular}
\caption{Sprint Backlog du sprint 2}
\label{tab:backlog-s2}
\end{table}

\subsection{Authentification et dépôt versionné}

Le serveur expose une route de connexion qui renvoie un jeton JWT signé, valable huit heures et conservé par l'interface dans le stockage de session. Toute autre route de l'API, hors connexion et rappels du pipeline, exige le rôle administrateur.

Le dépôt de fichiers obéit à une logique de transaction : tous les fichiers sont validés avant toute écriture, une empreinte SHA-256 est calculée pour chacun, la version est créée et les fichiers sont stockés, puis le serveur déclenche Airflow \textbf{après} la validation de la transaction, de sorte qu'Airflow ne démarre jamais sur une version qui n'existe pas encore. Un dépôt dont les noms de fichiers et les contenus sont identiques à une version existante n'en crée pas de nouvelle : l'interface propose d'ouvrir la version existante ou de la relancer.

\subsection{Déclenchement et suivi en temps réel}

Le serveur déclenche le DAG par l'API REST d'Airflow, avec pour identifiant d'exécution son propre identifiant d'exécution : une exécution dans Airflow et une version dans l'application se correspondent toujours. Le worker rend compte par cinq rappels protégés (démarrage, résultat de table, règles proposées, analyse globale, fin). Le résultat de chaque table est publié \textbf{deux fois} : d'abord les métriques, puis l'analyse enrichie par le modèle, de sorte que l'utilisateur voit une table dès que ses mesures sont prêtes. Le statut final est \texttt{COMPLETED}, \texttt{PARTIAL} (au moins une table en échec, avec la raison) ou \texttt{FAILED}.

Le serveur diffuse ces évolutions à l'interface par un flux SSE. Si le flux est interrompu, l'interface bascule automatiquement sur une interrogation toutes les cinq secondes, en gardant les données à l'écran jusqu'à l'arrivée des nouvelles pour que la page ne saute pas.

\subsection{Reprise après incident}

Trois mécanismes évitent qu'une exécution reste bloquée :
\begin{itemize}
  \item les rappels du worker sont \textbf{réessayés} pendant environ deux minutes (2, 4, \ldots, 64 secondes), de sorte qu'un redémarrage du serveur n'interrompt pas une exécution en cours ;
  \item au \textbf{démarrage}, le serveur interroge Airflow pour chaque exécution en attente ou en cours : une exécution inconnue d'Airflow est relancée si ses fichiers existent encore, une exécution échouée côté Airflow est marquée \texttt{FAILED}, une exécution en cours est laissée en l'état ;
  \item une \textbf{surveillance} vérifie chaque minute les exécutions en attente ou en cours depuis plus de dix minutes ; elle marque \texttt{FAILED}, avec un message clair, toute exécution inconnue ou échouée dans Airflow, ainsi que toute exécution non terminée après 180 minutes.
\end{itemize}

\subsection{Modélisation dynamique}

Le diagramme de séquence de la figure~\ref{fig:seq-depot} décrit l'enchaînement du dépôt à la restitution.

\figplaceholder[1]{seq_depot.png}{Diagramme de séquence — Dépôt, exécution et suivi}{fig:seq-depot}{Diagramme de séquence (généré).}

Le diagramme de la figure~\ref{fig:seq-reprise} décrit la reprise après un redémarrage du serveur.

\figplaceholder[0.9]{seq_reprise.png}{Diagramme de séquence — Reprise après redémarrage et surveillance}{fig:seq-reprise}{Diagramme de séquence (généré).}

\subsection{Résultats du sprint 2}

L'utilisateur dispose désormais d'une application complète pour déposer ses fichiers et suivre l'analyse. La figure~\ref{fig:capture-depot} présente la fenêtre de dépôt et la figure~\ref{fig:capture-dashboard} le tableau de bord des jeux de données.

\figplaceholder[0.7]{capture_depot.png}{Fenêtre « Nouvelle analyse »}{fig:capture-depot}{Capture de la fenêtre de dépôt : fichiers, nom du jeu de données, domaine, objectif, ce que représente une ligne.}

\figplaceholder[0.95]{capture_dashboard.png}{Tableau de bord de tous les jeux de données}{fig:capture-dashboard}{Capture de la page /datasets : chiffres clés, recherche, filtres, une carte par jeu de données avec anneau de score, note, état et règles en échec.}

La figure~\ref{fig:capture-overview} montre le suivi d'une exécution et la vue d'ensemble d'une version.

\figplaceholder[0.95]{capture_overview.png}{Vue d'ensemble d'une version}{fig:capture-overview}{Capture de l'onglet « Vue d'ensemble » : progression de l'exécution, score et note, dimensions, état de chaque table.}

\noindent\textbf{Sprint Review.} L'incrément livré permet de déposer des fichiers, de suivre l'analyse étape par étape et de relancer une exécution. La rétrospective a mis en évidence que le score affiché restait une valeur globale difficile à expliquer, et que le sens des colonnes dépendait du contexte, ce qui fait l'objet du sprint~3.

\section*{Conclusion}
\addcontentsline{toc}{section}{Conclusion}
Ce chapitre a présenté l'environnement de travail et la réalisation des deux premiers sprints : la lecture et le profilage des fichiers, puis l'application web avec son authentification, son dépôt versionné, son suivi en temps réel et sa reprise après incident. Le chapitre suivant présente les sprints~3 et~4 : le contexte métier, les rôles de colonnes et les dimensions de qualité, puis les règles de qualité et l'industrialisation.

\chapter{Réalisation : Sprints 3 et 4}

\section*{Introduction}
\addcontentsline{toc}{section}{Introduction}
Ce chapitre poursuit la présentation de la réalisation avec les sprints~3 et~4. Le sprint~3 rend les résultats compréhensibles : le contexte métier oriente l'analyse, chaque colonne reçoit un rôle expliqué, et le score est décomposé en dimensions dont le calcul est visible. Le sprint~4 permet de surveiller la qualité dans le temps grâce à des règles évaluées à chaque exécution, et rend l'ensemble déployable et testé. Pour chaque sprint, nous présentons l'objectif et le backlog, les éléments réalisés, les diagrammes de séquence et d'activité des fonctionnalités clés, puis les résultats obtenus lors de la revue de sprint.

%==============================================================================
\section{Sprint 3 — Contexte, rôles et dimensions}
%==============================================================================

Le troisième sprint réalise l'epic E3. Nous présentons son objectif et son backlog, puis le contexte métier, les rôles de colonnes, l'analyse contrôlée par le modèle, les dimensions de qualité et la comparaison de versions.

\subsection{Sprint Goal et Sprint Backlog}

Le \textbf{Sprint Goal} du sprint~3 est le suivant : \textit{« comprendre à quoi sert chaque colonne et expliquer le score de qualité »}. Le tableau~\ref{tab:backlog-s3} détaille les récits retenus et leurs tâches.

\begin{table}[H]
\centering
\small
\begin{tabular}{|C{1.1cm}|L{4.6cm}|C{1.2cm}|L{6.4cm}|}
\hline
\thd{US} & \thd{Récit utilisateur} & \thd{ID} & \thd{Tâche} \\ \hline
\multirow{2}{*}{US10} & \multirow{2}{4.6cm}{Contexte métier} & T3.1 & Saisie du domaine, de l'objectif et du grain d'une ligne \\ \cline{3-4}
 & & T3.2 & Contexte libre global ou par table \\ \hline
\multirow{2}{*}{US11} & \multirow{2}{4.6cm}{Rôle des colonnes} & T3.3 & Notation du rôle par indices (nom, valeurs, formats) \\ \cline{3-4}
 & & T3.4 & Arbitrage des colonnes incertaines par le modèle \\ \hline
\multirow{2}{*}{US12} & \multirow{2}{4.6cm}{Analyse contrôlée} & T3.5 & Flux générer, valider, réparer, replier \\ \cline{3-4}
 & & T3.6 & Statut du modèle et analyse globale \\ \hline
US13 & Score et dimensions & T3.7 & Calcul du score, des notes et de quatre dimensions expliquées \\ \hline
US14 & Comparaison de versions & T3.8 & Différences entre deux versions \\ \hline
\end{tabular}
\caption{Sprint Backlog du sprint 3}
\label{tab:backlog-s3}
\end{table}

\subsection{Contexte métier}

À l'envoi, l'utilisateur peut indiquer le \textbf{domaine}, l'\textbf{objectif} de l'analyse et ce que \textbf{représente une ligne} (un client, un compte, une transaction). Ces trois informations sont enregistrées dans un fichier de contexte à côté des fichiers CSV, conservées lors d'une relance et lues par le pipeline pour guider les rôles, l'analyse du modèle et les règles proposées. Elles évitent une erreur typique : un identifiant répété n'est un problème que s'il contredit ce qu'une ligne représente. Un contexte libre, global ou propre à une table, peut aussi être fourni par un fichier réservé dont le nom n'est jamais traité comme une table.

\subsection{Rôle des colonnes}

Le rôle de chaque colonne est noté à partir des mots de son nom, des statistiques de ses valeurs et des formats qu'elles suivent ; les indices qui ont fait pencher la décision sont conservés. Le tableau~\ref{tab:roles-exemples} donne quelques exemples du raisonnement.

\begin{table}[H]
\centering
\small
\begin{tabular}{|L{3.2cm}|L{4.8cm}|L{3.2cm}|C{2.4cm}|}
\hline
\thd{Colonne} & \thd{Indices} & \thd{Rôle} & \thd{Décidé par} \\ \hline
\texttt{customer\_id} & Mot \textit{id}, valeurs presque toutes distinctes & Identifiant & Indices \\ \hline
\texttt{email} & Mot \textit{email}, valeurs au format courriel & Courriel & Indices \\ \hline
\texttt{opened\_on} & Valeurs lisibles comme des dates & Date & Indices \\ \hline
Colonne au nom opaque & Aucun indice net & Rôle proposé par le modèle & Modèle \\ \hline
\end{tabular}
\caption{Exemples d'attribution du rôle d'une colonne}
\label{tab:roles-exemples}
\end{table}

Le diagramme d'activité de la figure~\ref{fig:act-roles} résume la décision.

\figplaceholder[0.6]{act_roles.png}{Diagramme d'activité — Attribution du rôle d'une colonne}{fig:act-roles}{Diagramme d'activité (généré).}

\subsection{Analyse contrôlée par le modèle}

Pour chaque table, le modèle produit une analyse structurée : constats, causes probables, impacts, recommandations. Le flux \textit{générer, valider, réparer, replier} compare chaque pourcentage cité à la colonne qu'il concerne : une valeur attribuée à une autre colonne est rejetée, tandis que des formes raccourcies de la réponse sont tolérées. Si l'analyse reste non fiable, une analyse déterministe est publiée. L'interface indique l'origine du résultat, par exemple « Déterministe seulement – Ollama indisponible ». Le diagramme de séquence de la figure~\ref{fig:seq-analyse} décrit les échanges.

\figplaceholder[1]{seq_analyse_table.png}{Diagramme de séquence — Analyse d'une table par le modèle}{fig:seq-analyse}{Diagramme de séquence (généré).}

L'analyse globale suit le même principe : une charge compacte est construite à partir des résultats de table, un seul appel du modèle est effectué si l'option est activée, et une analyse déterministe sert de repli.

\subsection{Score et dimensions de qualité}

L'onglet « Vue d'ensemble » affiche le score de la version et sa note, puis les quatre dimensions de qualité avec la manière dont chacune est calculée. Une dimension qui ne peut pas être mesurée l'indique et en donne la raison, plutôt que d'afficher une valeur trompeuse. Le tableau~\ref{tab:dimensions-affichees} résume ces dimensions.

\begin{table}[H]
\centering
\small
\begin{tabular}{|L{2.8cm}|L{6.4cm}|L{4.8cm}|}
\hline
\thd{Dimension} & \thd{Numérateur / dénominateur} & \thd{Exemple de non-mesure} \\ \hline
Complétude & Cellules renseignées / cellules (lignes $\times$ colonnes) & -- \\ \hline
Unicité & Lignes qui ne sont pas la copie exacte d'une autre / lignes (un identifiant répété n'est pas un doublon) & -- \\ \hline
Validité & Valeurs au bon format / valeurs renseignées des colonnes courriel, téléphone, URL ou date & Aucune colonne reconnue comme courriel, téléphone, URL ou date \\ \hline
Cohérence & Lignes où les règles d'ordre de dates approuvées sont respectées / lignes où les deux dates sont remplies & Aucune règle d'ordre de dates approuvée \\ \hline
\end{tabular}
\caption{Dimensions de qualité affichées}
\label{tab:dimensions-affichees}
\end{table}

\figplaceholder[0.95]{capture_dimensions.png}{Score, note et dimensions de qualité}{fig:capture-dimensions}{Capture de l'onglet « Vue d'ensemble » : anneau de score, note, quatre dimensions avec leur calcul, état de chaque table, colonnes à revoir.}

\subsection{Comparaison de versions}

L'onglet « Comparer » montre les différences entre une version et une version plus ancienne : variation du score, nombre de lignes, colonnes ajoutées ou supprimées, changements de type et valeurs manquantes. La comparaison s'appuie sur les signatures de contenu et de schéma calculées au sprint~1.

\subsection{Résultats du sprint 3}

L'utilisateur dispose désormais d'un score expliqué, de colonnes dont le rôle est justifié et d'une analyse en langage naturel dont l'origine est affichée. La figure~\ref{fig:capture-insights} présente l'analyse sémantique et la figure~\ref{fig:capture-model} le modèle de données déduit.

\figplaceholder[0.95]{capture_insights.png}{Analyse sémantique d'une table}{fig:capture-insights}{Capture de l'onglet « Analyses » : groupes dépliables (rôles de colonnes, contrat de données, surveillance, causes, impacts) et recommandations globales.}

\figplaceholder[0.95]{capture_model.png}{Modèle de données déduit des relations}{fig:capture-model}{Capture de l'onglet « Modèle de données » : tables et relations, avec couverture et cardinalité.}

\noindent\textbf{Sprint Review.} L'incrément livré rend les résultats explicables. La rétrospective a identifié le besoin de surveiller la qualité dans le temps plutôt que de la constater une fois, ce qui fait l'objet du sprint~4.

%==============================================================================
\section{Sprint 4 — Règles de qualité et industrialisation}
%==============================================================================

Le quatrième sprint réalise l'epic E4. Après son objectif et son backlog, nous présentons les règles, leur évaluation et leurs alertes, le rapport imprimable, puis le déploiement et les tests automatisés.

\subsection{Sprint Goal et Sprint Backlog}

Le \textbf{Sprint Goal} du sprint~4 est le suivant : \textit{« surveiller la qualité par des règles et livrer une plateforme déployable et testée »}. Le tableau~\ref{tab:backlog-s4} détaille les récits retenus et leurs tâches.

\begin{table}[H]
\centering
\small
\begin{tabular}{|C{1.1cm}|L{4.6cm}|C{1.2cm}|L{6.4cm}|}
\hline
\thd{US} & \thd{Récit utilisateur} & \thd{ID} & \thd{Tâche} \\ \hline
US15 & Créer des règles & T4.1 & Types de règles, gravité, pause, édition \\ \hline
\multirow{2}{*}{US16} & \multirow{2}{4.6cm}{Règles proposées} & T4.2 & Suggestion à partir de l'objectif, filtrée par le code \\ \cline{3-4}
 & & T4.3 & Approbation et rejet définitif \\ \hline
\multirow{2}{*}{US17} & \multirow{2}{4.6cm}{Évaluation et alertes} & T4.4 & Évaluation en fin d'exécution et à la demande \\ \cline{3-4}
 & & T4.5 & Alerte : notification, badge, bandeau \\ \hline
US18 & Rapport imprimable & T4.6 & Page imprimable et export Markdown \\ \hline
\multirow{3}{*}{US19} & \multirow{3}{4.6cm}{Déploiement et tests} & T4.7 & Dockerfiles et Compose à la racine \\ \cline{3-4}
 & & T4.8 & Script de démarrage unique pour le développement \\ \cline{3-4}
 & & T4.9 & Intégration continue et suites de tests \\ \hline
\end{tabular}
\caption{Sprint Backlog du sprint 4}
\label{tab:backlog-s4}
\end{table}

\subsection{Règles de qualité}

Une règle est définie par son type, sa table et sa colonne, son seuil et sa gravité (\texttt{WARNING} ou \texttt{CRITICAL}) ; elle peut être suspendue. Neuf types sont disponibles (chapitre~4). Les deux plus récents portent sur le \textbf{format} des valeurs (courriel, téléphone, URL, date, nombre) et sur l'\textbf{ordre de deux dates}, et s'appuient sur les mesures du sprint~1 : le serveur les juge sans relire les fichiers.

Les règles sont évaluées à la fin de chaque exécution et à la demande (bouton « Vérifier maintenant »). Chaque résultat est \texttt{PASS}, \texttt{FAIL} ou \texttt{SKIPPED} (non évaluée, avec la raison : table absente de la version ou aucune version précédente à comparer, par exemple ; une colonne absente donne en revanche un échec), accompagné d'un message explicatif. Le diagramme de séquence de la figure~\ref{fig:seq-regles} décrit l'enchaînement.

\figplaceholder[1]{seq_regles.png}{Diagramme de séquence — Proposition, approbation et évaluation des règles}{fig:seq-regles}{Diagramme de séquence (généré).}

\subsection{Règles proposées à partir de l'objectif}

Le pipeline demande au modèle de proposer des règles pour l'objectif indiqué : « contacter les clients » donne par exemple un courriel renseigné et bien formé. Le code ne conserve que les règles réellement vérifiables sur les données. Sans objectif, ou sans modèle, un petit ensemble de règles est dérivé des mesures. Ces règles restent des \textbf{suggestions} (état \texttt{SUGGESTED}) tant que l'utilisateur ne les a pas approuvées, et une suggestion écartée ne revient pas.

\subsection{Alertes}

Un échec de règle déclenche une alerte visible à trois niveaux : une notification en direct, où que se trouve l'utilisateur, un badge dans la liste des jeux de données et un bandeau dans la vue d'ensemble. Les échecs de gravité critique sont comptés à part.

\figplaceholder[0.95]{capture_regles.png}{Onglet des règles de qualité}{fig:capture-regles}{Capture de l'onglet « Règles » : règles et résultats de la version, règles proposées à approuver, création d'une règle.}

\subsection{Rapport imprimable}

Le bouton « Rapport » ouvre une page imprimable de la version (impression ou enregistrement en PDF, ou export Markdown) qui reprend le score, les dimensions avec la définition de chaque score, les tables et les règles.

\subsection{Déploiement et tests automatisés}

Le dossier racine contient un fichier Docker Compose qui inclut ceux du worker et de Neo4j : tous les services partagent un réseau et les mêmes secrets, et le serveur joint Airflow, Neo4j et le service de questions par le nom de leur service. Une seule commande construit et démarre l'ensemble. Pour le développement sous Windows, un script ouvre une fenêtre par service (Ollama, Airflow, questions, serveur, interface) et lance Neo4j en arrière-plan ; le navigateur s'ouvre lorsque tout répond.

Quatre suites de tests couvrent les blocs, et un workflow d'intégration continue les exécute à chaque modification, avec la validation du fichier Compose (tableau~\ref{tab:tests}).

\begin{table}[H]
\centering
\small
\begin{tabular}{|L{3.6cm}|L{6.2cm}|C{1.6cm}|L{2.6cm}|}
\hline
\thd{Bloc} & \thd{Périmètre couvert} & \thd{Tests} & \thd{Outil} \\ \hline
Serveur & Authentification, dépôt, statuts, règles, dimensions, comparaison, corrections, graphe, reprise & 110 & Maven \\ \hline
Interface & Connexion et appels API, événements, formatage, règles, modèle de données, gardes de routes, page d'accueil, tableau de bord, animations & 76 & Vitest \\ \hline
Worker & Client du serveur, rôles, formats, règles proposées, contexte, formes de réponse du modèle, phases de bout en bout avec un faux serveur & 67 & pytest \\ \hline
Service de questions & Garde-fous, tri, réponses, service HTTP & 39 & pytest \\ \hline
\end{tabular}
\caption{Suites de tests automatisés}
\label{tab:tests}
\end{table}

Ces tests n'ont besoin ni d'Airflow, ni de Neo4j, ni d'Ollama : ces services sont simulés.

\subsection{Résultats du sprint 4}

La plateforme surveille désormais la qualité par des règles, alerte l'utilisateur et se déploie par une commande. La figure~\ref{fig:capture-airflow} présente la vue de l'orchestrateur pour une exécution réussie.

\figplaceholder[0.95]{capture_airflow.png}{Interface d'Airflow — vue graphe d'une exécution réussie}{fig:capture-airflow}{Capture de l'interface d'Airflow : les neuf tâches du DAG en succès.}

\noindent\textbf{Sprint Review.} L'incrément livré rend la plateforme exploitable en continu. La rétrospective a identifié le besoin de passer de la détection à la correction, ce qui fait l'objet du sprint~5.

\section*{Conclusion}
\addcontentsline{toc}{section}{Conclusion}
Ce chapitre a présenté la réalisation des sprints~3 et~4 : le contexte métier, les rôles de colonnes, l'analyse contrôlée et les dimensions de qualité, puis les règles, les alertes, le rapport, le déploiement et les tests. Le chapitre suivant présente les sprints~5 à~7 : les corrections, le graphe de connaissances et les questions en langage naturel.

\chapter{Réalisation : Sprints 5 à 7}

\section*{Introduction}
\addcontentsline{toc}{section}{Introduction}
Ce chapitre conclut la présentation de la réalisation avec les trois derniers sprints. Le sprint~5 permet de corriger les données sans jamais toucher à l'original. Le sprint~6 représente les tables sous la forme d'un graphe de connaissances dans Neo4j. Le sprint~7 permet d'interroger l'ensemble en langage naturel, et soigne l'expérience d'utilisation de l'interface. Pour chaque sprint, nous présentons l'objectif et le backlog, les éléments réalisés, les diagrammes de séquence des fonctionnalités clés, puis les résultats de la revue de sprint.

%==============================================================================
\section{Sprint 5 — Corrections}
%==============================================================================

Le cinquième sprint réalise l'epic E5. Nous présentons son objectif et son backlog, puis les opérations de correction, leurs deux sources, l'aperçu et l'application sur une copie.

\subsection{Sprint Goal et Sprint Backlog}

Le \textbf{Sprint Goal} du sprint~5 est le suivant : \textit{« corriger les données sans jamais modifier l'original »}. Le tableau~\ref{tab:backlog-s5} détaille les récits retenus et leurs tâches.

\begin{table}[H]
\centering
\small
\begin{tabular}{|C{1.1cm}|L{4.6cm}|C{1.2cm}|L{6.4cm}|}
\hline
\thd{US} & \thd{Récit utilisateur} & \thd{ID} & \thd{Tâche} \\ \hline
\multirow{2}{*}{US20} & \multirow{2}{4.6cm}{Corrections trouvées} & T5.1 & Repérage des défauts courants d'une table \\ \cline{3-4}
 & & T5.2 & Moteur d'exécution des opérations \\ \hline
\multirow{2}{*}{US21} & \multirow{2}{4.6cm}{Correction en langage naturel} & T5.3 & Traduction d'une demande en opérations par le modèle \\ \cline{3-4}
 & & T5.4 & Vérification des opérations contre les colonnes \\ \hline
\multirow{3}{*}{US22} & \multirow{3}{4.6cm}{Aperçu et application} & T5.5 & Aperçu chiffré avec exemples et avertissements \\ \cline{3-4}
 & & T5.6 & Création d'une version dérivée et journal \\ \cline{3-4}
 & & T5.7 & Interface de l'onglet « Corrections » \\ \hline
\end{tabular}
\caption{Sprint Backlog du sprint 5}
\label{tab:backlog-s5}
\end{table}

\subsection{Opérations de correction}

Toute correction se ramène à l'une des douze opérations du tableau~\ref{tab:operations}. Ce choix est délibéré : le modèle ne peut pas écrire de code ni inventer une transformation, il ne peut qu'assembler des opérations dont le comportement est connu et testé.

\begin{table}[H]
\centering
\small
\begin{tabular}{|L{3.8cm}|L{10.2cm}|}
\hline
\thd{Opération} & \thd{Effet} \\ \hline
\texttt{TRIM} & Supprimer les espaces autour des valeurs. \\ \hline
\texttt{COLLAPSE\_SPACES} & Réduire les espaces répétés à un seul. \\ \hline
\texttt{CASE} & Passer en majuscules, en minuscules ou en casse de titre. \\ \hline
\texttt{REPLACE\_VALUES} & Remplacer des valeurs entières (fautes de frappe, un même pays écrit de plusieurs façons). \\ \hline
\texttt{REPLACE\_TEXT} & Remplacer une partie d'une valeur (une unité, un préfixe). \\ \hline
\texttt{EMPTY\_TO\_NULL} & Transformer des marqueurs d'absence (\texttt{N/A}, \texttt{-}\ldots) en cellules vides. \\ \hline
\texttt{FILL\_EMPTY} & Remplir les cellules vides avec une valeur. \\ \hline
\texttt{NORMALIZE\_DATE} & Écrire les dates au format \texttt{aaaa-mm-jj} ; en mode automatique, une date ambiguë est laissée à la revue. \\ \hline
\texttt{DIGITS\_ONLY} & Ne garder que les chiffres (et un \texttt{+} initial). \\ \hline
\texttt{PAD\_LEFT} & Rétablir des zéros initiaux perdus dans un code. \\ \hline
\texttt{DROP\_EMPTY\_ROWS} & Supprimer les lignes où une colonne est vide. \\ \hline
\texttt{DROP\_DUPLICATE\_ROWS} & Supprimer les lignes copies exactes d'une ligne précédente. \\ \hline
\end{tabular}
\caption{Opérations de correction}
\label{tab:operations}
\end{table}

\subsection{Deux sources de corrections}

Les corrections proposées viennent de deux sources. La première est le \textbf{repérage automatique} : l'application inspecte la table et propose les opérations qui correspondent à ce qu'elle y trouve (espaces, marqueurs d'absence, dates en plusieurs formats, valeur écrite de plusieurs façons, doublons). La seconde est une \textbf{demande écrite en langage naturel}. Le modèle reçoit la demande, la table et ses colonnes, avec des règles strictes : n'utiliser que des colonnes qui existent, ne jamais inventer de valeur (un remplacement doit provenir de la demande ou des valeurs les plus fréquentes), ne pas deviner le format d'une date ambiguë, et renvoyer une liste vide en expliquant pourquoi si la demande ne peut pas se faire avec ces opérations. La réponse est ensuite vérifiée par le code contre les colonnes de la table.

\subsection{Aperçu et application sur une copie}

Avant toute application, l'interface affiche un aperçu chiffré qui se met à jour à chaque modification de la liste : nombre de valeurs touchées, nombre à revoir, nombre non reconnues, exemples et avertissements. L'application crée alors une \textbf{nouvelle version}, dérivée de la première, analysée comme un dépôt et accompagnée d'un journal ; la version d'origine n'est jamais modifiée. Les fichiers corrigés sont écrits en UTF-8, séparés par des virgules, avec l'en-tête inchangé. Des essais sur de vraies données ont conduit à une règle supplémentaire : ne rien changer qui n'ait été demandé (en-tête, encodage, données tronquées, fuseaux horaires).

\subsection{Modélisation dynamique}

Le diagramme de séquence de la figure~\ref{fig:seq-corrections} décrit une correction écrite en langage naturel, de la demande à la nouvelle version.

\figplaceholder[1]{seq_corrections.png}{Diagramme de séquence — Correction d'une table}{fig:seq-corrections}{Diagramme de séquence (généré).}

\subsection{Résultats du sprint 5}

L'utilisateur peut désormais corriger une table en quelques clics ou en une phrase, en voyant exactement ce qui sera modifié, et retrouver sa version d'origine intacte. La figure~\ref{fig:capture-corrections} présente l'onglet « Corrections ».

\figplaceholder[0.95]{capture_corrections.png}{Onglet « Corrections »}{fig:capture-corrections}{Capture de l'onglet « Corrections » : corrections trouvées, demande en langage naturel, aperçu chiffré, bouton d'application sur une copie.}

\noindent\textbf{Sprint Review.} L'incrément livré permet de passer de la détection à la correction. La rétrospective a mis en évidence que les tables restaient vues isolément : les liens entre les entités n'étaient exploitables ni visuellement ni par des requêtes, ce qui fait l'objet du sprint~6.

%==============================================================================
\section{Sprint 6 — Graphe de connaissances}
%==============================================================================

Le sixième sprint réalise l'epic E6. Nous présentons son objectif et son backlog, puis la déduction du modèle de graphe, les cas particuliers, la gestion des conflits et la construction dans Neo4j.

\subsection{Sprint Goal et Sprint Backlog}

Le \textbf{Sprint Goal} du sprint~6 est le suivant : \textit{« représenter les tables comme des entités et des relations dans Neo4j »}. Le tableau~\ref{tab:backlog-s6} détaille les récits retenus et leurs tâches.

\begin{table}[H]
\centering
\small
\begin{tabular}{|C{1.1cm}|L{4.6cm}|C{1.2cm}|L{6.4cm}|}
\hline
\thd{US} & \thd{Récit utilisateur} & \thd{ID} & \thd{Tâche} \\ \hline
\multirow{2}{*}{US23} & \multirow{2}{4.6cm}{Modèle de graphe} & T6.1 & Déduction des entités, propriétés et relations \\ \cline{3-4}
 & & T6.2 & Noms proposés par le modèle, relecture avant construction \\ \hline
\multirow{4}{*}{US24} & \multirow{4}{4.6cm}{Cas courants} & T6.3 & Identification par une combinaison de références \\ \cline{3-4}
 & & T6.4 & Table de liaison transformée en relation \\ \cline{3-4}
 & & T6.5 & Colonne de numéros de ligne d'un fichier sans identifiant \\ \cline{3-4}
 & & T6.6 & Référence vers la même table \\ \hline
\multirow{3}{*}{US25} & \multirow{3}{4.6cm}{Construction et suppression} & T6.7 & Signalement des conflits et des liens incertains \\ \cline{3-4}
 & & T6.8 & Construction dans Neo4j avec périmètre par version \\ \cline{3-4}
 & & T6.9 & Suppression du graphe d'une version \\ \hline
\end{tabular}
\caption{Sprint Backlog du sprint 6}
\label{tab:backlog-s6}
\end{table}

\subsection{Déduction du modèle de graphe}

Le modèle est déduit des rôles de colonnes et des relations entre tables mesurées au sprint~1. Une ligne d'une table de comptes devient un nœud \textit{Compte}. Une colonne qui référence un client crée un nœud \textit{Client} à part entière, dont les propriétés sont les colonnes constantes pour un client, et une relation entre les deux. Le modèle proposé est présenté à l'utilisateur dans l'onglet « Graphe » avant la construction : entités, propriétés, relations, conflits. Le modèle peut aussi proposer des noms plus parlants pour les entités et les relations, que l'utilisateur peut accepter ou modifier.

\subsection{Cas particuliers}

Des essais sur des jeux de données variés ont montré que les fichiers réels n'ont pas toujours d'identifiant déclaré. Quatre cas sont donc gérés (chapitre~4) : l'identification par une combinaison de références, la colonne de numéros de ligne d'un autre fichier, la table de deux références transformée en relation, et la référence d'une table vers elle-même. Les marqueurs d'absence (\texttt{NULL}, \texttt{N/A}, \texttt{NaN}) sont lus comme « pas de valeur », comme dans l'analyse.

\subsection{Conflits, incertitudes et périmètre}

Lorsque les lignes d'une même clé se contredisent, le graphe conserve toutes les valeurs dans une propriété de candidats, marque le nœud comme conflictuel et ne choisit aucune valeur : la plateforme ne tranche pas à la place de l'utilisateur. Les références vers un nœud inexistant, les tables liées à rien et les liens incertains sont listés avant la construction. Chaque nœud porte une propriété de périmètre (\texttt{d<jeu>v<version>}) : les graphes de plusieurs versions coexistent dans la même base sans se mélanger. Un bouton permet de supprimer le graphe d'une version, et l'onglet signale la présence d'autres graphes dans Neo4j.

\subsection{Modélisation dynamique}

Le diagramme de séquence de la figure~\ref{fig:seq-graphe} décrit la proposition et la construction du graphe.

\figplaceholder[1]{seq_graphe.png}{Diagramme de séquence — Proposition et construction du graphe}{fig:seq-graphe}{Diagramme de séquence (généré).}

\subsection{Résultats du sprint 6}

L'utilisateur obtient pour chaque version un graphe qu'il a pu relire avant de le construire, et une requête prête à copier dans le navigateur de Neo4j, avec le filtre de périmètre déjà écrit. La figure~\ref{fig:capture-graphe} présente l'onglet « Graphe » et la figure~\ref{fig:capture-neo4j} le graphe dans Neo4j.

\figplaceholder[0.95]{capture_graphe.png}{Onglet « Graphe »}{fig:capture-graphe}{Capture de l'onglet « Graphe » : modèle proposé, conflits, liens incertains, construction et suppression.}

\figplaceholder[0.9]{capture_neo4j.png}{Graphe d'une version dans Neo4j Browser}{fig:capture-neo4j}{Capture de Neo4j Browser : nœuds et relations d'une version, avec le filtre de périmètre.}

\noindent\textbf{Sprint Review.} L'incrément livré met les données en relation et les rend interrogeables. La rétrospective a conclu que l'interroger demandait de connaître Cypher, ce qui fait l'objet du sprint~7.

%==============================================================================
\section{Sprint 7 — Questions en langage naturel}
%==============================================================================

Le septième sprint réalise l'epic E7. Nous présentons son objectif et son backlog, puis le tri des questions, la génération contrôlée des requêtes, la rédaction des réponses, l'évaluation rejouable et l'expérience d'utilisation de l'interface.

\subsection{Sprint Goal et Sprint Backlog}

Le \textbf{Sprint Goal} du sprint~7 est le suivant : \textit{« répondre en langage naturel aux questions sur les données, leur qualité et l'usage de l'outil »}. Le tableau~\ref{tab:backlog-s7} détaille les récits retenus et leurs tâches.

\begin{table}[H]
\centering
\small
\begin{tabular}{|C{1.1cm}|L{4.6cm}|C{1.2cm}|L{6.4cm}|}
\hline
\thd{US} & \thd{Récit utilisateur} & \thd{ID} & \thd{Tâche} \\ \hline
US26 & Tri des questions & T7.1 & Classement en données, qualité ou aide \\ \hline
\multirow{2}{*}{US27} & \multirow{2}{4.6cm}{Requêtes contrôlées} & T7.2 & Génération de Cypher par le modèle \\ \cline{3-4}
 & & T7.3 & Garde-fous : lecture seule, périmètre, \texttt{EXPLAIN}, durée, tentatives \\ \hline
\multirow{3}{*}{US28} & \multirow{3}{4.6cm}{Réponses sourcées} & T7.4 & Faits de qualité joints par le serveur \\ \cline{3-4}
 & & T7.5 & Guide d'utilisation de l'application \\ \cline{3-4}
 & & T7.6 & Restitution de la source, de la requête et des lignes \\ \hline
\multirow{3}{*}{US29} & \multirow{3}{4.6cm}{Expérience d'utilisation} & T7.7 & Page d'accueil publique \\ \cline{3-4}
 & & T7.8 & Système d'animation cohérent \\ \cline{3-4}
 & & T7.9 & Palette de commandes et raccourcis clavier \\ \hline
\end{tabular}
\caption{Sprint Backlog du sprint 7}
\label{tab:backlog-s7}
\end{table}

\subsection{Tri des questions}

Une question peut porter sur le contenu des données, sur leur qualité ou sur la manière d'utiliser l'application. Le service commence par la classer, ce qui détermine la source de la réponse (tableau~\ref{tab:types-questions}, chapitre~4). Cette étape est décisive : une question sur les règles en échec ne doit pas déclencher une requête sur le graphe, et elle doit pouvoir être traitée avant même que le graphe soit construit. Le diagramme d'activité de la figure~\ref{fig:act-questions} résume le parcours d'une question.

\figplaceholder[0.85]{act_questions.png}{Diagramme d'activité — Parcours d'une question}{fig:act-questions}{Diagramme d'activité (généré).}

\subsection{Génération contrôlée des requêtes}

Pour une question sur les données, le modèle écrit une requête Cypher à partir du schéma du graphe et d'exemples. Cette requête est ensuite soumise aux garde-fous décrits au chapitre~4 : instruction unique en lecture seule, filtre de périmètre sur chaque nœud, vérification du plan par \texttt{EXPLAIN}, durée limitée, nombre de lignes borné. En cas de refus, le motif est renvoyé au modèle, qui dispose de deux tentatives. Le diagramme de séquence de la figure~\ref{fig:seq-questions} décrit l'ensemble.

\figplaceholder[1]{seq_questions.png}{Diagramme de séquence — Question sur les données}{fig:seq-questions}{Diagramme de séquence (généré).}

\subsection{Rédaction des réponses et traçabilité}

La réponse est rédigée uniquement à partir des lignes obtenues, des faits de qualité que le serveur joint à la question (règles et message de chacune, dimensions, constats, rôles de colonnes) et du guide d'utilisation. Le guide est un fichier de l'application, à tenir à jour lorsque l'interface change ; il définit aussi les scores et les gravités. Chaque réponse indique sa provenance (graphe, résultats de qualité ou guide) et, pour le graphe, affiche la requête exécutée et les lignes obtenues, afin que l'utilisateur puisse la vérifier. Les questions de plus de 1\,000 caractères sont refusées.

\subsection{Évaluation rejouable}

Un jeu d'une vingtaine de questions (données, qualité, usage, demandes d'écriture à refuser) peut être rejoué contre les services en fonctionnement. Le script vérifie le type de la réponse et la présence des mots attendus, et renvoie un code d'erreur si un cas échoue, ce qui permet de comparer deux modèles. Les cas fournis sont écrits pour les jeux de démonstration du projet. Les résultats de cette évaluation sont présentés au chapitre~\ref{chap:benchmark}.

\subsection{Expérience d'utilisation}

Le dernier sprint a aussi soigné l'interface. Une \textbf{page d'accueil publique} présente la plateforme à un visiteur non technique : fonctionnalités, fonctionnement, aperçu et bouton de connexion. Le \textbf{tableau de bord} regroupe tous les jeux de données avec leurs chiffres clés, des filtres (tous, à surveiller, sains, en cours), un tri et un anneau de score animé par jeu de données. Un \textbf{système d'animation} unique (durées et courbes communes, transitions de page, apparition au défilement, sorties animées des menus, fenêtres et notifications) a été mis au point en privilégiant les propriétés qui ne déclenchent pas de recalcul de mise en page, et respecte la préférence « réduire les animations » du système. Une \textbf{palette de commandes} (\texttt{Ctrl/Cmd+K} ou \texttt{/}) donne accès aux jeux de données, aux tables, aux onglets et aux actions ; \texttt{N} ouvre une nouvelle analyse et les chiffres de 1 à 9 changent d'onglet.

\subsection{Résultats du sprint 7}

L'utilisateur peut interroger ses données, leur qualité et l'outil lui-même dans sa propre langue, et vérifier chaque réponse. La figure~\ref{fig:capture-ask} présente l'onglet « Questions », et la figure~\ref{fig:capture-accueil} la page d'accueil publique.

\figplaceholder[0.95]{capture_ask.png}{Onglet « Questions »}{fig:capture-ask}{Capture de l'onglet « Questions » : exemples cliquables, réponse avec sa provenance, requête exécutée et lignes obtenues.}

\figplaceholder[0.95]{capture_accueil.png}{Page d'accueil publique}{fig:capture-accueil}{Capture de la page d'accueil : présentation, fonctionnalités, fonctionnement et bouton de connexion.}

\noindent\textbf{Sprint Review.} L'incrément livré complète la plateforme. Il n'a pas été suivi d'un nouveau sprint dans le cadre du stage ; ses limites et les pistes d'évolution sont présentées dans la conclusion.

\section*{Conclusion}
\addcontentsline{toc}{section}{Conclusion}
Ce chapitre a présenté les sprints~5 à~7 : les corrections appliquées sur une copie, le graphe de connaissances dans Neo4j et les questions en langage naturel, ainsi que le soin apporté à l'expérience d'utilisation. Le chapitre suivant présente l'évaluation de la solution : les tests automatisés, l'évaluation des questions et le choix du modèle de langage.

\chapter{Évaluation de la Solution et Choix du Modèle}\label{chap:benchmark}

\section*{Introduction}
\addcontentsline{toc}{section}{Introduction}
Les chapitres précédents ont montré que la qualité des résultats de \projet{} repose sur deux piliers : un code déterministe, vérifiable par des tests automatisés, et un modèle de langage, dont le comportement ne peut pas se prouver mais doit se mesurer. Ce chapitre présente d'abord les tests automatisés et les tests fonctionnels, puis une évaluation comparative de trois modèles de langage sur un jeu de questions rejouable, qui éclaire le choix du modèle sur la tâche où il compte le plus : répondre à des questions en langage naturel.

%==============================================================================
\section{Tests automatisés}
%==============================================================================

Le code déterministe de la plateforme est couvert par 292 tests automatisés répartis en quatre suites (tableau~\ref{tab:tests}), exécutées à chaque modification par l'intégration continue. Elles vérifient en particulier les comportements dont dépend la confiance dans les résultats : le rejet d'un téléversement invalide avant toute écriture, la détection d'un dépôt identique, la reprise des exécutions bloquées, l'évaluation de chaque type de règle, le refus par les garde-fous des requêtes qui écrivent ou qui sortent du périmètre d'une version, l'attribution des rôles de colonnes et la validation des réponses du modèle. Ces tests n'ont besoin ni d'Airflow, ni de Neo4j, ni d'Ollama, qui sont simulés. Les suites du worker (67 tests) et du service de questions (39 tests) ont été rejouées lors de la rédaction de ce rapport et réussissent toutes.

%==============================================================================
\section{Tests fonctionnels}
%==============================================================================

Le tableau~\ref{tab:tests-fonctionnels} présente onze scénarios de test fonctionnel, joués sur la plateforme en fonctionnement. Le scénario TF4, qui suppose un modèle injoignable, a été vérifié par un test automatisé du worker plutôt qu'en arrêtant Ollama.

\begin{table}[H]
\centering
\small
\begin{tabular}{|C{1cm}|L{3.6cm}|L{4.2cm}|L{4cm}|}
\hline
\thd{ID} & \thd{Scénario} & \thd{Résultat attendu} & \thd{Résultat obtenu} \\ \hline
TF1 & Dépôt d'un fichier invalide parmi plusieurs & Aucun fichier écrit, fichier fautif signalé & Réussi : HTTP 400, aucun jeu de données créé \\ \hline
TF2 & Table à l'encodage et au séparateur non standard & Lecture correcte de toutes les colonnes & Réussi : 4 colonnes et 20 lignes lues \\ \hline
TF3 & Dépôt identique à une version existante & Proposition d'ouvrir ou de relancer & Réussi : doublon détecté, pas de relance \\ \hline
TF4 & Modèle indisponible & Analyse déterministe, statut du modèle \texttt{unavailable} & Réussi (test automatisé : modèle injoignable simulé) \\ \hline
TF5 & Redémarrage du serveur pendant une exécution & Exécution reprise ou signalée, sans blocage & Réussi : \texttt{COMPLETED} malgré le redémarrage, 6\,s après le dépôt \\ \hline
TF6 & Règle de format sur une colonne courriel & Échec avec nombre de valeurs invalides & Réussi : échec à 93,0\,\% avec un exemple \\ \hline
TF7 & Correction d'une date ambiguë & Date non devinée, avertissement affiché & Réussi : 2 dates converties, 3 ambiguës laissées à la revue \\ \hline
TF8 & Correction appliquée & Nouvelle version créée, version d'origine intacte & Réussi : version 2 dérivée, fichier d'origine identique (même empreinte) \\ \hline
TF9 & Question demandant de supprimer des données & Requête refusée, aucune écriture dans Neo4j & Réussi : aucune requête exécutée, 548 nœuds avant et après \\ \hline
TF10 & Question sur les données sans graphe construit & Message demandant de construire le graphe & Réussi : demande de construire le graphe \\ \hline
TF11 & Déploiement par Docker Compose & Services démarrés et sains & Réussi : conteneurs démarrés, ceux d'Airflow sains \\ \hline
\end{tabular}
\caption{Scénarios de tests fonctionnels}
\label{tab:tests-fonctionnels}
\end{table}

%==============================================================================
\section{Évaluation comparative de trois modèles sur les questions}
%==============================================================================

Les tests précédents vérifient le code, pas le comportement d'un modèle réel. Nous avons donc comparé trois modèles sur la tâche où leur qualité se voit le mieux : le service de questions, qui leur demande de classer une question, d'écrire une requête Cypher, puis de rédiger une réponse fidèle aux lignes obtenues.

\subsection{Modèles candidats}

Les modèles comparés sont ceux qui sont installés dans Ollama sur la machine de développement (tableau~\ref{tab:bm-modeles}) : un modèle de taille moyenne qui tient en mémoire graphique, un modèle plus grand qui ne tient pas entièrement en mémoire graphique, et un modèle servi à distance.

\begin{table}[H]
\centering
\small
\begin{tabular}{|L{4cm}|L{2.6cm}|C{2.6cm}|L{3.8cm}|}
\hline
\thd{Modèle} & \thd{Éditeur} & \thd{Taille sur disque} & \thd{Exécution} \\ \hline
Qwen 2.5 14B Instruct & Alibaba Cloud & 9,0 Go & Locale, mémoire graphique \\ \hline
Qwen 3.5 35B & Alibaba Cloud & 23 Go & Locale, avec déchargement vers la mémoire vive \\ \hline
Gemma 4 (\texttt{gemma4:cloud}) & Google & -- & \textbf{Distante} : servie par le service cloud d'Ollama \\ \hline
\end{tabular}
\caption{Modèles candidats}
\label{tab:bm-modeles}
\end{table}

\noindent\textbf{Remarque de confidentialité.} Le modèle \texttt{gemma4:cloud} n'est pas exécuté sur la machine : Ollama transmet les requêtes à un service distant. Il a été utilisé pendant le développement et dans cette évaluation sur des données fictives, mais il ne respecte pas l'exigence de confidentialité du projet et ne doit pas recevoir de données client. Pour des données sensibles, la plateforme doit être configurée avec un modèle local.

\subsection{Protocole}

Le jeu comprend 19 questions (une question par cas) portant sur trois jeux de données de démonstration : un magasin de vélos de neuf tables (\texttt{MyBikeStore}), un jeu de comptes (\texttt{check-accounts}) et un jeu de clients et de commandes (\texttt{graph-shop}). Le graphe de chaque version est construit avant l'évaluation, et deux règles de qualité sont approuvées sur \texttt{check-accounts}, comme le suppose l'un des cas. Les questions couvrent quatre catégories : questions sur les \textbf{données} (réponse par le graphe), sur la \textbf{qualité} (règles, scores, constats), sur l'\textbf{usage} de l'outil, et des cas \textbf{sans type attendu} (définitions, demandes d'écriture qui doivent être refusées, questions hors sujet).

Pour chaque modèle, le service de questions est redémarré avec le modèle voulu (variable \texttt{OLLAMA\_MODEL}), puis les 19 questions sont posées dans le même ordre, une fois chacune, par l'intermédiaire du serveur applicatif. Un cas est réussi si le type de la réponse est celui attendu et si elle contient les éléments attendus (par exemple le nombre exact de commandes) sans contenir de mot interdit (par exemple une confirmation de suppression). Pour une question sur les données, la requête exécutée doit aussi être fournie avec la réponse. Nous relevons en outre la latence de chaque cas, le nombre de requêtes Cypher produites et le nombre de tentatives de requêtes refusées par les garde-fous.

\subsection{Résultats}

Le tableau~\ref{tab:eval-modeles} et la figure~\ref{fig:eval-modeles} présentent les résultats.

\begin{table}[H]
\centering
\small
\begin{tabular}{|L{3.2cm}|C{1.7cm}|C{1.5cm}|C{1.4cm}|C{1.1cm}|C{1.3cm}|C{1.6cm}|C{1.5cm}|}
\hline
\thd{Modèle} & \thd{Total} & \thd{Données} & \thd{Qualité} & \thd{Aide} & \thd{Autres} & \thd{Latence moy. (s)} & \thd{Refus Cypher} \\ \hline
Gemma 4 (cloud) & 19/19 (100\,\%) & 6/6 & 4/4 & 2/2 & 7/7 & 1,4 & 0 \\ \hline
Qwen 3.5 35B & 18/19 (95\,\%) & 5/6 & 4/4 & 2/2 & 7/7 & 25,0 & 0 \\ \hline
Qwen 2.5 14B Instruct & 17/19 (89\,\%) & 5/6 & 4/4 & 2/2 & 6/7 & 14,4 & 2 \\ \hline
\end{tabular}
\caption{Résultats de l'évaluation des questions pour trois modèles}
\label{tab:eval-modeles}
\end{table}

\figplaceholder[1]{eval_modeles.png}{Taux de réussite par catégorie et latence moyenne des trois modèles}{fig:eval-modeles}{Deux graphiques : à gauche, cas réussis par catégorie pour chaque modèle ; à droite, latence moyenne par modèle.}

\subsection{Analyse des échecs}

Les trois échecs se répartissent sur deux cas seulement.
\begin{itemize}
  \item \textbf{« Combien de clients et de commandes ? »} (\texttt{graph-shop} : 40 clients et 120 commandes). Les deux modèles locaux échouent, de deux façons différentes. Qwen 3.5 35B produit une requête admissible, que les garde-fous laissent passer, mais sa réponse annonce « 4\,800 clients et 4\,800 commandes » : ce nombre est égal à $40 \times 120$, ce qui évoque un produit cartésien entre les deux ensembles. Les garde-fous contrôlent la sécurité d'une requête, pas son exactitude : c'est pourquoi la requête exécutée est toujours affichée avec la réponse, pour que l'utilisateur puisse repérer ce type d'erreur. Qwen 2.5 14B voit ses deux requêtes refusées par les garde-fous ; la réponse se rabat alors sur les faits de qualité, sans requête, et elle est fausse (le contrôle la signale comme « réponse du graphe sans requête »). Aucune requête dangereuse n'a été exécutée dans ce second cas.
  \item \textbf{« Quelle différence entre un avertissement et une règle critique ? »} Qwen 2.5 14B échoue le contrôle lexical : sa réponse, plausible, n'emploie aucun des termes attendus (« à corriger en premier », « rouge », « séparément »). Ce cas illustre la limite d'un contrôle par mots-clés, qui peut rejeter une réponse correcte formulée autrement.
\end{itemize}
Gemma 4 réussit les 19 cas. Aucun modèle ne s'est trompé sur le type d'une question ni ne s'est plié à une demande d'écriture ou à une question hors sujet.

\subsection{Latence}

Les réponses du modèle distant arrivent en moyenne en 1,4\,s, contre 14,4\,s pour Qwen 2.5 14B et 25,0\,s pour Qwen 3.5 35B, le plus lent des trois, ce qui est cohérent avec sa taille (23\,Go, déchargé en partie vers la mémoire vive). Ces durées, de l'ordre de la dizaine de secondes, sont acceptables pour une question posée dans l'interface.

%==============================================================================
\section{Choix du modèle}
%==============================================================================

Le modèle intervient dans quatre tâches de la plateforme, dont les exigences diffèrent (tableau~\ref{tab:bm-taches}).

\begin{table}[H]
\centering
\small
\begin{tabular}{|L{4cm}|L{6cm}|L{4cm}|}
\hline
\thd{Tâche} & \thd{Exigence principale} & \thd{Réglage} \\ \hline
Analyse d'une table & Validité du format, fidélité aux chiffres, faible latence (nombreuses tables) & \texttt{OLLAMA\_MODEL} du worker \\ \hline
Analyse globale & Raisonnement riche, grand contexte ; appel unique & Même modèle, appel désactivable \\ \hline
Correction en langage naturel & Respect strict des opérations autorisées, aucune valeur inventée & \texttt{OLLAMA\_MODEL} du serveur \\ \hline
Questions & Génération de Cypher admissible et exacte, réponses fidèles aux lignes obtenues & \texttt{OLLAMA\_MODEL} du service \\ \hline
\end{tabular}
\caption{Tâches confiées au modèle et exigences}
\label{tab:bm-taches}
\end{table}

Aucune de ces tâches n'est liée à un modèle précis : chacune lit le modèle dans la configuration. Les résultats appellent trois conclusions.
\begin{itemize}
  \item Le modèle distant est le plus juste et de loin le plus rapide, mais il ne respecte pas la confidentialité : il ne convient qu'à des données fictives ou publiques.
  \item Pour des données sensibles, le choix se fait entre les deux modèles locaux. Qwen 3.5 35B a réussi un cas de plus que Qwen 2.5 14B, mais il est près de deux fois plus lent et occupe 23\,Go. Avec 19 cas et une seule exécution, un écart d'un cas n'est pas significatif : nous gardons le modèle de taille moyenne par défaut, plus rapide et qui tient en mémoire graphique, et la plateforme laisse le choix à l'utilisateur.
  \item Quel que soit le modèle, la requête exécutée doit rester visible avec la réponse : c'est la seule défense contre une requête admissible mais erronée, comme le produit cartésien observé.
\end{itemize}

%==============================================================================
\section{Limites de l'évaluation}
%==============================================================================

Plusieurs limites doivent être gardées à l'esprit dans l'interprétation des résultats :
\begin{itemize}
  \item la \textbf{taille du jeu} : 19 questions, posées une seule fois à chaque modèle, sur trois jeux de démonstration ; les écarts de un ou deux cas ne sont pas statistiquement significatifs ;
  \item le \textbf{contrôle par mots-clés}, qui ne vérifie ni la qualité de rédaction ni la pertinence d'une explication, et peut rejeter une réponse correcte ;
  \item la \textbf{dépendance aux consignes} : les résultats sont mesurés avec un jeu de consignes donné, qu'un autre modèle pourrait mieux exploiter s'il était adapté ;
  \item la \textbf{dépendance au matériel} : la latence des modèles locaux dépend de la mémoire graphique disponible ;
  \item la \textbf{portée de l'évaluation} : seules les questions ont été comparées entre modèles. Une première mesure sur l'analyse d'une table, avec quatre tables fictives aux défauts injectés, n'a pas été retenue : les constats structurés étant bornés par les mesures du pipeline, les trois modèles obtenaient exactement les mêmes résultats, et cette mesure ne permettait pas de les départager ;
  \item l'\textbf{absence d'évaluation de la pertinence} des recommandations rédigées par le modèle, qui exigerait plusieurs évaluateurs humains.
\end{itemize}

\section*{Conclusion}
\addcontentsline{toc}{section}{Conclusion}
Ce chapitre a présenté l'évaluation de \projet{} : 292 tests automatisés qui couvrent le code déterministe, onze scénarios fonctionnels joués sur la plateforme, et une comparaison de trois modèles sur 19 questions rejouables. Cette comparaison montre des différences réelles de justesse et de latence, identifie le type d'erreur le plus difficile à détecter (une requête admissible mais inexacte) et justifie le choix de toujours montrer la requête exécutée avec la réponse.

\chapter*{Conclusion Générale et Perspectives}
\markboth{Conclusion Générale}{}
\addcontentsline{toc}{chapter}{Conclusion Générale et Perspectives}

Ce projet de fin d'études, réalisé au sein de Wevioo, avait pour objectif d'automatiser de bout en bout l'évaluation de la qualité de jeux de données tabulaires, en produisant non seulement des indicateurs fiables, mais aussi une interprétation, une surveillance, des corrections et une interrogation en langage naturel, dans le respect de la confidentialité des données.

Pour y répondre, nous avons conçu et réalisé \textbf{\projet}, une plateforme web fondée sur un principe unique : \textbf{le code mesure et contrôle, le modèle interprète et rédige}. Le traitement de profilage, orchestré par Airflow, lit de manière robuste des exports hétérogènes, calcule les métriques de chaque table, détermine le rôle de chaque colonne à partir d'indices explicites, mesure les formats, l'ordre des dates et les liens entre tables, puis fait analyser chaque table par un modèle local dont les affirmations sont confrontées aux mesures, avec une analyse déterministe de repli. Un serveur Spring Boot versionne les dépôts, sécurise l'accès par jeton, évalue des règles de qualité à chaque exécution, applique des corrections sur une copie et construit le graphe de chaque version dans Neo4j. Une interface Angular restitue ces résultats en temps réel, et un service de questions répond en langage naturel sur les données, leur qualité et l'usage de l'outil, avec des garde-fous qui n'exécutent que des requêtes en lecture seule limitées au périmètre d'une version.

Les contraintes initiales ont été respectées : la confidentialité est garantie par l'exécution locale du modèle, de la base de graphe et du service de questions ; la reproductibilité, par un déploiement en une commande ; la fiabilité, par des statuts explicites, une reprise après incident et 292 tests automatisés ; et la traçabilité, par le statut du modèle, le journal des corrections et la restitution de la requête et des lignes de chaque réponse.

Ce travail a également permis de tirer des enseignements sur l'usage des LLM dans un contexte de qualité des données. Un modèle est utile pour interpréter, rédiger et traduire une demande, mais il ne doit ni calculer ni décider seul : confronter chaque chiffre à sa colonne, limiter le modèle à des opérations connues et contrôler chaque requête avant exécution ont été plus efficaces que d'affiner indéfiniment les consignes. De même, signaler une dimension « non mesurée » plutôt que d'afficher une valeur flatteuse, ou conserver toutes les valeurs d'un conflit plutôt que d'en choisir une, renforce la confiance de l'utilisateur.

Enfin, le projet comporte des limites qu'il convient de souligner : un seul rôle utilisateur est défini à ce stade ; la base de métadonnées par défaut est une base de fichiers adaptée à un poste, non à une montée en charge ; la plateforme lit des fichiers CSV et non directement des bases de données ; et l'évaluation des modèles repose sur un jeu de 19 questions, posées une seule fois à chaque modèle.

Sur le plan personnel, ce stage m'a permis de consolider mes compétences en ingénierie des données et en développement d'applications web (Java, Angular, Python), de découvrir en profondeur l'exploitation des LLM en environnement local, les bases de graphe et la sécurisation d'un modèle qui écrit des requêtes, et de pratiquer une démarche agile, de la vision du produit à la revue de sprint. Il m'a aussi appris à évaluer un modèle de langage avec méthode : mesurer plutôt que supposer, documenter ses limites et contrôler ce qu'il produit avant de s'y fier.

Ces résultats et ces limites ouvrent plusieurs pistes d'évolution, présentées ci-après.

\section*{Perspectives}

Plusieurs évolutions permettraient d'enrichir la solution :
\begin{itemize}
  \item \textbf{Gestion de plusieurs rôles} : distinguer l'ingénieur de données, l'analyste et le responsable métier, avec des droits adaptés (lecture seule, approbation des règles, correction).
  \item \textbf{Sources de données élargies} : connecter directement la plateforme aux bases de données et aux entrepôts (SQL, Parquet) plutôt qu'aux seuls exports CSV.
  \item \textbf{Montée en charge} : migrer la base de métadonnées vers un serveur de base de données, traiter les tables en parallèle et suivre les temps d'exécution.
  \item \textbf{Évaluation plus large des modèles} : enrichir le jeu de questions, répéter chaque mesure, évaluer la qualité de rédaction par des évaluateurs humains et comparer les modèles sur les corrections en langage naturel.
  \item \textbf{Règles exécutables} : exporter les règles approuvées vers des outils de validation existants (par exemple Great Expectations), pour les vérifier dans les chaînes de données des clients.
  \item \textbf{Graphe enrichi} : exploiter des algorithmes de graphe (communautés, centralité) pour détecter les entités dupliquées et les anomalies de structure.
  \item \textbf{Mémoire de conversation} : permettre des questions de suivi dans l'onglet de questions, en conservant le contexte des réponses précédentes.
\end{itemize}


%------------------------------ Bibliographie ---------------------------------
\cleardoublepage
\markboth{Bibliographie}{}
\begin{thebibliography}{99}
\addcontentsline{toc}{chapter}{Bibliographie}

\bibitem{dama} DAMA International, \textit{DAMA-DMBOK: Data Management Body of Knowledge}, 2\textsuperscript{e} édition, Technics Publications, 2017.

\bibitem{iso25012} ISO/IEC 25012:2008, \textit{Software engineering — Software product Quality Requirements and Evaluation (SQuaRE) — Data quality model}, International Organization for Standardization, 2008.

\bibitem{wang} R. Y. Wang et D. M. Strong, « Beyond Accuracy: What Data Quality Means to Data Consumers », \textit{Journal of Management Information Systems}, vol. 12, n\textsuperscript{o} 4, p. 5--33, 1996.

\bibitem{abedjan} Z. Abedjan, L. Golab et F. Naumann, « Profiling relational data: a survey », \textit{The VLDB Journal}, vol. 24, p. 557--581, 2015.

\bibitem{sweetviz} F. Bertrand, \textit{Sweetviz — Visualize and compare datasets, target values and associations, with one line of code}, documentation officielle. \url{https://github.com/fbdesignpro/sweetviz}

\bibitem{ydata} YData, \textit{ydata-profiling (ex pandas-profiling) — Documentation}. \url{https://docs.profiling.ydata.ai}

\bibitem{gx} Great Expectations, \textit{Great Expectations Documentation}. \url{https://docs.greatexpectations.io}

\bibitem{deequ} S. Schelter \textit{et al.}, « Automating large-scale data quality verification », \textit{Proceedings of the VLDB Endowment}, vol. 11, n\textsuperscript{o} 12, p. 1781--1794, 2018.

\bibitem{soda} Soda, \textit{Soda Core Documentation}. \url{https://docs.soda.io}

\bibitem{vaswani} A. Vaswani \textit{et al.}, « Attention Is All You Need », \textit{Advances in Neural Information Processing Systems (NeurIPS)}, 2017.

\bibitem{brown} T. Brown \textit{et al.}, « Language Models are Few-Shot Learners », \textit{Advances in Neural Information Processing Systems (NeurIPS)}, 2020.

\bibitem{qwen25} Qwen Team, « Qwen2.5 Technical Report », arXiv:2412.15115, 2024.

\bibitem{ollama} Ollama, \textit{Ollama — Get up and running with large language models locally}, documentation et API REST. \url{https://github.com/ollama/ollama}

\bibitem{narayan} A. Narayan, I. Chami, L. Orr et C. Ré, « Can Foundation Models Wrangle Your Data? », \textit{Proceedings of the VLDB Endowment}, vol. 16, n\textsuperscript{o} 4, 2022.

\bibitem{ji} Z. Ji \textit{et al.}, « Survey of Hallucination in Natural Language Generation », \textit{ACM Computing Surveys}, vol. 55, n\textsuperscript{o} 12, 2023.

\bibitem{airflow} The Apache Software Foundation, \textit{Apache Airflow Documentation (version 2.10)}. \url{https://airflow.apache.org/docs}

\bibitem{docker} Docker Inc., \textit{Docker Compose Documentation}. \url{https://docs.docker.com/compose}

\bibitem{fielding} R. T. Fielding, \textit{Architectural Styles and the Design of Network-based Software Architectures}, thèse de doctorat, University of California, Irvine, 2000.

\bibitem{jwt} M. Jones, J. Bradley et N. Sakimura, « JSON Web Token (JWT) », RFC 7519, IETF, 2015.

\bibitem{spring} Broadcom, \textit{Spring Boot Reference Documentation}. \url{https://docs.spring.io/spring-boot}

\bibitem{angular} Google, \textit{Angular Documentation}. \url{https://angular.dev}

\bibitem{neo4j} Neo4j, Inc., \textit{Neo4j Cypher Manual}. \url{https://neo4j.com/docs/cypher-manual}

\bibitem{lewis} P. Lewis \textit{et al.}, « Retrieval-Augmented Generation for Knowledge-Intensive NLP Tasks », \textit{Advances in Neural Information Processing Systems (NeurIPS)}, 2020.

\bibitem{edge} D. Edge \textit{et al.}, « From Local to Global: A Graph RAG Approach to Query-Focused Summarization », arXiv:2404.16130, 2024.

\bibitem{wevioo} Wevioo, \textit{À propos de Wevioo} (présentation du groupe), site institutionnel, consulté le 7 novembre 2026. \url{https://wevioo.com/en/about-us/}

\bibitem{scrum} K. Schwaber et J. Sutherland, \textit{The Scrum Guide — The Definitive Guide to Scrum: The Rules of the Game}, 2020.

\bibitem{pandas} W. McKinney, « Data Structures for Statistical Computing in Python », \textit{Proceedings of the 9th Python in Science Conference}, p. 56--61, 2010.

\bibitem{cramer} H. Cramér, \textit{Mathematical Methods of Statistics}, Princeton University Press, 1946.

\end{thebibliography}


%------------------------------ Abstract (4e de couverture) -------------------
\chapter*{Abstract}
\markboth{Abstract}{}
\addcontentsline{toc}{chapter}{Abstract}
\begin{otherlanguage}{english}

In a context where data is a strategic asset, its quality directly determines the reliability of analyses, dashboards and decision models. Yet, in many organizations, data quality assessment remains a manual, time-consuming and expert-dependent task: profiling tools measure without interpreting, and validation tools require rules to be written in advance.

This end-of-studies project, carried out at \textbf{Wevioo}, addresses this problem through the design and implementation of \textbf{\projet}, a web platform for intelligent data quality assessment built on one principle: \textbf{code measures and controls, the model interprets and writes}. A processing pipeline orchestrated by \textbf{Apache Airflow} reads heterogeneous files, computes the metrics of each table, determines the role of each column from explicit clues, and measures formats, date order and links between tables. A \textbf{locally executed large language model} (Ollama) interprets these measurements; its claims are checked against the measurements, with a deterministic fallback analysis.

A \textbf{Spring Boot} server versions uploads, secures access with tokens, evaluates \textbf{quality rules} on each run and raises alerts on failure, applies \textbf{corrections} to a copy without ever modifying the original, and builds a \textbf{knowledge graph} in \textbf{Neo4j}. An \textbf{Angular} interface presents the results in real time. A natural-language question service (\textbf{GraphRAG}) answers about the data, its quality and the use of the tool; the Cypher queries written by the model are checked (read-only, scoped to one version, execution plan, time limit) before they run.

The platform can be deployed with a single Docker Compose command and is covered by 292 automated tests. A comparison of three language models on a replayable question set supports the choice of the model.

\vspace{0.4cm}
\noindent\textbf{Keywords:} data quality, data profiling, large language models, local LLM, Ollama, GraphRAG, Neo4j, Spring Boot, Angular, Apache Airflow, Docker, data governance.
\end{otherlanguage}


\end{document}
